\documentclass[11pt]{article}
% Shared preamble for manuscripts-v2 papers. Load with % Shared preamble for manuscripts-v2 papers. Load with % Shared preamble for manuscripts-v2 papers. Load with \input{../common/preamble.tex}
% immediately after \documentclass[11pt]{article}.
\ifdefined\pdftrailerid\pdftrailerid{}\fi
\usepackage[T1]{fontenc}
\usepackage{lmodern}
\usepackage[letterpaper,margin=1.02in]{geometry}
\usepackage{amsmath,amssymb,amsthm,mathtools,booktabs,array,microtype,enumitem}
\usepackage{tikz}
\usepackage{xurl}
\usepackage[colorlinks=true,allcolors=bluegreen]{hyperref}
\definecolor{bluegreen}{RGB}{22,83,91}
\usepackage[capitalize,nameinlink]{cleveref}

\theoremstyle{plain}
\newtheorem{theorem}{Theorem}[section]
\newtheorem{lemma}[theorem]{Lemma}
\newtheorem{proposition}[theorem]{Proposition}
\newtheorem{corollary}[theorem]{Corollary}
\newtheorem{conjecture}[theorem]{Conjecture}
\newtheorem*{theorem*}{Theorem}
\newtheorem{maintheorem}{Theorem}
\renewcommand{\themaintheorem}{\Alph{maintheorem}}
\theoremstyle{definition}
\newtheorem{definition}[theorem]{Definition}
\newtheorem{example}[theorem]{Example}
\newtheorem{question}[theorem]{Question}
\theoremstyle{remark}
\newtheorem{remark}[theorem]{Remark}
\crefname{maintheorem}{Theorem}{Theorems}
\crefname{conjecture}{Conjecture}{Conjectures}
\crefname{question}{Question}{Questions}
\crefname{example}{Example}{Examples}

\newcommand{\R}{\mathbb R}
\newcommand{\Q}{\mathbb Q}
\newcommand{\Z}{\mathbb Z}
\newcommand{\N}{\mathbb N}
\newcommand{\T}{\mathbb T}
\newcommand{\C}{\mathbb C}
\DeclareMathOperator{\den}{den}
\newcommand{\e}{\mathrm e}

% Immutable Lean source links at research-checkpoint-18.
\newcommand{\leancommit}{e4809cf561573aef588e6ffc491a35cdb79fd9ec}
\newcommand{\leanbase}{https://github.com/eldernos/MathematicsResearch/blob/\leancommit}
% \lean{File}{line}{label}: link to NormalityChecker/File.lean#Lline at the pinned commit.
\newcommand{\lean}[3]{\href{\leanbase/NormalityChecker/#1.lean\#L#2}{#3}}
% \repo{path}{label}: link to any repository path at the pinned commit.
\newcommand{\repo}[2]{\href{https://github.com/eldernos/MathematicsResearch/tree/\leancommit/#1}{#2}}

% Evidence labels (see papers/WAVE2-CONTRACT.md).
\newcommand{\evidence}[1]{\textup{[\emph{#1}]}}

\setlength{\emergencystretch}{1.5em}
\raggedbottom

% immediately after \documentclass[11pt]{article}.
\ifdefined\pdftrailerid\pdftrailerid{}\fi
\usepackage[T1]{fontenc}
\usepackage{lmodern}
\usepackage[letterpaper,margin=1.02in]{geometry}
\usepackage{amsmath,amssymb,amsthm,mathtools,booktabs,array,microtype,enumitem}
\usepackage{tikz}
\usepackage{xurl}
\usepackage[colorlinks=true,allcolors=bluegreen]{hyperref}
\definecolor{bluegreen}{RGB}{22,83,91}
\usepackage[capitalize,nameinlink]{cleveref}

\theoremstyle{plain}
\newtheorem{theorem}{Theorem}[section]
\newtheorem{lemma}[theorem]{Lemma}
\newtheorem{proposition}[theorem]{Proposition}
\newtheorem{corollary}[theorem]{Corollary}
\newtheorem{conjecture}[theorem]{Conjecture}
\newtheorem*{theorem*}{Theorem}
\newtheorem{maintheorem}{Theorem}
\renewcommand{\themaintheorem}{\Alph{maintheorem}}
\theoremstyle{definition}
\newtheorem{definition}[theorem]{Definition}
\newtheorem{example}[theorem]{Example}
\newtheorem{question}[theorem]{Question}
\theoremstyle{remark}
\newtheorem{remark}[theorem]{Remark}
\crefname{maintheorem}{Theorem}{Theorems}
\crefname{conjecture}{Conjecture}{Conjectures}
\crefname{question}{Question}{Questions}
\crefname{example}{Example}{Examples}

\newcommand{\R}{\mathbb R}
\newcommand{\Q}{\mathbb Q}
\newcommand{\Z}{\mathbb Z}
\newcommand{\N}{\mathbb N}
\newcommand{\T}{\mathbb T}
\newcommand{\C}{\mathbb C}
\DeclareMathOperator{\den}{den}
\newcommand{\e}{\mathrm e}

% Immutable Lean source links at research-checkpoint-18.
\newcommand{\leancommit}{e4809cf561573aef588e6ffc491a35cdb79fd9ec}
\newcommand{\leanbase}{https://github.com/eldernos/MathematicsResearch/blob/\leancommit}
% \lean{File}{line}{label}: link to NormalityChecker/File.lean#Lline at the pinned commit.
\newcommand{\lean}[3]{\href{\leanbase/NormalityChecker/#1.lean\#L#2}{#3}}
% \repo{path}{label}: link to any repository path at the pinned commit.
\newcommand{\repo}[2]{\href{https://github.com/eldernos/MathematicsResearch/tree/\leancommit/#1}{#2}}

% Evidence labels (see papers/WAVE2-CONTRACT.md).
\newcommand{\evidence}[1]{\textup{[\emph{#1}]}}

\setlength{\emergencystretch}{1.5em}
\raggedbottom

% immediately after \documentclass[11pt]{article}.
\ifdefined\pdftrailerid\pdftrailerid{}\fi
\usepackage[T1]{fontenc}
\usepackage{lmodern}
\usepackage[letterpaper,margin=1.02in]{geometry}
\usepackage{amsmath,amssymb,amsthm,mathtools,booktabs,array,microtype,enumitem}
\usepackage{tikz}
\usepackage{xurl}
\usepackage[colorlinks=true,allcolors=bluegreen]{hyperref}
\definecolor{bluegreen}{RGB}{22,83,91}
\usepackage[capitalize,nameinlink]{cleveref}

\theoremstyle{plain}
\newtheorem{theorem}{Theorem}[section]
\newtheorem{lemma}[theorem]{Lemma}
\newtheorem{proposition}[theorem]{Proposition}
\newtheorem{corollary}[theorem]{Corollary}
\newtheorem{conjecture}[theorem]{Conjecture}
\newtheorem*{theorem*}{Theorem}
\newtheorem{maintheorem}{Theorem}
\renewcommand{\themaintheorem}{\Alph{maintheorem}}
\theoremstyle{definition}
\newtheorem{definition}[theorem]{Definition}
\newtheorem{example}[theorem]{Example}
\newtheorem{question}[theorem]{Question}
\theoremstyle{remark}
\newtheorem{remark}[theorem]{Remark}
\crefname{maintheorem}{Theorem}{Theorems}
\crefname{conjecture}{Conjecture}{Conjectures}
\crefname{question}{Question}{Questions}
\crefname{example}{Example}{Examples}

\newcommand{\R}{\mathbb R}
\newcommand{\Q}{\mathbb Q}
\newcommand{\Z}{\mathbb Z}
\newcommand{\N}{\mathbb N}
\newcommand{\T}{\mathbb T}
\newcommand{\C}{\mathbb C}
\DeclareMathOperator{\den}{den}
\newcommand{\e}{\mathrm e}

% Immutable Lean source links at research-checkpoint-18.
\newcommand{\leancommit}{e4809cf561573aef588e6ffc491a35cdb79fd9ec}
\newcommand{\leanbase}{https://github.com/eldernos/MathematicsResearch/blob/\leancommit}
% \lean{File}{line}{label}: link to NormalityChecker/File.lean#Lline at the pinned commit.
\newcommand{\lean}[3]{\href{\leanbase/NormalityChecker/#1.lean\#L#2}{#3}}
% \repo{path}{label}: link to any repository path at the pinned commit.
\newcommand{\repo}[2]{\href{https://github.com/eldernos/MathematicsResearch/tree/\leancommit/#1}{#2}}

% Evidence labels (see papers/WAVE2-CONTRACT.md).
\newcommand{\evidence}[1]{\textup{[\emph{#1}]}}

\setlength{\emergencystretch}{1.5em}
\raggedbottom

% Author information approved by the repository owner on 2026-11-25:
%   name: Christopher Anderson; designation: Independent Researcher.
% No academic affiliation, ORCID or personal contact details have been supplied.
% Do not add any that the owner has not approved.
\newcommand{\OwnerAuthor}{Christopher Anderson}
\newcommand{\OwnerAffiliation}{Independent Researcher}
\newcommand{\OwnerContact}{Corrections and comments: \url{https://github.com/eldernos/MathematicsResearch/issues}}
\newcommand{\ReleaseLabel}{manuscripts-v2.0.0}
\newcommand{\LicenseLine}{\copyright\ 2026 Christopher Anderson. Licensed under \href{https://creativecommons.org/licenses/by/4.0/}{CC BY 4.0}.}
% \makebyline{version}{date}: centred author block used under the title.
\newcommand{\makebyline}[2]{%
  {\large \OwnerAuthor\par}%
  \ifx\OwnerAffiliation\empty\else{\small \OwnerAffiliation\par}\fi
  \vspace{0.3em}
  {\small Version #1 \enspace$\cdot$\enspace #2 \enspace$\cdot$\enspace \ReleaseLabel\par}%
  {\footnotesize \OwnerContact\par}%
  {\footnotesize \LicenseLine\par}%
  \hypersetup{pdfauthor={\OwnerAuthor}}}

% This manuscript is a candidate draft and is not part of a release.
\renewcommand{\ReleaseLabel}{candidate draft, not released}
\usepackage[numbers,sort&compress]{natbib}
\hypersetup{pdftitle={Digit distribution of arithmetic generating functions}}

\DeclareMathOperator{\supp}{supp}
\DeclareMathOperator{\sgn}{sgn}
\DeclareMathOperator{\ord}{ord}
\DeclareMathOperator{\Cl}{Cl}
\DeclareMathOperator{\tr}{tr}
\DeclareMathOperator{\rad}{rad}
\newcommand{\dist}[1]{\lVert #1\rVert}
\newcommand{\fr}[1]{\{#1\}}
\newcommand{\Ahat}{\widehat A}
\newcommand{\ind}{\mathbf 1}
\newcommand{\cnt}{\#}
\newcommand{\LL}{\log\log N}
\newcommand{\decl}[1]{\nolinkurl{#1}}
% Immutable commit at which every Lean declaration cited below was checked.
\newcommand{\pincommit}{347ec3fe631a3e7ab2ae25d53508a3a4c9ea9fa8}
% \leanP{File}{line}{declaration}: link to NormalityChecker/File.lean#Lline at \pincommit.
\newcommand{\leanP}[3]{\leanat{\pincommit}{#1}{#2}{\decl{#3}}}

\begin{document}

\begin{center}
{\LARGE\bfseries Digit Distribution of Arithmetic Generating Functions:\\[3pt]
Euler Products, Theta Values and Sparse Coefficient Transfer\par}
\vspace{10pt}
\makebyline{1.0 (preprint draft)}{26 November 2026}
\end{center}

\begin{abstract}
We study the base-$b$ digits of values at $q=1/b$ of arithmetic generating
functions: Euler's product $E_b=\prod_{k\ge1}(1-b^{-k})$ and its square and
cube, theta values of positive definite binary quadratic forms, weight-one
modular forms with integer coefficients, and Lambert series of periodic weights.
The common tool is a finite coefficient-to-discrepancy transfer inequality, valid
in every integer radix, for integer coefficient sequences of controlled mass and
density-zero support. It places all of these numbers in the numerical
translation kernel $K_b=\{z:D_N(z;b)\to1\}$, where $D_N$ is the extreme
discrepancy of $(\{b^nz\})_{n<N}$. For $E_b$ we determine the digits exactly,
show that the words $0^L$ and $(b-1)^L$ each have frequency $1/2$, and prove
$1-D_N(E_b;b)=\sqrt{2/3}\,(\log_bN)/\sqrt N+O_b(N^{-1/2})$; the cube $E_b^3$ is
disjunctive, with constant $\sqrt2$. A sub-threshold width law gives exact
leading constants for multiplicative Lambert series, for $E_b^2$ and for theta
values of every primitive form of fundamental discriminant, in every radix, and
exhibits a discontinuity between the radices $2^k$ and all other radices. We
classify digit frequencies inside $K_b$: for $b\ge3$ no element of $K_b$ is
simply normal in base $b$, while $E_2$, $E_2^3$ and their $q$-derivatives are
simply normal in base $2$. A self-contained Selberg sieve shows that
$(a/d)E_b^2\in K_b$ only when $a/d\in\Z[1/b]$, and an entropy--multiplicity count
explains why sparse transfer cannot reach modular forms of weight at least two.
Lean~4 verifies the finite transfer core; the arithmetic arguments are written
proofs.
\end{abstract}

\section{Introduction}\label{sec:intro}

\subsection{Setting}\label{subsec:setting}

Fix an integer radix $b\ge2$. For a real number $z$ write $x_n=\fr{b^nz}$ for
the fractional parts of its orbit under multiplication by $b$, and let
\[
D_N(z;b)=\sup_{0\le s<t\le1}\Bigl|\frac{\cnt\{0\le n<N:x_n\in[s,t)\}}{N}-(t-s)\Bigr|
\]
be the extreme discrepancy of $x_0,\dots,x_{N-1}$. A number is normal in base $b$
exactly when $D_N(z;b)\to0$ \cite{KN}. At the other extreme, $D_N(z;b)\le1$
always, and we call
\[
K_b=\{z\in\R:\ D_N(z;b)\to1\}
\]
the \emph{numerical translation kernel}. The deficit $1-D_N(z;b)$ has a
translation meaning: for every $N\ge1$,
\begin{equation}\label{eq:sharp}
\sup_{x\in\R}\bigl|D_N(x+z;b)-D_N(x;b)\bigr|=1-D_N(z;b),
\end{equation}
with the supremum attained at $x=0$ (\cref{thm:sharp}). So $1-D_N(z;b)$ is the
worst change that translation by $z$ can cause in the discrepancy of any input,
and $K_b$ is the set of numbers whose addition asymptotically preserves the
discrepancy of every real number. The kernel is a subgroup of $\R$ and equals
$(b-1)^{-1}\Z+\{w:M_N(w)\to0\}$, where $M_N(w)=N^{-1}\sum_{n<N}\dist{b^nw}$ and
$\dist{\cdot}$ is the distance to the nearest integer \cite{DiscrepancyCompanion}.
Elements of $K_b$ are never normal, and the rationals in $K_b$ are exactly those
whose reduced denominator divides $b^j(b-1)$ for some $j$.

Throughout, for an integer sequence $a$ with $\sum_m|a(m)|b^{-m}<\infty$ we write
\[
F_b(a)=\sum_{m\ge0}a(m)b^{-m},\qquad S(X)=\cnt\{1\le m\le X:a(m)\ne0\},\qquad
A(X)=\sum_{1\le m\le X}|a(m)|,
\]
for the value, the support count and the mass. Many classical constants have
this form. With $E(q)=\prod_{k\ge1}(1-q^k)$ and $\Theta_3(q)=\sum_{j\in\Z}q^{j^2}$:
\begin{itemize}[nosep]
\item $E_b=E(1/b)$ has coefficients in $\{0,\pm1\}$ supported on the generalized
pentagonal numbers (Euler);
\item $E_b^3$ has coefficients $(-1)^k(2k+1)$ at the triangular numbers (Jacobi);
\item $L_\chi(b)=\sum_{m\ge1}\chi(m)/(b^m-1)=\sum_n(1*\chi)(n)b^{-n}$ for a
Dirichlet character $\chi$, and $\Theta_3(1/b)^2=1+4L_{\chi_{-4}}(b)$;
\item $\Theta_Q(1/b)=\sum_{(m,n)\in\Z^2}b^{-Q(m,n)}$ for a positive definite
integral binary quadratic form $Q$.
\end{itemize}
Here and below $\chi_{-4}$ is the nontrivial character modulo $4$, and $E_b^2$,
$E_b^3$ denote powers of the number $E_b$.

\emph{Series versus product.} Euler's pentagonal number theorem and Jacobi's
identity identify $E_b$ and $E_b^3$ with explicit sparse series. The Lean
declarations cited in this paper speak about these series, which they define
coefficient by coefficient; the identities themselves are classical
\cite{EulerE244,EulerE541,Jacobi1829,Bell2005} and are not formalized.
\Cref{subsec:seriesproduct} states precisely which parts are formal.

\subsection{Main results}\label{subsec:main}

Let $\rho_N=(1+S(N-1))/N$ and $\log_b^+y=\max(0,\log_by)$. For $X\ge0$ let
\[
\Ahat(X)=\sum_{m\ge1}|a(m)|\min\bigl(1,b^{X-m}\bigr),
\]
a smoothed mass that
equals $A(X)$ up to the far tail.

\begin{maintheorem}[Sparse coefficient transfer]\label{thm:A}\evidence{Lean-supported paper proof}
Let $b\ge2$ and let $a$ be an integer sequence with
$\sum_m|a(m)|b^{-m}<\infty$; put $F=F_b(a)$.
\begin{enumerate}[label=\textup{(\roman*)},nosep]
\item For all $N\ge1$, $L\ge0$ and $\varepsilon>0$,
\[
1-D_N(F;b)\le2\varepsilon+\frac{L\,(1+S(N-1))}{N}+\frac{\Ahat(N+L)}{(b-1)\,b^L\varepsilon N}.
\]
\item If $A(X)\le CX$ for every $X\ge1$, then for every $N\ge1$
\[
1-D_N(F;b)\le\rho_N\Bigl(2\log_b\frac1{\rho_N}+\log_b^+\frac{4C}{b-1}+3\Bigr),
\]
and the constant $2$ cannot be lowered.
\item $F\in K_b$ in each of the cases: $A(X)=O(X)$ and $S(X)=o(X)$;
$A(X)\le C(X+1)^k$ and $S(X)\log X=o(X)$; $A(X)\le CX(\log(X+2))^\gamma$ and
$S(X)\log\log X=o(X)$. Density-zero support alone does not suffice.
\item Removing from each coefficient its largest power of $b$ lowers the constant
$2$ to $1$ when the remaining parts have bounded mean, and under an isolation
hypothesis the bound becomes an asymptotic formula
(\cref{thm:displacedwidth,thm:widthlaw}).
\end{enumerate}
\end{maintheorem}

\begin{maintheorem}[Euler's product]\label{thm:B}\evidence{Lean-supported paper proof}
For every $b\ge2$, the base-$b$ expansion of $E_b$ is
$0.(b-2)(b-1)00B_1B_2B_3\cdots$ with
$B_j=1\,0^{2j}(b-1)^{4j}(b-2)(b-1)^{2j+1}0^{4j+2}$. For every $L\ge1$ the words $0^L$
and $(b-1)^L$ each have frequency $1/2$, with error $O(L\sqrt N+L^2)$ at every
prefix $N$, and every other word occurs $O(L\sqrt N+L^2)$ times. Hence $E_2$ is
simply normal in base $2$, $E_b$ is not simply normal in base $b$ for $b\ge3$, and
$E_b$ is never disjunctive. Moreover
\[
1-D_N(E_b;b)=\sqrt{\tfrac23}\,\frac{\log_bN}{\sqrt N}+O_b\Bigl(\frac1{\sqrt N}\Bigr),
\]
so $E_b\in K_b$, and by \eqref{eq:sharp} translation by $E_b$ changes the
discrepancy of any input by at most this amount.
\end{maintheorem}

\begin{maintheorem}[Cube, square and rational scalings]\label{thm:C}\evidence{Lean-supported paper proof}
Let $b\ge2$.
\begin{enumerate}[label=\textup{(\roman*)},nosep]
\item $E_b^3$ has explicit digits; the words $0^L$ and $(b-1)^L$ each have
frequency $1/2$, all other words frequency $0$; $E_b^3$ is disjunctive; and
$1-D_N(E_b^3;b)=\sqrt2\,(\log_bN)/\sqrt N+O_b(N^{-1/2})$.
\item $E_b^2\in K_b$ with $1-D_N(E_b^2;b)\ll_b(\LL)/\sqrt{\log N}$. Granting the
classical identification of $\eta(12\tau)^2$ as a Hecke eigenform with complex
multiplication \cite{Serre1985}, the coefficient of $b^{-n}$ in $E_b^2$ is the
coefficient $c(12n+1)$ of $\eta(12\tau)^2$; it is nonzero exactly when every prime
$p\not\equiv1\pmod{12}$ divides $12n+1$ to an even power; and the number of
nonzero coefficients up to $N$ is asymptotic to $C_\eta N/(\log N)^{3/4}$ with
$C_\eta=2.418529\ldots$ given in \eqref{eq:Ceta}.
\item If $P\in\Z[1/b][x]$ has degree at most $3$, then $P(E_b)\in K_b$. For
$r\in\{1,2,3\}$ and a reduced fraction $a/d$, $(a/d)E_b^r\in K_b$ if and only if
$a/d\in\Z[1/b]$; for $r=2$ and odd prime factors of $d$ this is proved granting the
eigenform identification of \textup{(ii)}, and rests on a uniform pair sieve for the
support of $E_b^2$ (\cref{thm:S}), proved here in full.
\end{enumerate}
\end{maintheorem}

For the next theorem write $LL=\log\log N$, let $C_\chi$ be the constant in
$\cnt\{n\le N:(1*\chi)(n)>0\}\sim C_\chi N/\sqrt{\log N}$, and for a radix
$b=2^sb'$ with $b'$ odd put $\beta_b=\log b/\log b'\in[1,\infty]$.

\begin{maintheorem}[Exact discrepancy constants in every radix]\label{thm:D}\evidence{Lean-supported paper proof}
\leavevmode
\begin{enumerate}[label=\textup{(\roman*)},nosep]
\item For every real nonprincipal Dirichlet character $\chi$,
\[
1-D_N(L_\chi(b);b)\sim
\begin{cases}
\dfrac{C_\chi}{2\log b}\,\dfrac{LL}{\sqrt{\log N}}, & b=2^k,\\[2ex]
\dfrac{C_\chi(1+\log2)}{2\log b}\,\dfrac{LL}{\sqrt{\log N}}, & b\text{ not a power of }2.
\end{cases}
\]
For $b=2$ and $\chi=\chi_{-4}$ the constant is $K/(2\log2)$ with $K=0.76422365\ldots$
the Landau--Ramanujan constant, and the same asymptotic holds for
$\Theta_3(1/2)^2=1+4L_{\chi_{-4}}(2)$.
\item Granting the eigenform identification of \cref{thm:C}(ii),
$1-D_N(E_b^2;b)\sim\frac{3C_\eta}{4\log b}\frac{LL}{(\log N)^{3/4}}$ for $b=2^k$ and
$\sim\frac{C_\eta(3+\log2)}{4\log b}\frac{LL}{(\log N)^{3/4}}$ otherwise.
\item Let $Q$ be a positive definite integral binary quadratic form and
$\Psi_N=C_QLL/(2\log b\sqrt{\log N})$, where
$\cnt\{n\le N:r_Q(n)>0\}\sim C_QN/\sqrt{\log N}$. Then
$\liminf_N(1-D_N(\Theta_Q(1/b);b))/\Psi_N\ge1$ for every $Q$ and $b$. If $Q$ is
primitive of fundamental discriminant $\Delta<0$, then
$1-D_N\sim(1+\log2)\Psi_N$ for every $b$ that is not a power of $2$, and
$1-D_N\sim\bigl(1+\log2\cdot\mathbb E_\Delta[1-1/d]\bigr)\Psi_N$ for $b=2^k$, where
$\mathbb E_\Delta[1-1/d]$ is an explicit Euler product depending only on the odd part
of the class group (\cref{thm:classgroups}); it vanishes when $\Cl(\Delta)$ is a
$2$-group.
\item For multiplicative $c$ with $0\le c(p^e)\le e+1$ and $c(p)=2$ exactly on a set
of primes of natural density $\delta\in(0,1)$ (\cref{subsec:primedensity}), the
deficit of $\sum_nc(n)b^{-n}$ is asymptotic to $(1-\delta)\Phi_N$ for $b=2^k$ and
to $(1-\delta+\delta\log2)\Phi_N$ for every other $b$ provided
$\delta\le\beta_b/(\beta_b+\log2)$; here $\Phi_N$ is the normalized support density
\eqref{eq:PhiN}. For larger $\delta$ explicit unconditional bounds hold.
\end{enumerate}
\end{maintheorem}

\begin{maintheorem}[Theta values, weight one and periodic weights]\label{thm:E}\evidence{Lean-supported paper proof}
Let $b\ge2$.
\begin{enumerate}[label=\textup{(\roman*)},nosep]
\item For every positive definite integral binary form $Q$,
$\Theta_Q(1/b)\in K_b$ with $1-D_N\ll_{Q,b}LL/\sqrt{\log N}$; the same holds for
integral shifted lattices, congruence subsums, bounded integer weights and
polynomials of total degree at most $2$ in finitely many unary quadratic
series.
\item Every holomorphic modular form $f=\sum_{n\ge0}a(n)q^n$ of weight one on some
$\Gamma_1(M)$ with all $a(n)\in\Z$ satisfies $f(1/b)\in K_b$ with
$1-D_N\ll_{f,b}LL/(\log N)^{\alpha_f}$, where $\alpha_f>0$ is defined in
\cref{thm:weightone}.
\item Every integer-valued combination of dilated divisor sums $1*\chi$ of
characters of even order lies in $K_b$; and for every odd $M$-periodic integer
weight $w$, the Lambert series $\sum_mw(m)/(b^m-1)$ lies in $K_b$. In particular,
for $q=b^{-M}$ and $1\le r<M$ the normalized $q$-digamma reflection
$(\psi_q(r/M)-\psi_q(1-r/M))/\log q$ lies in $K_b$.
\end{enumerate}
\end{maintheorem}

\begin{maintheorem}[Digit frequencies in the kernel]\label{thm:F}\evidence{Lean-supported paper proof}
Let $b\ge2$.
\begin{enumerate}[label=\textup{(\roman*)},nosep]
\item For $z\in K_b$, either some interior constant word $m^L$ ($1\le m\le b-2$) has
frequency tending to $1$ for every $L$, or the limit points of the word
distribution of length $L$ are exactly the measures
$\alpha\delta_{0^L}+(1-\alpha)\delta_{(b-1)^L}$ with $\alpha$ in a closed interval.
For $b\ge3$ no element of $K_b$ is simply normal in base $b$, and no element of
$K_b$ is normal in base $b$.
\item Every series $F_b(a)$ with integers $a\ge0$ that satisfies one of the three
hypotheses of \cref{thm:A}(iii) has $f_{0^L}(N)\to1$ for every $L$; this applies to $\Theta_Q(1/b)$ and $L_\chi(b)$ for real
nonprincipal $\chi$, which are therefore not simply normal in any base.
\item With $D=q\,d/dq$, the numbers $D^jE_b$, $D^jE_b^3$ $(j\ge0)$ and $E_b^2DE_b$
have $f_{0^L},f_{(b-1)^L}\to\frac12$ with error $O_{b,j}((L+\log N)/\sqrt N)$; they
are simply normal in base $2$ and not simply normal in base $b\ge3$; and
$D^jE_b$ $(j\ge1)$, $D^jE_b^3$ and $E_b^2DE_b$ are disjunctive.
\item $E_b^2$ is not simply normal for $b\ge3$, and, granting the eigenform
identification, both $0^L$ and $(b-1)^L$ have positive lower frequency in $E_b^2$.
\item For $d\ge2$ prime to $b$ and $a\ne0$ with $\gcd(a,d)=1$, the words of
$(a/d)E_b$ and $(a/d)E_b^3$ have exact rational limiting frequencies computed by a
finite automaton. In base $2$ these numbers are simply normal whenever $\ord_d(2)$
is odd; they are never normal.
\end{enumerate}
\end{maintheorem}

\begin{maintheorem}[Limits of sparse and signed transfer]\label{thm:G}\evidence{Lean-supported paper proof}
\leavevmode
\begin{enumerate}[label=\textup{(\roman*)},nosep]
\item $z\in K_b$ if and only if $z-c-\sum_{t\ge1}A(t)b^{-t}\in\Z$ for some
$c\in(b-1)^{-1}\Z$ and some integers $A(t)$ with $|A(t)|\le\lfloor b/2\rfloor$ and
density-zero support. Hence signed-digit transfer with unrestricted
cancellation between coefficients is equivalent to membership in $K_b$.
\item For $E_b^4$, $E_b^6$, $E_b^8$ and $D(\Theta_3(1/b)^2)$ the minimal signed-digit
width of the coefficients up to $N$ is of exact order $N\sqrt{\log N}$, and of
exact order $N(\log N)^{1/4}$ for $D(E_b^2)$, in every radix (granting the classical
CM descriptions of the forms involved, \cref{thm:etainput,thm:cminput}); every decomposition of
the individual coefficients into boundedly many signed digits uses
$\gg N\sqrt{\log N}$ digits for the first four. No distributional statement about
these numbers follows.
\item $K_b$ is not closed under squaring, and modulo $K_b$ one has
$\Theta_3(1/b)^4\equiv16L_{\chi_{-4}}(b)^2$ and
$\Theta_3(1/b)^3\equiv4\Theta_3(1/b)L_{\chi_{-4}}(b)$.
\end{enumerate}
\end{maintheorem}

\subsection{Relation to earlier work}\label{subsec:literature}

The identities used are classical: Euler's pentagonal number theorem, proved by
Euler in a letter to Goldbach of 1750 and published in
\cite{EulerE244,EulerE541} (see \cite{Bell2005} for the history), and Jacobi's
cube identity \cite{Jacobi1829}. The counting functions of integers represented
by binary forms are due to Landau \cite{Landau1908} and Bernays \cite{Bernays1912}
(see \cite{BMO2011}); lacunarity of powers of $\eta$ and the Chebotarev
description of weight-one forms are due to Serre \cite{Serre1981,Serre1985} and
Deligne--Serre \cite{DeligneSerre1974}. Transcendence of $E_b$, of $\Theta_Q(1/b)$
and of $f(1/b)$ in \cref{thm:E}(ii) follows from Nesterenko's theorem
\cite{Nesterenko1996} through \cite{DNNS1996,Bertrand1997,HamiehMurty2016}; we
claim nothing new about transcendence.

For digits, the closest published result is due to Bailey, Borwein, Crandall and
Pomerance \cite[\S9]{BBCP2004}: the number of ones among the first $M$ binary digits
of $\Theta_3(1/2)^2$ is asymptotic to $C_0M/\sqrt{\log M}$ with $C_0\approx0.7996$,
and they expect $\Theta_3(1/2)^3$ to be normal. Their isolated-coefficient argument
is close to ours. They do not consider discrepancy, other radices, $E_b$ or its
powers. A companion paper \cite{MultLambert} proved
$1-D_N(L;2)\asymp(\LL)/(\log N)^{1-\delta}$ for the binary multiplicative Lambert
series of \cref{thm:D}(iv), together with disjunctivity; \cref{thm:D} replaces
the order by an exact constant and extends it to all radices. Rauzy
\cite{Rauzy1976} characterized the numbers whose addition preserves normality, and
Bergelson and Downarowicz \cite{BD2025} studied preservation of normality and
determinism under arithmetic operations; multiplication by nonzero rationals
preserves normality \cite{Wall1949}. Those results concern normal inputs, whereas
$K_b$ is defined through the numerical discrepancy of every input
\cite{DiscrepancyCompanion}, and its elements are far from normal. We located no
source that states the digit distribution of $E_b$, $E_b^2$ or $E_b^3$, an exact
discrepancy constant for any of the numbers above, or the radix dependence in
\cref{thm:D}; this records a failed search, not a proof of novelty. Recent
preprints on related constants include \cite{Campbell2026} (the word $11$ in the
Erd\H{o}s--Borwein constant) and \cite{BecherLew2026} (numbers that are normal in
odd bases and deterministic in base $2$).

\subsection{What is not claimed}\label{subsec:notclaimed}

No classical constant is proved normal or non-normal here beyond the numbers named
in the theorems. Membership in $K_b$, normality and even $D^+<1$ remain open for
the partition constant $P_b=1/E_b$, for $E_b^r$ with $r\ge4$, for
$\Theta_3(1/b)^3$ and $\Theta_3(1/b)^4$, for the Eisenstein values
$E_2(1/b),E_4(1/b),E_6(1/b)$ and for $\Delta(1/b)$. \Cref{thm:G}(ii) is a statement
about the method, not about these numbers. The exact constants of
\cref{thm:D}(iv) above the threshold $\beta_b/(\beta_b+\log2)$, and of
\cref{thm:D}(iii) for non-fundamental or imprimitive forms beyond the lower bound,
remain open, as do the simple normality of $E_2^2$ in base $2$ and the converse of
\cref{thm:F}(v). All exact constants in \cref{thm:D} have relative error of order
$(\log\log\log N)/\log\log N$ or worse, so finite computations, which reach
$\log\log N<2.7$, do not test them.

\subsection{Organization and evidence}\label{subsec:evidence}

\Cref{sec:prelim} collects finite tools. \Cref{sec:transfer} proves the transfer
theorems and the sub-threshold width law. \Cref{sec:euler} treats $E_b$, $E_b^2$,
$E_b^3$ and polynomials in them. \Cref{sec:radix} proves the exact constants of
\cref{thm:D}. \Cref{sec:theta} treats theta values, weight-one forms and periodic
weights. \Cref{sec:digits} classifies digit frequencies inside $K_b$ and proves
\cref{thm:F}. \Cref{sec:sieve} proves the uniform pair sieve for $E_b^2$ and the
rational-scaling criterion. \Cref{sec:limits} proves \cref{thm:G}, and
\cref{sec:open} lists open problems.

Each numbered statement carries one evidence label.
\evidence{Kernel-verified}: the complete statement is a Lean~4 declaration checked
by the kernel, linked at the pinned commit
\texttt{347ec3fe631a3e7ab2ae25d53508a3a4c9ea9fa8}.
\evidence{Lean-supported paper proof}: a written proof whose finite core is such a
declaration. \evidence{Paper proof}: a written proof only.
\evidence{External input}: a published theorem quoted without proof. Finite
computations are reported only as illustrations; none is used in a proof. The
written proofs have been checked by AI models from several providers, which is not
independent human peer review (\cref{sec:ai}).

\section{Finite tools}\label{sec:prelim}

Write $\T=\R/\Z$ and $\dist t$ for the distance from $t$ to $\Z$. For $\varepsilon\ge0$
let
\[
E_N(\varepsilon)=E_N(z,\varepsilon)=\cnt\{0\le n<N:\ \dist{b^nz}>\varepsilon\}.
\]
A half-open arc of $\T$ of length $\ell\in[0,1]$ is a set $J=\{t:\fr{t-\alpha}<\ell\}$,
and $\mu_N(J)$ is the proportion of $n<N$ with $x_n\in J$. All logarithms are natural
unless a base is shown.

\begin{lemma}[Arcs and concentration]\label{lem:arcs}\evidence{Lean-supported paper proof}
For every half-open arc $J$ of length $\ell$, $\mu_N(J)-\ell\le D_N(z;b)$.
Conversely every interval deviation in the definition of $D_N$ equals
$\mu_N(J)-\ell$ for some arc $J$. Hence
\[
1-D_N(z;b)=\inf_J\bigl(1-\mu_N(J)+\ell\bigr).
\]
\end{lemma}

\begin{proof}
The first claim is \leanP{SharpTranslation}{123}{arc_le}; the second is
\leanP{SharpTranslation}{167}{interval_eq_arc_excess} (the complement of an interval with
negative deviation is an arc with the same positive excess). The identity follows.
\end{proof}

\begin{lemma}[Deficit bound]\label{lem:deficit}\evidence{Kernel-verified}
For every $\varepsilon\ge0$ and $N\ge1$, $1-D_N(z;b)\le2\varepsilon+E_N(\varepsilon)/N$.
\end{lemma}

\begin{proof}
This is \leanP{Phase18Synthetic}{524}{discrepancy_deficit_le}; the arc of length
$1-2\varepsilon$ avoiding the closed $\varepsilon$-neighbourhood of $0$ contains only far
points.
\end{proof}

\begin{theorem}[Sharp translation identity {\cite{DiscrepancyCompanion}}]\label{thm:sharp}\evidence{Kernel-verified}
For every real $z$ and $N\ge1$,
$\sup_x|D_N(x+z;b)-D_N(x;b)|=1-D_N(z;b)$, attained at $x=0$.
\end{theorem}

\begin{proof}
\leanP{SharpTranslation}{386}{sharp_translation_isGreatest}.
\end{proof}

By \cref{thm:sharp}, every upper bound for $1-D_N(z;b)$ proved below is also a bound for
the worst translation error, and every lower bound shows that this error is attained.
We do not repeat this remark.

\begin{lemma}[Separation]\label{lem:separation}\evidence{Paper proof}
For every $\varepsilon\ge0$ and $N\ge1$,
\[
1-D_N(z;b)\ge\min\Bigl\{1-\frac{E_N(\varepsilon)}N,\ \inf_{\ell\ge0}\Bigl(\ell+\frac{E_N(\ell+\varepsilon)}N\Bigr)\Bigr\}.
\]
\end{lemma}

\begin{proof}
Let $J$ be a half-open arc of length $\ell$. If $J$ misses the closed ball
$\{\dist t\le\varepsilon\}$, every sample in $J$ is far, so $1-\mu_N(J)+\ell\ge1-E_N(\varepsilon)/N$.
Otherwise $J$ contains a point $y$ with $\dist y\le\varepsilon$, and every $t\in J$ has
$\dist t<\ell+\varepsilon$; samples with $\dist{x_n}>\ell+\varepsilon$ lie outside $J$, so
$1-\mu_N(J)+\ell\ge\ell+E_N(\ell+\varepsilon)/N$. Take the infimum over $J$ and use
\cref{lem:arcs}.
\end{proof}

\begin{lemma}[Circle separation]\label{lem:circle}\evidence{Kernel-verified}
Let $0\le\zeta$ and $0\le\Delta\le\eta-\zeta$. If at least $\Delta N$ indices $n<N$ satisfy
$\dist{x_n}\le\zeta$ and at least $\Delta N$ satisfy $\dist{x_n}>\eta$, then
$1-D_N(z;b)\ge\Delta$.
\end{lemma}

\begin{proof}
\leanP{CoefficientTransfer}{686}{one_sub_discrepancy_ge_of_bulk_and_far} states
$\min\bigl(\eta-\zeta,\ \cnt\mathrm{far}/N,\ \cnt\mathrm{bulk}/N\bigr)\le1-D_N$ in every radix.
\end{proof}

\begin{lemma}[Linear transfer]\label{lem:linear}\evidence{Kernel-verified}
If $b^kx-y=\sum_ic_iz_i$ is a finite combination with integers $c_i$, then for every
$N\ge1$
\[
|D_N(x;b)-D_N(y;b)|\le\frac kN+\sum_i|c_i|\bigl(1-D_N(z_i;b)\bigr).
\]
Moreover $D_N(z+m;b)=D_N(z;b)$ for $m\in\Z$.
\end{lemma}

\begin{proof}
\leanP{CharacterLambertTransfer}{175}{radix_pow_linear_transfer}; integer translation
invariance is \leanP{Phase18Synthetic}{681}{discrepancy_add_int}.
\end{proof}

Letting $N\to\infty$ in \cref{lem:linear} shows that $K_b$ is a $\Z[1/b]$-module; this
deduction is written, not formal.

\section{Sparse coefficient transfer}\label{sec:transfer}

Throughout this section $a$ is an integer sequence with $\sum|a(m)|b^{-m}<\infty$,
$F=F_b(a)$, and $x_n=\fr{b^nF}$. The term $a(0)$ is an integer and plays no role.

\subsection{The finite transfer inequality}\label{subsec:finite}

\begin{lemma}[Integral terms drop out]\label{lem:integral}\evidence{Paper proof}
Fix $n\ge0$ and let $I$ be a set of indices $m>n$ with $b^{m-n}\mid a(m)$. Then
$\dist{b^nF}\le\sum_{m>n,\ m\notin I}|a(m)|b^{n-m}$. In particular, if
$a(n+1)=\dots=a(n+L)=0$, then $\dist{b^nF}\le\tau_L(n):=\sum_{j>L}|a(n+j)|b^{-j}$.
\end{lemma}

\begin{proof}
Every term with $m\le n$ or $m\in I$ contributes an integer to $b^nF$, and a summable
family of integers has finitely many nonzero terms.
\end{proof}

\begin{lemma}[Shadows]\label{lem:shadow}\evidence{Paper proof}
Let $W_L(N)$ be the number of $0\le n<N$ whose next occupied index lies in $(n,n+L]$.
Then $W_L(N)\le L(1+S(N-1))$.
\end{lemma}

\begin{proof}
The shifts whose next occupied index is $m$ form an interval of length at most $L$ ending
at $m-1$. Occupied indices $m\ge N$ meet $[0,N)$ only through the first of them, which
shadows at most $L$ shifts.
\end{proof}

\begin{lemma}[Mean tail]\label{lem:meantail}\evidence{Paper proof}
For $N\ge1$ and $L\ge0$,
$\sum_{n<N}\tau_L(n)\le\Ahat(N+L)/((b-1)b^L)$, where
$\Ahat(X)=\sum_{m\ge1}|a(m)|\min(1,b^{X-m})$, and for integers $X\ge0$,
$\Ahat(X)=(1-b^{-1})\sum_{i\ge0}b^{-i}A(X+i)$. If $A(X)\le CX$ for all $X\ge1$, then $\Ahat(X)\le C(X+1/(b-1))$.
\end{lemma}

\begin{proof}
Exchanging the sums, $a(m)$ is paired with the shifts $0\le n\le\min(N-1,m-L-1)$, and the
geometric sum of $b^{n-m}$ over them is at most $b^{-\max(L,m-N)}/(b-1)$. The second
formula for $\Ahat$ is Abel summation, and the last bound follows from it.
\end{proof}

\begin{theorem}[Finite all-prefix transfer]\label{thm:finitetransfer}\evidence{Paper proof}
For all $N\ge1$, integers $L\ge0$ and $\varepsilon>0$,
\[
E_N(\varepsilon)\le L(1+S(N-1))+\frac{\Ahat(N+L)}{(b-1)b^L\varepsilon},\qquad
1-D_N(F;b)\le2\varepsilon+\frac{L(1+S(N-1))}N+\frac{\Ahat(N+L)}{(b-1)b^L\varepsilon N}.
\]
\end{theorem}

\begin{proof}
Shadowed shifts are counted by \cref{lem:shadow}; at every other shift
$\dist{b^nF}\le\tau_L(n)$ by \cref{lem:integral}. Markov's inequality and
\cref{lem:meantail} bound the number of $n<N$ with $\tau_L(n)>\varepsilon$. Finish with
\cref{lem:deficit}.
\end{proof}

The Lean library contains a version with a slightly weaker shadow term.

\begin{theorem}[Lean form of the transfer]\label{thm:leantransfer}\evidence{Kernel-verified}
If $A(X)\le CX$ for all $X$, then for all $N\ge1$, $L\ge0$ and $\varepsilon>0$,
\[
1-D_N(F;b)\le2\varepsilon+\frac{L\,S(N+L)}N+\frac{C(N+L+2)}{N\varepsilon b^L},
\]
and the same bound holds for $|D_N(x+F;b)-D_N(x;b)|$ for every real $x$. Consequently
$A(X)=O(X)$ and $S(X)=o(X)$ imply $F\in K_b$ and $D_N(x+F;b)-D_N(x;b)\to0$ for every
$x$. If instead $A(X)\le C(X+1)^k$ with an integer $k\ge0$ and $S(X)\log X=o(X)$, or
$A(X)\le CX(\log(X+2))^\gamma$ with an integer $\gamma\ge0$ and
$S(X)\log\log(X+3)=o(X)$, then again $F\in K_b$.
\end{theorem}

\begin{proof}
\leanP{CoefficientTransfer}{392}{finite_sparse_mass_transfer},
\leanP{CoefficientTransfer}{423}{finite_sparse_mass_translation},
\leanP{CoefficientTransfer}{434}{linear_mass_sparse_tendsto_one},
\leanP{CoefficientTransfer}{512}{linear_mass_sparse_preserves_discrepancy},
\leanP{GeneralMassTransfer}{253}{finite_poly_mass_transfer},
\leanP{GeneralMassTransfer}{299}{poly_mass_sparse_tendsto_one},
\leanP{GeneralMassTransfer}{657}{finite_log_mass_transfer} and
\leanP{GeneralMassTransfer}{693}{log_mass_sparse_tendsto_one}. A real exponent $\gamma$ may
be replaced by $\lceil\gamma\rceil$.
\end{proof}

\begin{theorem}[Linear mass: explicit rate]\label{thm:linearrate}\evidence{Paper proof}
If $A(X)\le CX$ for every $X\ge1$, then for every $N\ge1$, with
$\rho_N=(1+S(N-1))/N$,
\[
1-D_N(F;b)\le\rho_N\Bigl(2\log_b\frac1{\rho_N}+\log_b^+\frac{4C}{b-1}+3\Bigr).
\]
\end{theorem}

\begin{proof}
Write $\rho=\rho_N\in(0,1]$. If $\rho=1$ the right side is at least $3$. Otherwise put
$\varepsilon=\rho/2$ and $L=\max(0,\lceil\log_b(4C/((b-1)\rho^2))\rceil)$. If
$L+1/(b-1)\le N$, then $\Ahat(N+L)\le2CN$ by \cref{lem:meantail}, so the last term of
\cref{thm:finitetransfer} is at most $4Cb^{-L}/((b-1)\rho)\le\rho$ and
$1-D_N\le(L+2)\rho$. Otherwise $L\ge N$, so $(L+2)\rho>1$ and the bound is trivial.
Finally $L+2$ is at most the bracket.
\end{proof}

\begin{proposition}[The constant $2$ is sharp]\label{prop:sharp2}\evidence{Paper proof}
Let $K\asymp q$ be the integer with base-$b$ digits $1,0,1,\dots,0,1$ and
$z=K/(b^q-1)=\sum_{k\ge1}Kb^{-qk}$. Then $A(X)\le(K/q)X$, $\rho_q=1/q$, and
$(1-D_q(z;b))\,q=2\log_bq+O_b(\log\log q)$.
\end{proposition}

\begin{proof}
The orbit has period $q$. Let $w$ be the number of digits of $K$. For $1\le j<w$ the point
$x_{q-j}$ is a digit suffix of $K$ times $b^{-j}$ up to $O(b^{-q})$, and since the digits
alternate it lies at distance at least $c_b>0$ from $0$. For $j\ge w$ the point is
$Kb^{-j}(1+O(b^{-q}))$, at distance at least $\eta$ from $0$ for
$j\le w+\log_b(1/\eta)-O(1)$ and at most $\zeta$ for $j\ge w+\log_b(1/\zeta)+O(1)$. Take
$\zeta=1/q$ and $\eta=6\log_bq/q$ in \cref{lem:circle}: this gives
$1-D_q\ge(w+\log_b(1/\eta)-O(1))/q=(2\log_bq-O(\log\log q))/q$. \Cref{thm:linearrate} gives
the upper bound.
\end{proof}

Concatenating such periodic blocks on rapidly growing ranges gives a single number with
$\rho_N\to0$ for which the constant $2$ is attained along a subsequence.

\begin{theorem}[Larger mass]\label{thm:largermass}\evidence{Paper proof}
For every $N\ge1$, with $\rho=\rho_N$ and $\mu_N=\max(1,\Ahat(2N)/N)$,
\[
1-D_N(F;b)\le\rho\Bigl(2\log_b\frac1\rho+\log_b\mu_N+\log_b^+\frac2{b-1}+3\Bigr).
\]
If $S(X)\le C_1X(\log X)^{-\kappa}$ for $X\ge3$ with $\kappa>0$, and
$A(X)\ll X(\log X)^{\beta}$ with $\beta\ge0$, then
$1-D_N(F;b)\le(2\kappa+\beta+o(1))\,C_1\log_b\log N/(\log N)^\kappa$.
\end{theorem}

\begin{proof}
As for \cref{thm:linearrate} with $L=\max(0,\lceil\log_b(2\mu_N/((b-1)\rho^2))\rceil)$ and
$\Ahat(N+L)\le\Ahat(2N)\le\mu_NN$ when $L\le N$. Under the second hypothesis
$\mu_N\ll(\log N)^\beta$, and $\rho\mapsto\rho\log(1/\rho)$ increases on $(0,1/e)$, so the
upper bound for $\rho$ may be inserted.
\end{proof}

\begin{proposition}[Density zero is not enough]\label{prop:counterexamples}\evidence{Paper proof}
Let $0=m_0<m_1<\cdots$.
\begin{enumerate}[label=\textup{(\alph*)},nosep]
\item If all $m_k$ are even and $a(m_k)=(b^{m_k-m_{k-1}}-1)/(b^2-1)$, with $a=0$ elsewhere,
then $F=1/(b^2-1)\notin K_b$, while $S(X)$ may grow as slowly as $\log_2X$.
\item If $\xi=\sum_i\xi_ib^{-i}$ is normal in base $b$ and
$a(m_k)=\sum_{m_{k-1}<i\le m_k}\xi_ib^{m_k-i}$, then $F=\xi$; for $0<\gamma\le1$ and
$m_k-m_{k-1}=\max(1,\lfloor\gamma\log_b(m_{k-1}+1)\rfloor)$ one has $|a(m)|\le bm^{\gamma}$,
$S(X)\sim X/(\gamma\log_bX)$ and $A(X)\ll X^{1+\gamma}$.
\end{enumerate}
\end{proposition}

\begin{proof}
(a) $\sum_ka(m_k)b^{-m_k}=\sum_k\sum_{m_{k-1}<2i\le m_k}b^{-2i}=\sum_{i\ge1}b^{-2i}$, and the
reduced denominator $b^2-1$ does not divide $b^j(b-1)$. (b) The blocks reproduce the digits
of $\xi$; the bounds are immediate. Normal numbers satisfy $D_N\to0$.
\end{proof}

In both examples the digit widths of the coefficients fill the gaps between occupied
indices. The width form below shows that this is the only obstruction.

\subsection{Displaced sources}\label{subsec:displaced}

For $a(m)\ne0$ write $a(m)=b^{V(m)}U(m)$ with $b\nmid U(m)$, let $u(m)$ be the number of
base-$b$ digits of $|U(m)|$, and put $t(m)=m-V(m)$ (the \emph{arrival}) and
$\tilde t(m)=t(m)-u(m)$. The radix may be composite; then $U(m)$ may still be divisible by
a prime factor of $b$, but its last base-$b$ digit is nonzero.

\begin{lemma}[Radix factorization]\label{lem:radixval}\evidence{Kernel-verified}
In every radix $b\ge2$, every nonzero integer $a$ factors as $b^{V}U$ with
$V=\mathtt{radixVal}(b,a)\ge0$ and $b\nmid U$, and $b^n(a/b^m)$ is an integer if $m\le n+V$ and equals $U/b^{m-V-n}$ if
$n+V<m$.
\end{lemma}

\begin{proof}
\leanP{CoefficientTransfer}{588}{radixVal_mul_radixUnit} with \decl{not_dvd_radixUnit},
and \leanP{CoefficientTransfer}{617}{displaced_term_eq}.
\end{proof}

The factorization is unique, since $V$ must be the exact power of $b$ dividing $a$; this
remark is not part of the Lean statement.

\begin{theorem}[Displaced reindexing]\label{thm:reindex}\evidence{Lean-supported paper proof}
Each fibre $\{m:t(m)=t\}$ is finite and only finitely many $m$ have $t(m)\le0$. With
$C(t)=\sum_{t(m)=t}U(m)\in\Z$,
\[
F=\sum_{t\in\Z}C(t)b^{-t},\qquad x_n=\Bigl\{\sum_{d\ge1}C(n+d)b^{-d}\Bigr\}\quad(n\ge0).
\]
\end{theorem}

\begin{proof}
Since $|U|\ge1$, a term with $t(m)\le0$ has $|a(m)|b^{-m}\ge1$, so there are finitely many,
and the size of a fibre times $b^{-t}$ is at most $\sum|a(m)|b^{-m}$. Absolute convergence
permits the rearrangement, and at shift $n$ the terms with $t\le n$ are integers by
\cref{lem:radixval}.
\end{proof}

\begin{theorem}[Displaced width inequality]\label{thm:displacedwidth}\evidence{Lean-supported paper proof}
For $N\ge1$, $L\ge0$ and $\varepsilon>0$,
\[
E_N(\varepsilon)\le W_U(N)+L(1+P(N))+\Lambda_N+\frac{\widehat P(N+L)}{(b-1)b^L\varepsilon},
\]
where $P(N)=\cnt\{m:1\le t(m)\le N\}$, $W_U(N)=\sum_{1\le t(m)\le N}u(m)$,
$\Lambda_N=\max_{t(m)>N}(u(m)-t(m)+N)^+$ and
$\widehat P(X)=\sum_{a(m)\ne0}|U(m)|b^{-u(m)}\min(1,b^{X-\tilde t(m)})$.
Suppose moreover that $P(N)\to\infty$, $\pi_N=P(N)/N\to0$,
$\sum_{1\le t(m)\le X}|U(m)|\le K_0P(X)$ for all $X$, $\Lambda_N=o(P(N)\log_b(1/\pi_N))$, and
that for some $\eta_N\to0$, with $L_N=\lceil(1+\eta_N)\log_b(1/\pi_N)\rceil$,
$\widehat P(N+L_N)\le P(N)\pi_N^{-\eta_N}$. Then
$1-D_N(F;b)\le(1+o(1))\pi_N\log_b(1/\pi_N)$.
\end{theorem}

\begin{proof}
At shift $n$ a source with $t(m)\le n$ contributes an integer (\cref{lem:radixval}). If $n$
lies in no interval $[\tilde t(m)-L,t(m)-1]$, every source with $t(m)>n$ has
$\tilde t(m)>n+L$ and contributes $|U|b^{n-t}<b^{n-\tilde t}\le b^{-L-1}$. The union of the
intervals inside $[0,N)$ has size at most $W_U(N)+L(1+P(N))+\Lambda_N$, since intervals of
sources with $t(m)>N$ end at $N-1$. On the remaining shifts apply Markov's inequality to the
truncated tail as in \cref{lem:meantail}.

For the second claim, Jensen's inequality (concavity of the logarithm) gives
$W_U(N)\le P(N)(1+\log_bK_0)=O(P(N))$. Take $\varepsilon=\pi_N$ and $L=L_N$. Since
$b^{-L_N}\le\pi_N^{1+\eta_N}$, the Markov term is at most
$P(N)\pi_N^{-\eta_N}\pi_N^{1+\eta_N}/((b-1)\pi_N)=P(N)/(b-1)$. Dividing by $N$ and using
\cref{lem:deficit},
\[
1-D_N\le2\pi_N+O(\pi_N)+(1+\eta_N)\log_b\frac1{\pi_N}\cdot\pi_N\Bigl(1+\frac1{P(N)}\Bigr)
+O(\pi_N)+\frac{\Lambda_N}N,
\]
and each term other than the third is $o(\pi_N\log_b(1/\pi_N))$ because
$\log(1/\pi_N)\to\infty$, $P(N)\to\infty$ and $\Lambda_N/N=o(\pi_N\log_b(1/\pi_N))$. The
Lean declaration
\leanP{RadixWidthLaw}{466}{width_law_upper} checks the final arithmetic step from the
far-count bound.
\end{proof}

\subsection{The sub-threshold width law}\label{subsec:widthlaw}

By \cref{thm:reindex} we may work with the merged sequence $C$. For a source $t$ (that is,
$C(t)\ne0$) let $g_t$ be the gap to the previous source, $A(t)=\sum_{d\ge1}C(t+d)b^{-d}$
its tail and $Y(t)=C(t)+A(t)$ its local value.

\begin{lemma}[Source decomposition]\label{lem:source}\evidence{Paper proof}
Assume that there are infinitely many sources $t\ge1$, and for the first of them, $t_1$,
put $g_{t_1}=t_1$. The map $n\mapsto(t,j)$, with $t$ the first source after $n$ and $j=t-n$, is a bijection
from $\N$ onto $\{(t,j):1\le j\le g_t\}$, and $x_n=\fr{b^{-j}Y(t)}$. If $(t,t+h]$ contains
no source, then $A(t)=b^{-h}A(t+h)$.
\end{lemma}

\begin{proof}
At shift $n=t-j$ all sources $\le n$ are integral; the others are $t$ and its successors.
\end{proof}

\begin{lemma}[One source]\label{lem:onesource}\evidence{Lean-supported paper proof}
Suppose $|A(t)|\le\frac12$ and let $u$ be the number of digits of $|C(t)|$. If $b\nmid C(t)$,
then $\dist{b^{-j}Y(t)}\ge b^{-j}/2$ for every $j\ge1$. If $j\ge u+1$, then
$\dist{b^{-j}Y(t)}=b^{-j}|Y(t)|\ge b^{-j}(b^{u-1}-\frac12)$. Both statements hold for either
sign of $C(t)$.
\end{lemma}

\begin{proof}
If $b\nmid C(t)$ then $C(t)\not\equiv0\pmod{b^j}$, so $C(t)+A(t)$ is at distance at least
$\frac12$ from $b^j\Z$. For $j\ge u+1$, $|Y(t)|\le b^u-\frac12<b^{j}/2$. The two statements
are \leanP{RadixWidthLaw}{69}{dist0_div_pow_ge} and
\leanP{RadixWidthLaw}{102}{dist0_div_pow_eq}; the resulting count of at least
$\min(u,\lambda)+\lambda-1$ far positions $j\le u+\lambda-1$ when $2\eta b^\lambda\le1$ is
\leanP{RadixWidthLaw}{139}{card_far_in_gap_ge}.
\end{proof}

Let $\mathcal A$ be the support of $a$, $S(N)=\cnt\mathcal A\cap[1,N]$, $\rho_N=S(N)/N$,
$\lambda_N=\log_b(1/\rho_N)$, $\ell_N=\log_b\log N$ and $W(N)=\sum_{n\in\mathcal A,\,n\le N}u(n)$.

\begin{theorem}[Sub-threshold width law]\label{thm:widthlaw}\evidence{Lean-supported paper proof}
Suppose that for constants $B\ge0$, $0<\kappa<1$ and $\omega\ge0$:
\begin{enumerate}[label=\textup{(W\arabic*)},start=0,nosep]
\item $\rho_N=(\log N)^{-\kappa+o(1)}$;
\item $A(X)\le X(\log X)^B$ for large $X$, and $|a(n)|\le n^{o(1)}$;
\item with $H_V=\lceil(B+2)\ell_N\rceil$, $R_N=\lceil(2+\omega)\lambda_N+(B+2)\ell_N\rceil$ and
$H_N=R_N+H_V$, all but $o(S(N))$ indices $n\in\mathcal A\cap[1,N]$ have no other element
of $\mathcal A$ within distance $H_N$;
\item $W(N)\le(\omega+o(1))S(N)\lambda_N$, and for each $\varepsilon>0$ only $o(S(N))$
indices $n\in\mathcal A\cap[1,N]$ have $u(n)<(\omega-\varepsilon)\lambda_N$.
\end{enumerate}
Then
\[
\bigl(1+\min(\omega,1)-o(1)\bigr)\rho_N\lambda_N\le1-D_N(F;b)\le(1+\omega+o(1))\rho_N\lambda_N.
\]
In particular $1-D_N(F;b)\sim(1+\omega)\rho_N\lambda_N$ whenever $\omega\le1$.
\end{theorem}

\begin{proof}
\emph{Upper bound.} Apply \cref{thm:displacedwidth} with $\varepsilon=\rho_N$ and
$L=\lceil\lambda_N\rceil+2$. Since $b^{V(n)}\le|a(n)|=n^{o(1)}$, $t(n)\ge n/2$ for large $n$.
Hence $P(N)\le S(N)+O(\log N)$, $W_U(N)\le W(N)+O(\log^2N)$ and $\Lambda_N=O(\log N)$.
Since $|U|b^{-u}<1$,
$\widehat P(N+L)\le S(N)+O(\log N)+b\sum_{i\ge1}|a(N+L+i)|b^{-i}=S(N)(1+o(1))$, and as
$b^{-L}\le\rho_N/b^2$ the Markov term counts $O(S(N))$ shifts. So
$1-D_N\le2\rho_N+(W(N)+S(N)\lambda_N+O(S(N)))/N$.

\emph{Lower bound.} Fix $\varepsilon\in(0,1)$, write $R=R_N$, $H=H_N$, and call
$n\in\mathcal A\cap(2H,N]$ \emph{good} if
(G1) $V(n)\le H_V$; (G2) no $n'\in\mathcal A$ has $0<|n'-n|\le H$; (G3) no $n'\in\mathcal A$
with $V(n')>H_V$ has $|t(n')-t(n)|\le R$; (G4) $|A(t(n))|\le\frac12$; (G5)
$u(n)\ge(\omega-\varepsilon)\lambda_N$; (G6) $u(n)\le(\omega+\varepsilon)\lambda_N$.

Since $b^V\le|a|$, Markov's inequality and (W1) give
$\cnt\{n\le3N:V(n)>H_V\}\le b^{-H_V}A(3N)\ll N(\log N)^{-2}$, which is $o(S(N)/\ell_N)$
because $\kappa<1$; this controls (G1). (G2) fails for $o(S(N))$ indices by (W2). Each
exceptional $n'$ in (G3) has $n'\le3N$ and spoils at most $2R+H_V+1=O(\ell_N)$ indices, a total
of $o(S(N))$. For $n$ satisfying (G1)--(G3), $t=t(n)$ is the only arrival in $[t-R,t+R]$
and $C(t)=U(n)$, so $A(t)=b^{-R}A(t+R)$ by \cref{lem:source}; $n\mapsto t(n)+R$ is
injective, and \cref{lem:meantail} with $L=0$ applied to $C$ gives
$\sum_{s<M}|A(s)|\le(A(2M)+O(M^{o(1)}))/(b-1)$, the second term bounding the sources
$n>2M$ (for which $t(n)>M$ and $|a(n)|b^{M-n}\le n^{o(1)}b^{-n/2}$). With $M=N+R+1$ and
$b^{-R}\le\rho_N(\log N)^{-B-2}$, Markov's inequality
(\leanP{RadixWidthLaw}{225}{card_large_tail_le}) shows that (G4) fails for
$O(\rho_NN/(\log N)^2)=o(S(N))$ indices. (G5) fails for $o(S(N))$ indices by (W3), and
together with the first-moment bound in (W3) this gives $u(n)/\lambda_N\to\omega$ in $L^1$
on $\mathcal A\cap[1,N]$, so (G6) also fails for $o(S(N))$ indices. Thus the good set $G$
has $|G|=(1-o(1))S(N)$.

Put $\Delta^+=(1+\omega+\varepsilon)\rho_N\lambda_N$, $\zeta=\rho_N$, $\eta=2\rho_N+2\Delta^+$
and $\lambda'=\lfloor\log_b(1/(2\eta))\rfloor=\lambda_N-\log_b\lambda_N-O(1)$. For $n\in G$
with $t=t(n)$ and $u=u(n)$, \cref{lem:onesource} applies, since $b\nmid U(n)=C(t)$. It marks
the positions $1\le j\le\lambda'$ and $u+1\le j\le u+\lambda'-1$, at which
$\dist{x_{t-j}}\ge\eta$; there are at least $\min(u,\lambda')+\lambda'-1$ of them, and by
(G6) all lie in the arrival-free interval $[t-R,t)$, because
$u+\lambda'-1\le(1+\omega+\varepsilon)\lambda_N\le R$. So the marked sets of distinct good
sources are disjoint, and by (G5) their proportion is at least
$\Delta:=(1+\min(\omega,1)-2\varepsilon-o(1))\rho_N\lambda_N$. The upper-bound count shows that at
least $N(1-\Delta^+)\ge\Delta N$ points satisfy $\dist{x_n}\le\zeta$. Put
$\eta'=\frac32(\rho_N+\Delta^+)$; then every marked point satisfies
$\dist{x_{t-j}}\ge\eta>\eta'$, and $\Delta\le\Delta^+$ gives $\zeta+\Delta<\eta'$.
\Cref{lem:circle} with $\eta'$ in place of $\eta$ gives $1-D_N\ge\Delta$; let
$\varepsilon\to0$.
\end{proof}

\begin{remark}[Signs, mixtures and the threshold]\label{rem:widthlaw}\evidence{Paper proof}
Nothing in the proof uses $a\ge0$: \cref{lem:onesource,lem:circle} and the Markov bounds
are sign-blind, and (G2) excludes cancellation in the fibre of a good source. If
$u(n)/\lambda_N$ converges in distribution to a law $\mu$ on $[0,1)$ and
$W(N)\le(\int x\,d\mu+o(1))S(N)\lambda_N$, replacing (G5)--(G6) by $u(n)\le\lambda_N$ gives
$1-D_N\sim(1+\int x\,d\mu)\rho_N\lambda_N$ (\emph{mixture form}). When $\omega>1$ the
upper value is attained if the internal positions $\lambda'<j\le u$ at which
$\dist{b^{-j}Y(t)}<\eta$ contribute $o(S(N)\lambda_N)$ in total; such a position occurs
exactly when the digits of $C(t)$ below position $j$ begin with a run of about
$\lambda_N$ zeros or of about $\lambda_N$ digits $b-1$.
\end{remark}

For bounded coefficients ($\omega=0$) the law reads $1-D_N\sim\rho_N\log_b(1/\rho_N)$. Two
classical examples follow.

\begin{corollary}[Squares and primes]\label{cor:primes}\evidence{Paper proof}
For every $b\ge2$, $1-D_N(\sum_{j\ge1}b^{-j^2};b)\sim\frac12(\log_bN)/\sqrt N$, and
\[
1-D_N\Bigl(\sum_pb^{-p};b\Bigr)\sim\frac{\log\log N}{\log b\cdot\log N}.
\]
\end{corollary}

\begin{proof}
For the squares, the sources are isolated at every radius $o(\sqrt N)$ and the power-law
version of the argument above (with $\rho_N\sim N^{-1/2}$, $\lambda_N\sim\frac12\log_bN$)
applies verbatim. For the primes take $\kappa=1$ in the form $\rho_N\sim1/\log N$,
$\lambda_N\sim\log_b\log N$; the Selberg upper-bound sieve
\cite[Ch.~3]{HalberstamRichert} gives
$\cnt\{p\le x:p+h\text{ prime}\}\ll\mathfrak S(h)x/(\log x)^2$ uniformly in $h$, with
$\sum_{h\le H}\mathfrak S(h)\ll H$, so for $H=O(\log\log x)$ all but
$O(x\log\log x/(\log x)^2)=o(\pi(x))$ primes have no other prime within distance $H$. The coefficients equal $1$, so (G1), (G4)--(G6) are
automatic, and the proof of \cref{thm:widthlaw} goes through with $\kappa=1$ because the
only use of $\kappa<1$ was in (G1).
\end{proof}

\section{Euler's product and its powers}\label{sec:euler}

Let $P_k^-=k(3k-1)/2$ and $P_k^+=k(3k+1)/2$ $(k\ge1)$ be the generalized pentagonal
numbers, $T_k=k(k+1)/2$ the triangular numbers, and $P(M)=\cnt\{p<M\}$ the number of
generalized pentagonal numbers $p\ge1$ below $M$. Since $k(3k\mp1)/2<M$ is equivalent to
$k<(\pm1+\sqrt{24M+1})/6$,
\begin{equation}\label{eq:pentcount}
\tfrac13\sqrt{24M+1}-2\le P(M)\le\tfrac13\sqrt{24M+1}.
\end{equation}

\subsection{Series and products}\label{subsec:seriesproduct}

\begin{theorem}[Euler, Jacobi]\label{thm:eulerjacobi}\evidence{External input}
For $|q|<1$,
\[
\prod_{k\ge1}(1-q^k)=1+\sum_{k\ge1}(-1)^k\bigl(q^{P_k^-}+q^{P_k^+}\bigr),\qquad
\prod_{k\ge1}(1-q^k)^3=\sum_{k\ge0}(-1)^k(2k+1)q^{T_k}.
\]
\end{theorem}

The first identity is Euler's pentagonal number theorem
\cite{EulerE244,EulerE541,Bell2005}, the second Jacobi's \cite{Jacobi1829}. In Lean the
\emph{series} on the right are defined directly:
\leanP{EulerPentagonal}{133}{eulerPentagonalSum} is
$\sum_m\mathtt{pentCoeff}(m)b^{-m}$, where $\mathtt{pentCoeff}(0)=1$,
$\mathtt{pentCoeff}(m)=(-1)^k$ for a chosen $k$ with $m\in\{P_k^-,P_k^+\}$, and $0$
otherwise; and \leanP{EulerPentagonal}{355}{eulerCubeSum} is Jacobi's triangular series.
Neither identity is formalized, and no Lean theorem below uses the signs of the
coefficients: each holds for every coefficient sequence with the same support and the same
absolute values. Every statement about the products $E_b$ and $E_b^3$ is obtained from the
Lean statement about the series through \cref{thm:eulerjacobi}. Nothing about $E_b^2$ is
formalized.

\begin{proposition}[Lean statements for the series]\label{prop:leaneuler}\evidence{Kernel-verified}
For every $b\ge2$: the pentagonal series has at most $2(\lfloor\sqrt{2M}\rfloor+1)$ nonzero
coefficients up to $M$ and mass at most $M$; it lies in $K_b$ and
$D_N(x+\cdot;b)-D_N(x;b)\to0$ for every real $x$; and for all $N\ge1$, $L\ge0$,
$\varepsilon>0$,
\[
1-D_N\le2\varepsilon+\frac{2L(\lfloor\sqrt{2(N+L)}\rfloor+1)}N+\frac{N+L+2}{N\varepsilon b^L}.
\]
Jacobi's triangular series has at most $\lfloor\sqrt{2M}\rfloor$ nonzero coefficients up
to $M$ and mass at most $4M$, and it lies in $K_b$ and preserves discrepancy under
translation. For integer sequences $a,c$ the Cauchy convolution satisfies
$\sum_{1\le m\le M}|(a*c)(m)|\le(|a(0)|+A(M))(|c(0)|+C(M))$ with $C$ the mass of $c$.
\end{proposition}

\begin{proof}
\leanP{EulerPentagonal}{206}{pentagonal_supportCount_le},
\leanP{EulerPentagonal}{184}{pentCoeff_absMass_le},
\leanP{EulerPentagonal}{301}{eulerPentagonal_discrepancy_tendsto_one},
\leanP{EulerPentagonal}{316}{eulerPentagonal_preserves_discrepancy},
\leanP{EulerPentagonal}{283}{eulerPentagonal_finite_gap_upper},
\leanP{EulerPentagonal}{372}{jacobiCube_supportCount_le},
\leanP{EulerPentagonal}{409}{jacobiCube_absMass_le},
\leanP{EulerPentagonal}{464}{eulerCube_discrepancy_tendsto_one},
\leanP{EulerPentagonal}{476}{eulerCube_preserves_discrepancy} and
\leanP{EulerPentagonal}{497}{cauchyConv_absMass_le}.
\end{proof}

With $L\approx\log_bN$ and $\varepsilon\approx N^{-1/2}$ the Lean bound has the correct order
$(\log N)/\sqrt N$ but leading constant $2\sqrt2$, a factor $2\sqrt3$ above the sharp
constant of \cref{thm:eulerdeficit}.

\subsection{Digits and words of $E_b$}\label{subsec:eulerdigits}

The two members of the pair $k$ share the sign $(-1)^k$, and
$P_k^+-P_k^-=k$, $P_{k+1}^--P_k^+=2k+1$.

\begin{theorem}[Digits]\label{thm:eulerdigits}\evidence{Paper proof}
For $j\ge1$ put $a_j=P_{2j}^-=j(6j-1)$. Then
\[
E_b=0.(b-2)(b-1)\,0\,0\,B_1B_2B_3\cdots,\qquad B_j=1\,0^{2j}(b-1)^{4j}(b-2)(b-1)^{2j+1}0^{4j+2},
\]
where $B_j$ occupies the positions $a_j,\dots,a_{j+1}-1$ and has length $12j+5$.
\end{theorem}

\begin{proof}
Group the positive pair $2j$ with the negative pair $2j+1$, and let $a=P_{2j}^-$,
$c=P_{2j}^+$, $d=P_{2j+1}^-$, $e=P_{2j+1}^+$, so that $c-a=2j$, $d-c=4j+1$ and $e-d=2j+1$. For
$1\le m<n$,
\[
1-b^{-m}-b^{-n}=\sum_{i=1}^{m-1}(b-1)b^{-i}+(b-2)b^{-m}+\sum_{i=m+1}^n(b-1)b^{-i}.
\]
Hence $b^{-a}+b^{-c}(1-b^{-(d-c)}-b^{-(e-c)})$ has digit $1$ at $a$, zeros on $(a,c]$,
$b-1$ on $(c,d)$, $b-2$ at $d$, $b-1$ on $(d,e]$ and zeros elsewhere. The head
$1-b^{-1}-b^{-2}$ is the case $c=0$, $m=1$, $n=2$. The digit supports are disjoint and all
digits lie in $[0,b-1]$, so there are no carries, and the expansion does not end in
$(b-1)^\infty$. Finally $a_{j+1}-a_j=12j+5$.
\end{proof}

\begin{corollary}[Missing words]\label{cor:missing}\evidence{Paper proof}
Inside the blocks $B_j$ $(j\ge1)$ the maximal runs of zeros have lengths $2j$ and $4j+2$
(plus the isolated zero at $P_{2j+1}^-$ when $b=2$), and the maximal runs of $b-1$ have
lengths $4j$ and $2j+1$; the four-digit head $(b-2)(b-1)00$ contributes one further run of
$b-1$ of length $1$ and one run $00$. In base
$2$ the word $10^m1$ never occurs for odd $m\ge3$; for $b\ge3$ the word $11$ never occurs,
and for $b\ge5$ the digits $2,\dots,b-3$ never occur. In particular $E_b$ is not
disjunctive.
\end{corollary}

\begin{proof}
Read off \cref{thm:eulerdigits}.
\end{proof}

Let $Z_w(N)$ be the number of $0\le n<N$ at which the base-$b$ word of length $L$ read
from position $n+1$ equals $w$.

\begin{theorem}[Word frequencies]\label{thm:eulerwords}\evidence{Paper proof}
For $L,N\ge1$ and $w\in\{0^L,(b-1)^L\}$,
\[
\Bigl|Z_w(N)-\frac N2\Bigr|\le(2L+3)\sqrt{N/6}+3L^2+5L+7,
\]
and every other word of length $L$ occurs at most $(4L-2)\sqrt{N/6}+2L^2+12L+4$ times.
\end{theorem}

\begin{proof}
For $0^L$, a valid start reads a window inside one zero run, and runs never merge across
blocks because every block begins with $1$. Block $j$ contributes
$g_j=(2j-L+1)^++(4j+3-L)^+$ valid starts (one more when $b=2$, $L=1$) among its $12j+5$
starts. Put $\Delta(N)=Z(N)-N/2$. For $j\ge(L-1)/2$ the increment of $\Delta$ over block $j$
is $\frac32-2L$ (or $\frac52-2L$), at most $2L$ in absolute value; each smaller block changes
$\Delta$ by at most half its length, in total at most $\frac34L^2+\frac{11}4L$, and the first
four digits contribute at most $2$. Inside block $J$ the function $\Delta$ drops by
$\frac12$, rises over the first zero run to at most $J$, falls to $-2J-L-\frac12$, and
rises again, so its excursion is at most $2J+L+\frac12$. If $N$ lies in block $J$ then
$6J^2-J-1\le N$, so $J\le\sqrt{N/6}+\frac12$, and
$|\Delta(N)|\le(2L+2)\sqrt{N/6}+0.75L^2+6.75L+5.5$. The word $(b-1)^L$ is handled in the
same way, with block gain $(4j-L+1)^++(2j+2-L)^+$. Words other than $0^L$ and $(b-1)^L$
start only at the
$4L-2$ remaining starts of each block with $j\ge L/2$ and at the starts of smaller blocks.
\end{proof}

In base $2$ the word $(b-1)^1$ is the digit $1$, so $E_2$ is simply normal in base $2$;
for $b\ge3$ the digits $1,\dots,b-2$ have frequency $O(N^{-1/2})$, so $E_b$ is not simply
normal in base $b$.

\subsection{The sharp deficit of $E_b$}\label{subsec:eulerdeficit}

\begin{theorem}[Sharp deficit]\label{thm:eulerdeficit}\evidence{Lean-supported paper proof}
For $b\ge2$ and $N\ge2$,
\[
1-D_N(E_b;b)=\sqrt{\tfrac23}\,\frac{\log_bN}{\sqrt N}+O_b\Bigl(\frac1{\sqrt N}\Bigr).
\]
\end{theorem}

\begin{proof}
\emph{Upper bound.} $E_b-1$ is a sparse series with unit coefficients at the generalized
pentagonal numbers. The kernel-verified counting lemma
\leanP{Phase18Synthetic}{768}{sparse_farCount_le} (with $B=1$ and $\varepsilon=2b^{-w}$) and
\cref{lem:deficit} give, for every $w\ge1$,
\[
1-D_N(E_b)\le4b^{-w}+\frac{w\,P(N+w)}N.
\]
This instantiation to the pentagonal positions is a written step. With
$w=\lceil\log_b\sqrt N\rceil+1$ and \eqref{eq:pentcount}, the right side is
$\sqrt{2/3}\log_bN/\sqrt N+O(N^{-1/2})$.

\emph{Lower bound.} Fix $n$, let $p>n$ be the next generalized pentagonal number,
$r=p-n$, $s=\pm1$ its sign and $g$ the gap after $p$. All later exponents are distinct and at
least $r+g$, so $b^nE_b\equiv sb^{-r}+\theta\frac b{b-1}b^{-r-g}\pmod1$ with $|\theta|\le1$.
If $r,g\ge2$ then $\dist{x_n}\ge b^{-r}(1-2b^{-g})\ge b^{-r}/2$. Fix $t\in(0,\frac18]$ and let
$R$ be the largest integer with $b^{-R}/2>t$. All but at most $2R$ of the $P(N+1)$
pentagonal numbers $p\le N$ have both neighbouring gaps at least $R$, and for them the shifts
$p-r$, $2\le r\le R$, are distinct elements of $[0,N)$. Hence
$E_N(t)\ge(P(N+1)-2R)(\log_b\frac1{2t}-2)$. Apply \cref{lem:separation} with
$\varepsilon=1/N$. By the upper-bound count $E_N(1/N)=O(\sqrt N\log N)$, so the first branch
is $1-o(1)$. In the second branch put $M=P(N+1)-2\log_bN=2\sqrt{2N/3}+O(\log N)$; then
$\ell+E_N(\ell+1/N)/N\ge h(\ell)=\ell+\frac MN(\log_b\frac1{2(\ell+1/N)}-2)$, a convex
function minimized at $\ell+1/N=M/(N\log b)$ with value at least
$\frac MN(\log_b\frac{N\log b}{2M}-2)-\frac1N=\sqrt{2/3}\log_bN/\sqrt N-O_b(N^{-1/2})$. Arcs with
$\ell+1/N>\frac18$ cost at least $\frac18-\frac1N$.
\end{proof}

\Cref{thm:B} follows from
\cref{thm:eulerdigits,cor:missing,thm:eulerwords,thm:eulerdeficit,thm:sharp}.

\paragraph{Numerical illustration.}
Both signs matter: positive pairs place samples just above $0$ and negative pairs just
below $1$. The second-order term $(1-D_N)N/P(N)-\frac12\log_bN$ oscillates because only
integer distances $r$ occur. At $N=2^{20}$ the exact deficit is $0.016247$ in base $2$, a
ratio $1.0188$ to the main term; this is a finite measurement.

\subsection{The cube}\label{subsec:cube}

\begin{theorem}[Cube]\label{thm:cube}\evidence{Paper proof}
For every $b\ge2$:
\begin{enumerate}[label=\textup{(\alph*)},nosep]
\item for $j\ge2$ the positions $(T_{2j-1},T_{2j+1}]$ of $E_b^3$ read the $2j$-digit
zero-padded expansion of $4j$ followed by the $(2j+1)$-digit expansion of
$b^{2j+1}-4j-3$;
\item for $w\in\{0^L,(b-1)^L\}$, $Z_w(N)=N/2+O((L+\log N)\sqrt N)$, and every other word has
frequency $0$;
\item $E_b^3$ is disjunctive;
\item $1-D_N(E_b^3;b)=\sqrt2\,\log_bN/\sqrt N+O_b(N^{-1/2})$.
\end{enumerate}
\end{theorem}

\begin{proof}
(a) By \cref{thm:eulerjacobi},
$(4j+1)b^{-T_{2j}}-(4j+3)b^{-T_{2j+1}}=4jb^{-T_{2j}}+(b^{2j+1}-4j-3)b^{-T_{2j+1}}$, and the
left side is less than $b^{-T_{2j-1}}$ because $4j+1<b^{2j}$ for $j\ge2$; the terms with
$j\le1$ sum to a nonnegative multiple of $b^{-6}$.

(b) Per block the zero run has length $2j-O(\log j)$, the $(b-1)$-run has length
$2j+1-O(\log j)$, and $O(\log j)$ digits are spent on the expansions of $4j$ and of
$4j+3$; the proof of \cref{thm:eulerwords} applies with an extra $O(\log j)$ per block.

(c) Given a word $w$, let $W$ be the integer with digits $1w$. For $b^t\ge4$ the interval
$[Wb^t,(W+1)b^t)$ contains a multiple $4j$ of $4$, whose expansion begins with $1w$, and (a)
places it verbatim in the expansion; let $t\to\infty$.

(d) \emph{Upper bound.} For $n\in[T_{k-1},T_k)$ let $r=T_k-n\in[1,k]$; then
$|b^nE_b^3-(\text{integer})|\le2(2k+1)b^{-r}$, so
$E_N(\varepsilon)\le\sum_{k:\,T_k-\min(R_k,k)<N}\min(R_k,k)$ with
$R_k=\lceil\log_b(2(2k+1)/\varepsilon)\rceil$. With $\varepsilon=N^{-1/2}$ there are
$\sqrt{2N}+O(\log N)$ such $k$, each with $R_k=\log_bN+O(1)$, and \cref{lem:deficit} gives
$1-D_N\le\sqrt2\log_bN/\sqrt N+O(N^{-1/2})$.

\emph{Lower bound.} Let $w_k$ be the number of digits of $2k+1$. For $w_k+2\le r\le k$,
$\dist{x_{T_k-r}}\ge(2k+1)b^{-r}/2$. For $r\le w_k+1$ the sample lies within
$\eta_k=4(2k+1)b^{-k-2}$ of $\pm(2k+1)b^{-r}\bmod1$, which is far from $0$ unless
$(2k+1)\bmod b^r$ falls in a set of at most $2\tau b^r+2$ residues, $\tau=t+\eta_k$. As $k$
runs over $K$ consecutive integers, each class modulo $b^r$ contains at most $2K/b^r+1$ of the
odd numbers $2k+1$, so the exceptional pairs $(k,r)$ number $O(K(1+t\log K))$. With
$K'=\max\{k:T_k\le N\}=\sqrt{2N}+O(1)$ and
$\sum_{k\le K'}\log_b(2k+1)=K'\log_b(2K')-K'/\log b+O(\log K')$,
\[
E_N(t)\ge K'\Bigl(\log_b(2K')+\log_b\frac1t-C\Bigr)-O\bigl(K'(t\log K'+1)\bigr).
\]
\Cref{lem:separation} with $\varepsilon=1/N$ gives
$1-D_N\ge\min_\ell\bigl(\ell/2+\frac{K'}N(\log_b(2K')+\log_b\frac1{\ell+1/N}-C)\bigr)-\frac1N$,
the term $t\log K'$ being absorbed by $\ell/2$; the minimum, at $\ell\asymp K'/N$, is
$\frac{K'}N(\log_bN-O(1))$.
\end{proof}

The constant $\sqrt2$ is the product of the insertion density $\sqrt{2N}/N$ and the
$\log_bN$ shifts that each insertion spoils: $\frac12\log_bN$ for precision and
$\frac12\log_bN$ for the digits of $2k+1$, which are spread around the circle.

\subsection{The square}\label{subsec:square}

Since $24P_k^\mp+1=(6k\mp1)^2$ and $(-1)^k=\chi_{12}(|6k\mp1|)$, where $\chi_{12}(m)=1$ for
$m\equiv\pm1$ and $-1$ for $m\equiv\pm5\pmod{12}$,
\[
E(q)^2=\sum_{n\ge0}a_2(n)q^n,\qquad \eta(12\tau)^2=\sum_{n\ge0}a_2(n)q^{12n+1}=\sum_{m\ge1}c(m)q^m,
\]
where
\[
c(m)=\sum_{\substack{m_1,m_2\ge1\\ m_1^2+m_2^2=2m}}\chi_{12}(m_1)\chi_{12}(m_2).
\]

\begin{proposition}[Square, elementary bounds]\label{prop:squareelem}\evidence{Paper proof}
If $a_2(n)\ne0$, then $12n+1$ is a sum of two squares. The mass satisfies
$\sum_{n\le X}|a_2(n)|\le(P(X+1)+1)^2\le(2\sqrt{2X/3}+2)^2$, the support count is
$\ll X/\sqrt{\log X}$, and $E_b^2\in K_b$ with $1-D_N(E_b^2;b)\ll_b\LL/\sqrt{\log N}$.
\end{proposition}

\begin{proof}
$24n+2=m_1^2+m_2^2$ gives $12n+1=u^2+v^2$ with $u=(m_1+m_2)/2$, $v=(m_1-m_2)/2$. The mass bound
counts pairs of pentagonal exponents. The integers that are sums of two squares have
counting function $\ll x/\sqrt{\log x}$ by Landau's theorem (or by the Halberstam--Richert
bound for multiplicative indicators), and $n\mapsto12n+1$ is affine. Apply
\cref{thm:largermass} with $\kappa=\frac12$, $\beta=0$.
\end{proof}

\begin{theorem}[CM identification {\cite{Serre1985,Martin1996}}]\label{thm:etainput}\evidence{External input}
$\eta(12\tau)^2$ is a normalized Hecke eigenform of weight one with complex multiplication
by both $\Q(\sqrt{-1})$ and $\Q(\sqrt{-3})$; its projective Galois image is the Klein four
group.
\end{theorem}

\begin{theorem}[Square]\label{thm:square}\evidence{Paper proof}
Granting \cref{thm:etainput}:
\begin{enumerate}[label=\textup{(\alph*)},nosep]
\item $c$ is multiplicative, $c(2^e)=c(3^e)=0$ for $e\ge1$, $c(p^e)=(e+1)\epsilon_p^e$ with
$\epsilon_p=\pm1$ for $p\equiv1\pmod{12}$, and $c(p^{2e+1})=0$, $c(p^{2e})=(-\chi_{-4}(p))^e$
for the other primes $p\ge5$. Hence $a_2(n)\ne0$ if and only if every prime
$p\not\equiv1\pmod{12}$ divides $12n+1$ to an even power, and then
$|a_2(n)|=\prod_{p\equiv1\,(12)}(v_p(12n+1)+1)$.
\item The number $S_\eta(N)$ of $1\le n\le N$ with $a_2(n)\ne0$ satisfies
$S_\eta(N)\sim C_\eta N/(\log N)^{3/4}$ with
\begin{equation}\label{eq:Ceta}
C_\eta=\frac{12}{\Gamma(1/4)}\Bigl(\frac{\pi^2\log(2+\sqrt3)}{54}
\prod_{\substack{p\ge5\\p\not\equiv1\,(12)}}\Bigl(1-\frac1{p^2}\Bigr)^{-2}\Bigr)^{1/4}=2.418529\ldots,
\end{equation}
and $\sum_{n\le X}|a_2(n)|\asymp X/(\log X)^{1/2}$.
\item For $p\equiv1\pmod{12}$ write $p=u^2+v^2$ with $u$ odd; then $c(p)=+2$ if $3\mid v$
(that is, $p=x^2+36y^2$) and $c(p)=-2$ if $3\mid u$ ($p=4x^2+9y^2$). No congruence
condition modulo a fixed integer determines $\epsilon_p$ for all large $p$.
\end{enumerate}
\end{theorem}

\begin{proof}
(a) A weight-one eigenform has $c(p)=\tr\rho(\mathrm{Frob}_p)$ for an odd representation
$\rho$ with projective image $V_4$ \cite{DeligneSerre1974}. The three nontrivial elements of
$V_4$ lift to trace-zero matrices, and Frobenius is central exactly when $p$ splits
completely in $\Q(i,\sqrt3)$, that is, $p\equiv1\pmod{12}$, when $c(p)=\pm2$. The values at
prime powers follow from $c(p^{e+1})=c(p)c(p^e)-\chi_{-4}(p)c(p^{e-1})$.

(b) The support indicator is multiplicative, equal to $1$ at the primes $\equiv1\pmod{12}$
and at all $p^2$ with $p\ge5$. Let $F(s)=\sum_{m\in\mathcal S}m^{-s}$ for the set
$\mathcal S$ of such $m$. At a prime $p\ge5$, $p\not\equiv1\pmod{12}$, two of the four
characters modulo $12$ take the value $1$ and two the value $-1$, so
\[
F(s)^4=L(s,\chi_0^{(12)})L(s,\chi_{-4}^{(12)})L(s,\chi_{-3}^{(12)})L(s,\chi_{12})
\prod_{\substack{p\ge5\\p\not\equiv1\,(12)}}(1-p^{-2s})^{-2}.
\]
With $L(1,\chi_{-4})=\pi/4$, $L(1,\chi_{-3})=\pi/(3\sqrt3)$, $L(1,\chi_{12})=\log(2+\sqrt3)/\sqrt3$
and the Euler factors of the imprimitive characters at $2$ and $3$,
$\lim_{s\to1^+}(s-1)F(s)^4=\frac{\pi^2\log(2+\sqrt3)}{54}\prod(1-p^{-2})^{-2}$. The
Selberg--Delange theorem \cite[II.5]{Tenenbaum} with $z=\frac14$ gives
$\cnt\mathcal S\cap[1,x]\sim\lambda_0x(\log x)^{-3/4}$ with
$\lambda_0=\lim(s-1)^{1/4}F(s)/\Gamma(1/4)$, and $S_\eta(N)=\cnt\mathcal S\cap[1,12N+1]-1$
gives $C_\eta=12\lambda_0$. For $|c|$, $\sum_{p\le x}|c(p)|/p=\frac12\log\log x+O(1)$ and
$|c(p^e)|\le e+1$, so Wirsing's theorem \cite{Wirsing1961} and Shiu's bound \cite{Shiu1980}
give the order of the mass.

(c) Exactly one of $u,v$ is divisible by $3$. The representations $2p=m_1^2+m_2^2$ with
$m_i\ge1$ are $(u+v,|u-v|)$ and its swap, so $c(p)=2\chi_{12}(u^2-v^2)$, and
$\chi_{12}=\chi_{-4}\chi_{-3}$ with $u^2-v^2\equiv1\pmod4$. Modulo $3$, $u^2-v^2\equiv1$ if
$3\mid v$ and $\equiv-1$ if $3\mid u$. The reduced forms of discriminant $-144$ show
$\Cl(-144)\cong\Z/4$, and $p=x^2+36y^2$ if and only if $p$ splits completely in the ring
class field $H$ of $\Z[6i]$, whose Galois group over $\Q$ is dihedral of order $8$
\cite[Thm.~9.4, Lemma~9.3]{Cox}. If $\epsilon_p$ depended only on $p\bmod Q$, then, up to
finitely many primes, every prime splitting completely in $\Q(\zeta_{Q'})$,
$Q'=\mathrm{lcm}(12,Q)$, would split completely in $H$, and Bauer's theorem
\cite[VII.13.9]{Neukirch} would give $H\subseteq\Q(\zeta_{Q'})$, which is impossible because
$\mathrm{Gal}(H/\Q)$ is nonabelian.
\end{proof}

At $N=10^6$ there are $307{,}819$ nonzero coefficients $a_2(n)$, $153{,}878$ of them negative;
the ratio of $S_\eta(N)$ to the two-term Selberg--Delange approximation is $1.00576$. These are
finite measurements.

\subsection{Polynomials, rational scalings and unary quadratic series}\label{subsec:poly}

\begin{theorem}[Polynomials and scalings]\label{thm:polyscal}\evidence{Lean-supported paper proof}
\leavevmode
\begin{enumerate}[label=\textup{(\alph*)},nosep]
\item If $P=b^{-j}\sum_{i\le3}q_ix^i$ with $j\ge0$ and $q_i\in\Z$, then $P(E_b)\in K_b$ and
$1-D_N(P(E_b))\le\sum_{i=1}^3|q_i|(1-D_N(E_b^i))+j/N$.
\item For $r\in\{1,3\}$ and a reduced fraction $a/d$, $(a/d)E_b^r\in K_b$ if and only if
$a/d\in\Z[1/b]$.
\item If $b$ is odd and $a$ is odd, then $D_N(\frac a2E_b^2;b)\to\frac23$; in particular
$\frac a2E_b^2\notin K_b$.
\end{enumerate}
\end{theorem}

\begin{proof}
(a) is \cref{lem:linear} with \cref{thm:eulerdeficit,thm:cube,prop:squareelem}.

For (b) and (c) we use the residue-density test of \cite{DiscrepancyCompanion}: if $z\in K_b$
concentrates at the fixed point $0$ and $k_n$ is an integer nearest to $b^nz$, then
$(a/d)z\in K_b$ if and only if $z/d\in K_b$, if and only if some residue modulo $d$ carries
density one among the $k_n$. Since division by powers of $b$ preserves $K_b$
(\cref{lem:linear}) and $K_b$ is a group, it suffices to show $z/\ell\notin K_b$ for every
prime $\ell\nmid b$. Let $I_n=\sum_{m\le n}a(m)b^{n-m}$ be the integral head of
$z=F_b(a)$. By the transfer bounds, $k_n=I_n$ outside a set of density zero. Between
consecutive support points $I_{n+1}=bI_n$, and at a support point $p$,
$I_p=bI_{p-1}+a(p)$. Multiplication by $b$ permutes $\Z/\ell$, so on a run between support
points the residue of $I_n$ either stays at a fixed residue $r_0$ with $br_0\equiv r_0$ or
equals $r_0$ at most at every other position. Two adjacent runs both fixed at $r_0$ force
$\ell\mid a(p)$ at the support point between them.

For $E_b$, $a(p)=\pm1$, so no two adjacent runs are fixed at $r_0$; the runs alternate in
length between $k$ and $2k+1$, so any family of pairwise nonadjacent runs has density at most
$\frac23$, and no residue has density one. For $E_b^3$ the runs $[T_k,T_{k+1})$ have length
$k+1$ and the support point $T_{k+1}$ carries $\pm(2k+3)$. Among $k\in[K,2K]$, at most
$K/\ell+O(1)$ have $\ell\mid2k+3$; for each other $k$ one of the two runs meeting at
$T_{k+1}$ is not fixed at $r_0$, and each run serves at most two values of $k$. These runs
have lengths in $[K+1,2K+2]$ and lie in $[T_K,T_{2K+2})$, and each of them has at least
half of its positions away from $r_0$, so a proportion $\ge(1-1/\ell)/16-o(1)>0$ of
$[T_K,T_{2K+2})$ is not at $r_0$. This proves (b).

(c) Coefficientwise $E(q)^2\equiv E(q^2)\pmod2$, so $a_2(n)$ is odd exactly when $n=2p$ with $p$
a generalized pentagonal number or $p=0$. For odd $b$, $I_n\equiv I_{n-1}+a_2(n)\pmod2$, so
$I_n\bmod 2$ is the parity of $\cnt\{p\ge1:2p\le n\}$. The runs between consecutive points
$2p$ have lengths $2k$ and $2(2k+1)$ and carry alternating parities, with weights $\frac13$
and $\frac23$. The orbit of $\frac a2E_b^2$ therefore concentrates at the two points $0$ and
$\frac12$ with masses $\frac13$ and $\frac23$ (in some order), and $D_N\to\frac23$.
\end{proof}

The rational-scaling criterion for $E_b^2$ and odd primes $\ell$ is
\cref{thm:squarescaling}. Explicit limits: $\{2^nE_2/3\}$ has limit law
$\frac16\delta_0+\frac5{12}(\delta_{1/3}+\delta_{2/3})$, so $D_N(\frac13E_2;2)\to\frac12$,
and $D_N(\frac12E_3;3)\to\frac23$; the complete word laws are in \cref{thm:automaton}.

Call $F=\sum_{t\in T}\epsilon(t)b^{-f(t)}$ a \emph{unary quadratic series} if $T\subseteq\Z$,
$|\epsilon(t)|\le B$ and $f(t)=(\alpha t^2+\beta t+\gamma)/\delta$ is a nonnegative integer
for $t\in T$, with integers $\alpha>0$, $\delta\ge1$. Examples are $E_b-1$,
$\Theta_3(1/b)$, $\sum_{t\ge0}b^{-t(t+1)}$ and Rogers' false theta series
$\sum_{t\ge0}(-1)^tb^{-t(t+1)/2}$.

\begin{theorem}[Quadratic polynomials in unary quadratic series]\label{thm:unary}\evidence{Lean-supported paper proof}
For unary quadratic series $F_1,\dots,F_s$, every polynomial of total degree at most $2$
in them with coefficients in $\Z[1/b]$ lies in $K_b$, with
$1-D_N\ll\LL/\sqrt{\log N}$.
\end{theorem}

\begin{proof}
Linear terms have $S(X)=O(\sqrt X)$ and bounded coefficients. For a product $F_iF_j$ the
coefficient of $b^{-m}$ is $\sum_{f_i(t)+f_j(u)=m}\epsilon_i(t)\epsilon_j(u)$, and
$\delta_i\delta_jm=g(t,u)$ is an integral binary quadratic polynomial with positive definite
quadratic part. By the shifted-lattice reduction of \cref{prop:extensions}(a) its support
has counting function $\ll X/\sqrt{\log X}$, and the mass is $O(X)$ by the area of an ellipse
(\cref{prop:leaneuler} gives the convolution mass bound in Lean). Apply
\cref{thm:largermass} and \cref{lem:linear}.
\end{proof}

Degree $3$ is out of reach: $\Theta_3(1/b)^3=\sum r_3(n)b^{-n}$ has support
$\{n\ne4^a(8k+7)\}$ of density $\frac56$ by Legendre's theorem. $E_b^3$ is the exception
because Jacobi's identity collapses the triple convolution onto the triangular numbers.

\section{Exact constants in every radix}\label{sec:radix}

In this section $LL=\log\log N$, $\ell_N=\log_b\log N$, and the external inputs are the
following classical theorems.

\begin{theorem}[Mean values of multiplicative functions]\label{thm:meanvalues}\evidence{External input}
\leavevmode
\begin{enumerate}[label=\textup{(\alph*)},nosep]
\item \textup{(Shiu, Halberstam--Richert)} For multiplicative $g\ge0$ with $g(p^e)\le A^e$ and
$g(n)\ll_\epsilon n^\epsilon$, $\sum_{n\le x}g(n)\ll\frac x{\log x}\exp\sum_{p\le x}g(p)/p$
\cite[Thm.~1]{Shiu1980}, \cite{HalberstamRichert}.
\item \textup{(Wirsing, Selberg--Delange)} If moreover $\sum_{p\le x}g(p)\log p\sim\tau x$ with
$\tau>0$ and $g(p^\nu)\le AB^\nu$, $B<2$, then
$\sum_{n\le x}g(n)\sim\frac{e^{-\gamma\tau}}{\Gamma(\tau)}\frac x{\log x}\prod_{p\le x}\sum_\nu g(p^\nu)p^{-\nu}$
\cite{Wirsing1961}; for sets defined by congruence conditions or by Chebotarev conditions in a
fixed field, the Selberg--Delange method gives the same with power-saving refinements and
finitely many local conditions \cite[II.5]{Tenenbaum}.
\item \textup{(Brun)} The upper-bound sieve \cite[Thm.~2.2]{HalberstamRichert}: let a set of
$X$ consecutive integers be sifted, for each prime $p$ of a set $\mathcal P$ with
$z_0<p\le z$, by $\omega(p)$ residue classes, where $0\le\omega(p)\le\min(A,p-1)$,
$\sum_{w\le p<w',\,p\in\mathcal P}\omega(p)\log p/p\le\kappa_0\log(w'/w)+A$ for
$2\le w\le w'$, and the remainders satisfy $|R_d|\le\omega(d)$. If $z_0<z\le X$, the
number of survivors is $\ll_{\kappa_0,A}X\prod_{z_0<p\le z,\,p\in\mathcal P}(1-\omega(p)/p)$.
\item \textup{(Landau, Bernays, Mertens)} $\cnt\{n\le x:r_Q(n)>0\}\sim C_Qx/\sqrt{\log x}$ for
every positive definite primitive integral binary $Q$
\cite{Landau1908,Bernays1912,BMO2011}; and Mertens' theorem holds in arithmetic progressions
and in Chebotarev classes of a fixed Galois extension.
\end{enumerate}
\end{theorem}

\subsection{Multiplicative Lambert series of real characters}\label{subsec:lambertradix}

Let $\chi$ be a real nonprincipal character modulo $q$, $c=1*\chi$ and
$S(N)=\cnt\{n\le N:c(n)>0\}\sim C_\chi N/\sqrt{\log N}$, where
\[
C_\chi=\Bigl(\frac{L(1,\chi)}{\pi}\prod_{\chi(p)=-1}\Bigl(1-\frac1{p^2}\Bigr)^{-1}\prod_{p\mid q}\Bigl(1-\frac1p\Bigr)^{-1}\Bigr)^{1/2},
\]
as follows from $F(s)^2=\zeta(s)L(s,\chi)\prod_{\chi(p)=-1}(1-p^{-2s})^{-1}\prod_{p\mid q}(1-p^{-s})^{-1}$
for $F(s)=\sum_{c(n)>0}n^{-s}$ and \cref{thm:meanvalues}(b). For $\chi=\chi_{-4}$,
$C_\chi=K=2^{-1/2}\prod_{p\equiv3\,(4)}(1-p^{-2})^{-1/2}=0.76422365\ldots$ (Landau--Ramanujan).
Locally $c(p^e)=1$ for $p\mid q$, $e+1$ if $\chi(p)=1$ (split), and $(1+(-1)^e)/2$ if
$\chi(p)=-1$ (inert), so $c\ge0$.

\begin{lemma}[First moment]\label{lem:firstmoment}\evidence{Lean-supported paper proof}
$\sum_{n\le N,\,c(n)>0}\log c(n)=\bigl(\frac{\log2}2+o(1)\bigr)S(N)\,LL$.
\end{lemma}

\begin{proof}
For a split prime $p$, $n=p^em$ with $p\nmid m$ is supported if and only if $m$ is, so
\[
\sum_{\substack{n\le N\\c(n)>0}}\log c(n)=\sum_{\chi(p)=1}\sum_{e\ge1}\log(e+1)\Bigl(S\bigl(\tfrac N{p^e}\bigr)-S\bigl(\tfrac N{p^{e+1}}\bigr)\Bigr);
\]
this exact identity for every multiplicative function is
\leanP{RadixWidthLaw}{327}{sum_log_eq_sum_primePow}. Fix $\theta\in(0,\frac12)$. The terms
$e=1$, $p\le N^{1-\theta}$ contribute
$\log2\,(C_\chi+o(1))\frac N{\sqrt{\log N}}\sum_{\chi(p)=1,\,p\le N^{1-\theta}}\frac1{p\sqrt{1-\log p/\log N}}$,
and partial summation with Mertens' theorem in progressions evaluates the sum as
$\frac12\int_{\log3/\log N}^{1-\theta}\frac{du}{u\sqrt{1-u}}+O_\theta(1)=\frac12LL+O(1)$. The
primes $N^{1-\theta}<p\le N$ contribute $\ll\sqrt\theta\,N/\sqrt{\log N}$, and the prime powers
$e\ge2$ contribute $O(S(N))$. This gives the upper bound; the lower bound follows from
\cref{lem:concentration}.
\end{proof}

\begin{lemma}[Concentration]\label{lem:concentration}\evidence{Paper proof}
For each $\varepsilon>0$ there is $c(\varepsilon)>0$ with
$\cnt\{n\le N:c(n)>0,\ |\log c(n)-\frac{\log2}2LL|>\varepsilon LL\}\ll S(N)(\log N)^{-c(\varepsilon)}$.
\end{lemma}

\begin{proof}
For $|t|\le1$, the multiplicative function $\ind_{c>0}c^t$ has split-prime value $2^t$ on
primes of density $\frac12$, so \cref{thm:meanvalues} gives
$\mathbb E[c^t\mid c>0]\asymp(\log N)^{(2^t-1)/2}$. Chebyshev's inequality with small $t>0$
bounds the upper tail, since $\frac12(2^t-1)-t(\frac{\log2}2+\varepsilon)<0$ for small $t$. For
the lower tail use $c(n)\ge2^{\omega_{\rm split}(n)}$ and $z^{\omega_{\rm split}}$ with
$0<z<1$: the count is $\ll S(N)(\log N)^{h(z)}$ with
$h(z)=\frac{z-1}2-(\frac12-\varepsilon')\log z$, and $h(1)=0$, $h'(1)=\varepsilon'>0$.
\end{proof}

\begin{lemma}[Radix valuations]\label{lem:valuations}\evidence{Lean-supported paper proof}
Let $V_b(n)=v_b(c(n))$, $U_b(n)=c(n)/b^{V_b(n)}$.
\begin{enumerate}[label=\textup{(\alph*)},nosep]
\item If $b=2^k$, then $U_b(n)\le2^{k-1}\,\mathrm{odd}(c(n))$ and
$\sum_{n\le x}\mathrm{odd}(c(n))\ll x/\sqrt{\log x}$; hence the widths $u_b(n)$ have bounded
mean on the support.
\item If $b$ is not a power of $2$ and $q_0$ is an odd prime dividing $b$, then
$V_b(n)\le v_{q_0}(c(n))=\sum_{p^e\|n,\ \chi(p)=1}v_{q_0}(e+1)$, every nonzero term has
$e\ge q_0-1\ge2$, and $\sum_{n\le x,\,c(n)>0}v_{q_0}(c(n))\ll x/\sqrt{\log x}$.
\item In every radix, $\cnt\{n\le x:V_b(n)>H\}\le b^{-H}\sum_{n\le x}c(n)\ll xb^{-H}$.
\end{enumerate}
\end{lemma}

\begin{proof}
(a) uses \cref{thm:meanvalues}(a) for $\mathrm{odd}(c)$, whose split-prime value is $1$.
(b) $q_0\mid e+1$ forces $e\ge q_0-1$ (the finite step is the declaration
\decl{padicVal_succ_pos_imp} of \texttt{RadixWidthLaw}); sum the prime-power decomposition of
\cref{lem:firstmoment} over $e\ge2$. (c) is Markov's inequality with $b^{V_b}\le c$ and the
hyperbola bound $\sum_{n\le x}c(n)=L(1,\chi)x+O(\sqrt x)$.
\end{proof}

\begin{lemma}[Pairs]\label{lem:pairs}\evidence{Paper proof}
Uniformly for $1\le|h|\le\sqrt x$, $\cnt\{n\le x:c(n)c(n+h)>0\}\ll x\log\log x/\log x$.
\end{lemma}

\begin{proof}
If $n,n+h$ are supported and an inert prime $p\nmid h$ divides $n(n+h)$, then $p^2$ divides
$n$ or $n+h$. Discard the $n$ with $p^2\mid n(n+h)$ for some inert $p>z_0=(\log x)^2$
($\le2x/z_0+2\sqrt x$ of them). The rest avoid two classes modulo every inert
$p\in(z_0,\sqrt x]$, $p\nmid h$; Brun's sieve in dimension $1$ and Mertens give
$\ll x\log z_0/\log x$. Dropping the primes $p\mid h$ costs a factor
$(1-2/z_0)^{-\log_2x}\to1$.
\end{proof}

\begin{theorem}[Lambert series in every radix]\label{thm:lambertradix}\evidence{Lean-supported paper proof}
For every real nonprincipal $\chi$ and every $b\ge2$,
\[
1-D_N(L_\chi(b);b)\sim
\begin{cases}\dfrac{C_\chi}{2\log b}\dfrac{LL}{\sqrt{\log N}},&b=2^k,\\[2ex]
\dfrac{C_\chi(1+\log2)}{2\log b}\dfrac{LL}{\sqrt{\log N}},&\text{otherwise.}\end{cases}
\]
The same holds for $\Theta_3(1/b)^2=1+4L_{\chi_{-4}}(b)$ with $\chi=\chi_{-4}$.
\end{theorem}

\begin{proof}
Apply \cref{thm:widthlaw} with $B=0$, $\kappa=\frac12$; then
$\lambda_N=\frac12\ell_N-\log_bC_\chi+o(1)$ and
$\rho_N\lambda_N\sim C_\chi LL/(2\log b\sqrt{\log N})$. (W0) and (W1) hold by the support
asymptotic, the hyperbola bound and $c(n)\le\tau(n)$. (W2) follows by summing
\cref{lem:pairs} over $|h|\le H_N=O(\ell_N)$, which removes $O(N\ell_NLL/\log N)=o(S(N))$
indices. For $b=2^k$, (W3) holds with $\omega=0$ by \cref{lem:valuations}(a). Otherwise
$u_b(n)\in[\log_bc(n)-V_b(n),\log_bc(n)-V_b(n)+1]$, so \cref{lem:firstmoment} gives
$W(N)\le(\log2+o(1))S(N)\lambda_N$, while \cref{lem:valuations}(b) and \cref{lem:concentration}
give the lower tail; so (W3) holds with $\omega=\log2<1$. For $\Theta_3(1/b)^2$, multiplication
by $4$ changes $V_b,U_b,u_b$ by $O(1)$ and adding $1$ changes nothing, so the proof applies
verbatim; in base $2$, $|D_N(\Theta_3(1/2)^2)-D_N(L_{\chi_{-4}}(2))|\le2/N$.
\end{proof}

In base $2$ and for $\chi_{-4}$ the constant is $K/(2\log2)$. The cost decomposition is
transparent: every isolated source needs $\frac12\log_b\log N$ digits of precision, and when the
radix cannot absorb the powers of $2$ in $c(n)\approx2^{\omega_{\rm split}(n)}$ the coefficient
itself occupies another $\frac{\log2}2\log_b\log N$ digits. At $N=2^{17}$ the exact deficit of
$L_{\chi_{-4}}(b)$ is $0.59476$, $0.59045$, $0.39883$ for $b=2,3,4$; since $\log\log N<2.5$
there, these finite values do not test the constants.

\subsection{The prime-density class}\label{subsec:primedensity}

Let $c$ be multiplicative with $0\le c(p^e)\le e+1$, $c(p)=2$ on a set of primes $\mathcal S$
of natural density $\delta\in(0,1)$ with
$\sum_{p\le x,\,p\in\mathcal S}1/p=\delta\log\log x+M_{\mathcal S}+o(1)$, and $c(p)=0$ at the
other primes not dividing a fixed modulus $q_0$. Put $\kappa=1-\delta$, let
$S(N)\sim C_fN/(\log N)^\kappa$ be the support count (\cref{thm:meanvalues}(b)),
$L_f(b)=\sum_nc(n)b^{-n}$, and
\begin{equation}\label{eq:PhiN}
\Phi_N=\frac{C_f\log\log N}{\log b\,(\log N)^\kappa},\qquad \omega=\frac{\delta\log2}\kappa .
\end{equation}
For $b=2^sb'$ with $b'$ odd let $\beta_b=\log b/\log b'$ ($\beta_b=1$ for odd $b$,
$\beta_b=\infty$ for $b=2^k$), $\gamma=1-1/\beta_b$, and
$g_\beta(y)=\sum_{i\ge0}(y\gamma^i/\beta-1)_+$, with $g_1(y)=(y-1)_+$.

\begin{lemma}[$2$-adic run barrier]\label{lem:runbarrier}\evidence{Paper proof}
Let $b=2^sb'$ with $s\ge1$, $b'>1$ odd, and $U\ge1$ with $b\nmid U$. If the base-$b$ digits of
$U$ in positions $s_0,\dots,s_0+\ell-1$, with $\ell\ge1$ (position $0$ is the units digit), are all $0$ or all
$b-1$, and $s_0+\ell\le v_2(U)/s$, then $s_0\ge1$ and $\ell<s_0/(\beta_b-1)$.
\end{lemma}

\begin{proof}
Put $J=s_0+\ell$ and $r=U\bmod b^J$. Since $sJ\le v_2(U)$, both $r$ and $b^J-r$ are positive
multiples of $2^{sJ}$, so $\min(r,b^J-r)\ge2^{sJ}$. With zero digits $r<b^{s_0}$; with digits
$b-1$, $b^J-r=b^{s_0}-(U\bmod b^{s_0})<b^{s_0}$, and $s_0=0$ would give $b^J-r=1$. Hence
$2^{sJ}<b^{s_0}$, that is, $\ell s\log2<s_0\log b'$, and $s\log2/\log b'=\beta_b-1$.
\end{proof}

\begin{lemma}[Packing]\label{lem:packing}\evidence{Paper proof}
Let $\beta>1$, $\kappa_0>0$, and let $[s_i,s_i+\ell_i)$ be pairwise disjoint intervals in
$[0,x]$ with $0<\ell_i<s_i/(\beta-1)$. Then $\sum_i(\ell_i-\kappa_0)_+\le\kappa_0g_\beta(x/\kappa_0)$,
and this is sharp.
\end{lemma}

\begin{proof}
Scale to $\kappa_0=1$. The function $g=g_\beta$ is continuous, nondecreasing, $1$-Lipschitz
and satisfies $g(y)=(y/\beta-1)_++g(\gamma y)$. Induct on the number of intervals: let
$[s,s+\ell)$ have the largest left endpoint; the others contribute at most $g(s)$, and
$\ell<(x-\ell)/(\beta-1)$ gives $\ell<x/\beta$. If $\ell\le1$ the total is at most $g(x)$;
otherwise it is at most $(\ell-1)+g(x-\ell)\le x/\beta-1+g(\gamma x)=g(x)$, because
$\ell\mapsto\ell+g(x-\ell)$ is nondecreasing. Geometric families with tops $\gamma^ix$ show
sharpness.
\end{proof}

\begin{theorem}[Prime-density class]\label{thm:primedensity}\evidence{Lean-supported paper proof}
\leavevmode
\begin{enumerate}[label=\textup{(\alph*)},nosep]
\item If $b=2^k$, then $1-D_N(L_f(b);b)\sim\kappa\Phi_N$ for every $\delta\in(0,1)$.
\item If $b$ is not a power of $2$, then unconditionally
\[
\kappa\bigl(1+\omega-g_{\beta_b}(\omega)\bigr)\le\liminf_N\frac{1-D_N}{\Phi_N}\le\limsup_N\frac{1-D_N}{\Phi_N}\le\kappa+\delta\log2 .
\]
In particular $1-D_N\sim(\kappa+\delta\log2)\Phi_N$ whenever
$\delta\le\delta_b^*=\beta_b/(\beta_b+\log2)$. For odd $b$, $\delta_b^*=1/(1+\log2)=0.590616\ldots$
and the lower value is $\kappa(1+\min(\omega,1))$; for $b=6,10,12$, $\delta_b^*=0.701754$,
$0.673632$, $0.765433$.
\end{enumerate}
\end{theorem}

\begin{proof}
The changes to the proof of \cref{thm:lambertradix} are the following. The mass is
$\ll X(\log X)^{2\delta-1}$ by \cref{thm:meanvalues}(a), so (W1) holds with
$B=\max(0,2\delta-1)$. The pair sieve uses the primes outside $\mathcal S$ in dimension $2\kappa$
and gives $\ll x(\log\log x)^{2\kappa}/(\log x)^{2\kappa}$ pairs per shift; summed over
$|h|\le H$ this is $o(S(N))$ for every $H\le(\log N)^\kappa/(LL)^{2\kappa+1}$, so distinct
sources are isolated far beyond the radius needed, at every density. On the support
$\log c(n)=\omega_{\mathcal S,1}(n)\log2+\log M(n)$, where $\omega_{\mathcal S,1}$ counts primes of
$\mathcal S$ with exponent exactly $1$ and $M(n)=\prod_{p^e\|n,\,e\ge2\text{ or }p\mid q_0}c(p^e)$ is
tight on the support; the first moment and concentration go through with centre $\delta\log2$
and weight $(1-\log p/\log N)^{-\kappa}$. For odd $q_0\mid b$, $v_{q_0}(c(p))=0$, so
\cref{lem:valuations} holds verbatim. Hence (W3) holds with $\omega=0$ for $b=2^k$ and with
$\omega=\delta\log2/\kappa$ otherwise, and $\lambda_N=\kappa\ell_N+O(1)$; \cref{thm:widthlaw}
gives (a), the upper bound in (b), and the lower value $\kappa(1+\min(\omega,1))$.

For even composite $b$, add to the good set the condition (G7) $v_2(U(n))\ge s\,u(n)$. It holds
whenever $\omega_{\mathcal S,1}(n)\log b'/\log b\ge s(\log_bM(n)+1)$, which fails for $o(S(N))$
indices by tightness of $M$ and the lower tail of $\omega_{\mathcal S,1}$. For a good source,
an internal position $\lambda'<j\le u$ with $\dist{Y/b^j}<\eta$ forces the $\lambda'$ digits of
$U$ below position $j$ to be all $0$ or all $b-1$; a maximal such run of length $\ell$ yields at
most $\ell-\lambda'+1$ near positions. By (G7) every run lies below $v_2(U)/s$, so
\cref{lem:runbarrier,lem:packing} (with $\kappa_0=\lambda'-1$, $x=u$) bound the number of near
internal positions by $(\lambda'-1)g_\beta(u/(\lambda'-1))$. Each good source therefore marks at
least $(1+\omega-g_\beta(\omega+\varepsilon)-2\varepsilon-o(1))\lambda_N$ far shifts in its
arrival-free interval, and \cref{lem:circle} concludes as in \cref{thm:widthlaw}. Finally
$g_\beta(\omega)=0$ if and only if $\omega\le\beta$, that is, $\delta\le\delta_b^*$.
\end{proof}

For odd $b$ and $\delta>\delta^*$ the exact constant is decided by how often the integers
$M\cdot2^m$ have base-$b$ runs of $0$ or of $b-1$ whose length is a fixed proportion of their
number of digits, for $m$ in windows of width $\asymp\sqrt{LL}$; linear forms in logarithms and
$q_0$-adic logarithms exclude only $o(\lambda_N)$ positions, and we know no result that decides
it.

\subsection{Euler's square in every radix}\label{subsec:squareradix}

\begin{theorem}[Euler's square]\label{thm:squareradix}\evidence{Lean-supported paper proof}
Granting \cref{thm:etainput},
\[
1-D_N(E_b^2;b)\sim\begin{cases}\dfrac{3C_\eta}{4\log b}\dfrac{LL}{(\log N)^{3/4}},&b=2^k,\\[2ex]
\dfrac{C_\eta(3+\log2)}{4\log b}\dfrac{LL}{(\log N)^{3/4}},&\text{otherwise.}\end{cases}
\]
\end{theorem}

\begin{proof}
Apply \cref{thm:widthlaw} to $a_2$ with $\kappa=\frac34$, $B=0$. (W1) follows from
\cref{thm:square}(b) and $|a_2(n)|\le\tau(12n+1)$. For (W2), if $a_2(n)a_2(n+h)\ne0$ then
$12n+1$ and $12n+12h+1$ lie in $\mathcal S$; discard the $n$ for which $p^2$ divides either
number for some $p>z_0=(\log N)^2$ with $p\not\equiv1\pmod{12}$, and sift the rest by two
classes modulo each such $p\le\sqrt N$, $p\nmid h$; these primes have density $\frac34$, and
Brun's sieve in dimension $\frac32$ gives $\ll N(\log z_0/\log N)^{3/2}$ uniformly in $h$ (or
use \cref{thm:S}). For (W3), $\log|a_2(n)|=\sum_{p^e\|12n+1,\,p\equiv1\,(12)}\log(e+1)$ has first
moment $(\frac{\log2}4+o(1))S_\eta(N)LL$ by the decomposition of \cref{lem:firstmoment} with
weight $(1-\log p/\log x)^{-3/4}$; for $b=2^k$ the remainder is at most
$2^{k-1}\mathrm{odd}(|a_2|)$, of bounded mean, and for odd $q_0\mid b$, $q_0\mid e+1$ forces
$e\ge2$. So $\omega=0$ for $b=2^k$ and $\omega=\frac{\log2/4}{3/4}=\frac{\log2}3<1$ otherwise.
Half of the coefficients are negative; by \cref{rem:widthlaw} the signs are harmless.
\end{proof}

\subsection{Binary quadratic forms}\label{subsec:forms}

Let $Q=gQ_0$ with $Q_0$ primitive of discriminant $\Delta_0$, $r_Q(n)=r_{Q_0}(n/g)$, and
$S_Q(N)\sim C_QN/\sqrt{\log N}$. Put $\Psi_N=C_QLL/(2\log b\sqrt{\log N})$. We use: the mass
$\sum_{n\le X}r_Q(n)=2\pi X/\sqrt{|\Delta|}+O(\sqrt X)$ (\cref{lem:lattice}); the pointwise bound
$r_{Q_0}(m)\le C\,2^{\omega_{\rm split}(m)}\prod_{p^e\|m}(\lfloor e/2\rfloor+1)=m^{o(1)}$; the
support condition $v_p(n)$ even whenever $(\Delta_0/p)=-1$ (\cref{lem:bernays}); and, for
fundamental $\Delta$ and $Q$ in the ideal class $C$, $r_Q=w\,a_C$, where $a_C(n)$ counts ideals
of norm $n$ in $C$ and $\sum_Ca_C=1*\chi_\Delta$.

\begin{proposition}[Lower bound for every form]\label{prop:formlower}\evidence{Lean-supported paper proof}
For every positive definite integral binary $Q$ and every $b\ge2$,
$1-D_N(\Theta_Q(1/b);b)\ge(1-o(1))\,\frac{S_Q(N)}N\cdot\frac12\log_b\log N$.
\end{proposition}

\begin{proof}
Run the lower-bound half of \cref{thm:widthlaw} with $\rho_N=S_Q(N)/N$, $H_V=\lceil2\ell_N\rceil$,
$R=\lceil\lambda_N+2\ell_N\rceil$ and conditions (G1)--(G4) only. (G1) fails for
$\ll N/(\log N)^2$ indices by the linear mass. For (G2), two sources $gj<gj'$ within distance
$H$ have $j,j+h'\in\{v_p\text{ even at inert }p\}$ with $1\le h'\le H$; Brun's sieve in dimension
$1$, as in \cref{lem:pairs}, removes $o(S_Q(N))$ indices. (G3) and (G4) follow as before, using
$\sum_{s<M}|A(s)|\ll M$. Every good source marks the positions $1\le j\le\lambda'$ by
\cref{lem:onesource}, in disjoint gaps, and \cref{lem:circle} with
$\zeta=(\log N)^{-1/2}$ and $\eta=\zeta+2\rho_N\lambda_N$ gives the claim. No multiplicativity is
used.
\end{proof}

\begin{lemma}[Order-$4$ annihilation]\label{lem:annihilation}\evidence{Paper proof}
Let $\Delta<0$ be fundamental, $G=\Cl(\Delta)=G_2\times G_{\rm odd}$, $h_0=|G^2|$, and for
$n\in\supp(1*\chi_\Delta)$ let
$X_n=\prod_{p^e\|n\text{ split}}\sum_{i=0}^e[\mathfrak p]^{2i-e}\prod_{p^e\|n\text{ ramified}}[\mathfrak p]^e\in\Z[G]$,
so that $a_C(n)$ is the coefficient of $C$ in $X_n$. Suppose that for every character $\chi$
of $G$ with $4\mid\ord\chi$ some split $p\|n$ has $\chi(\mathfrak p)^2=-1$. Then
$|G_2^2|\,a_C(n)=\ind[n\in\mathrm{genus}(C)]\,b_{C_{\rm o}}(n)$, where $b_c(n)$ is the coefficient
of $c$ in the image of $X_n$ in $\Z[G_{\rm odd}]$. The hypothesis fails for
$\ll S_\Delta(N)(\log N)^{-c}$ integers $n\le N$, $c=c(\Delta)>0$.
\end{lemma}

\begin{proof}
If $\chi(\mathfrak p)=\pm i$ and $p\|n$, the local factor $\chi(\mathfrak p)+\chi(\mathfrak p)^{-1}$
vanishes, so only characters of $G/G_2^2\cong(G/G^2)\times G_{\rm odd}$ survive, and orthogonality
on that group gives the identity. The exceptional set is the support of a multiplicative
function vanishing on a Chebotarev set of primes of positive density; apply
\cref{thm:meanvalues}(a).
\end{proof}

\begin{lemma}[Strassmann]\label{lem:strassmann}\evidence{Paper proof}
Let $\mathcal O$ be the ring of integers of a finite extension of $\Q_\ell$,
$\beta_1,\dots,\beta_r\in\mathcal O^\times$ with no ratio $\beta_i/\beta_j$ $(i\ne j)$ a root of
unity, and $\mathcal C\subset\mathcal O^r$ compact with all coordinates nonzero. For
$A\in\mathcal C$ put $f_A(j)=\sum_iA_i\beta_i^j$. Then there is $T\ge1$ such that for each $H$ the
event $v(f_A(j))\ge H$ depends only on $j$ modulo $T\ell^{a(H)}$, and its density tends to $0$ as
$H\to\infty$, uniformly in $A\in\mathcal C$.
\end{lemma}

\begin{proof}
Choose $T$ with $\beta_i^T\in1+\mathfrak m^\nu$, where $\log$ and $\exp$ are inverse isometries.
On $j=j_0+Tk$, $f_A(j)=\sum_iA_i\beta_i^{j_0}\exp(kL_i)$ with pairwise distinct
$L_i=\log\beta_i^T$, an analytic function of $k\in\Z_\ell$ whose derivatives at $0$ cannot all
vanish (Vandermonde). By Strassmann's theorem \cite{Strassmann1928} it has finitely many zeros,
so the Haar measure $\phi_H(A)$ of $\{v(f_A)\ge H\}$ decreases to $0$; each $\phi_H$ depends only
on $A$ modulo $\mathfrak m^H$, hence is continuous, and Dini's theorem makes the convergence
uniform on $\mathcal C$.
\end{proof}

\begin{theorem}[Class groups]\label{thm:classgroups}\evidence{Lean-supported paper proof}
Let $\Delta<0$ be a fundamental discriminant, $Q$ a primitive form of discriminant $\Delta$ in any
class, and $G=\Cl(\Delta)=G_2\times G_{\rm odd}$.
\begin{enumerate}[label=\textup{(\alph*)},nosep]
\item For every odd prime $q_0$, $v_{q_0}(r_Q(n))$ is tight on the support of $r_Q$.
Consequently $1-D_N(\Theta_Q(1/b);b)\sim(1+\log2)\Psi_N$ for every $b$ that is not a power of $2$.
\item For $b=2^k$, $1-D_N(\Theta_Q(1/b);b)\sim(1+\log2\cdot\mathbb E_\Delta[1-1/d])\Psi_N$, where
$d=d(n)$ is the largest order of a character of $G_{\rm odd}$ not annihilated by a prime power
exactly dividing $n$, and $\mathbb E_\Delta$ is its limiting law on the support, an explicit Euler
product. If $G$ is a $2$-group the factor is $1$, and $r_Q(n)=w(1*\chi_\Delta)(n)/h_0$ for almost
all $n$ in the genus of $Q$. If $G_{\rm odd}\cong\Z/\ell$ with $\ell$ prime, the factor is
$1+(1-\frac1\ell)(1-\pi_{J,\ell})\log2$ with
\[
1-\pi_{J,\ell}=\prod_{\substack{p\text{ split}\\ [\mathfrak p]\notin G_2}}\Bigl(1-\frac{(1-1/p)\,p^{1-\ell}}{1-p^{-\ell}}\Bigr);
\]
for $\Delta=-23$ $(\Z/3)$ it equals $1.362117\ldots$, for $\Delta=-47$ $(\Z/5)$ $1.531999\ldots$,
and for $\Delta=-71$ $(\Z/7)$ $1.588879\ldots$ (Euler products truncated at $p\le2\cdot10^5$).
\end{enumerate}
\end{theorem}

\begin{proof}
Write $a=1*\chi_\Delta$, so that $r_Q=w\,a_C$ and $a_C\le a$. Both parts use
\cref{thm:widthlaw} in its mixture form (\cref{rem:widthlaw}), whose hypotheses are: the
support asymptotic (Bernays), isolation (as in \cref{prop:formlower}), convergence in
probability of $u_b(n)/\lambda_N$ to a law on $[0,1)$, and the matching first-moment bound for
$W(N)$. The lower value $1$ is \cref{prop:formlower}.

\emph{Three common facts.} (i) Genus characters show, as in \cref{lem:firstmoment}, that the first
moment of $\log a$ over the support of $r_Q$ is $(\frac{\log2}2+o(1))S_Q(N)LL$. (ii)
$a_C(n)=(1+o(1))a(n)/h_0$ in probability on the support of $r_Q$: by orthogonality on
$\widehat G$, the genus characters contribute exactly $a(n)/h_0$, and every character $\psi$
with $\psi^2\ne1$ contributes at most
$a(n)\prod_{p\|n\text{ split},\,[\mathfrak p]\notin\ker\psi^2}|\cos\arg\psi(\mathfrak p)|$; these
primes have positive Chebotarev density, so the product tends to $0$ in probability
(Rankin's argument of \cref{lem:concentration}). (iii) $\log_2a(n)=\omega_{\rm split}(n)+T(n)$
with $T(n)=\sum_{p^e\|n,\ e\ge2,\ p\text{ split}}(\log_2(e+1)-1)$, a tight (bounded in
probability, not uniformly bounded) quantity, and $\omega_{\rm split}(n)=(\frac12+o(1))LL$ in
probability. Hence $\log_2a_C(n)=(\frac12+o(1))LL$ in probability. Finally, the exceptional
set $\mathcal E_N$ of \cref{lem:annihilation} has total width $o(S_Q(N)\lambda_N)$ for
\emph{every} $G$: by the pointwise bound on $r_{Q_0}$ and \cref{thm:meanvalues}(a) with
exponent $s=1/LL$, $\sum_{n\le N}r_Q(n)^{1/LL}\ll S_Q(N)$, hence
$\sum(\log_br_Q(n))^2\ll S_Q(N)(LL)^2$, and Cauchy--Schwarz gives
$\sum_{n\in\mathcal E_N}u(n)\le|\mathcal E_N|+(|\mathcal E_N|\sum(\log_br_Q)^2)^{1/2}
\ll S_Q(N)LL(\log N)^{-c/2}$.

(a) Off $\mathcal E_N$, fix a local pattern $\pi$ at the primes $\le P_0$ (outside probability
$\varepsilon$) and a genus. With $\eta_c=c+c^{-1}$ for $c\in G_{\rm odd}$, so that
$\psi(\eta_c)=\psi(c)+\psi(c)^{-1}$ is $2$ or a cyclotomic unit, \cref{lem:annihilation} and
Fourier inversion on $G_{\rm odd}$ give, after pairing $\psi$ with $\bar\psi$,
\[
h_0a_C(n)=\sum_{\{\psi,\bar\psi\}}\epsilon_\psi\,\psi(\eta_{C_{\rm o}})\,w_\psi(\pi)\prod_{c\in G_{\rm odd}/\pm}\psi(\eta_c)^{m^{\rm o}_c(n)},
\]
over the pairs of characters surviving $\pi$, where $\epsilon_1=\frac12$, $\epsilon_\psi=1$
otherwise, $w_\psi(\pi)\ne0$ collects the primes $\le P_0$ and the ramified primes, and
$m^{\rm o}_c(n)$ counts split $p\|n$, $p>P_0$, with $[\mathfrak p]_{\rm o}\in\{c,c^{-1}\}$.
Split $m^{\rm o}_c=m_{(1,c)}+m^{\rm rest}_c$, where $m_{(1,c)}$ counts the primes whose class is
$(1,c^{\pm1})$, which lies in $G^2$, and $m^{\rm rest}_c$ the primes with nontrivial
$2$-component. Condition on all $m^{\rm rest}_c$: their factors are units at every prime above
an odd $q_0$, so they are absorbed into coefficients ranging over a compact set with nonzero
coordinates, and moving the coordinates $m_{(1,c)}$ does not change the genus. By
Selberg--Delange for Hecke $L$-functions of class-group characters, these coordinates are
asymptotically uniform modulo every $M$, given the conditioned data. If
$\{\psi,\bar\psi\}\ne\{\psi',\bar\psi'\}$, some $\psi(\eta_c)/\psi'(\eta_c)$ is not a root of
unity (squaring is a bijection of $G_{\rm odd}$, and characters are linearly independent), and a
lattice direction $\mathbf w$ with first coordinate $1$ separates all pairs simultaneously.
Along the lines $(0,\mathbf u)+j\mathbf w$ the sum is $f_A(j)$ of \cref{lem:strassmann} over the
completion of $\Q(\zeta_{|G_{\rm odd}|})^+$ at a prime above $q_0$, so
$\mathbb P(v_{q_0}(h_0a_C(n))\ge H)\le2\varepsilon$ for $H\ge H(\varepsilon)$. Hence
$V_b\le v_{q_0}(r_Q)$ is tight. With fact (ii)--(iii), $u_b(n)/\lambda_N\to\log2$ in
probability, and fact (i) with $a_C\le a$ gives $W(N)\le(\log2+o(1))S_Q(N)\lambda_N$; the
constant is $1+\log2$.

(b) For a surviving character $\psi$ of order $d\ge3$, the same expansion shows that the
$2$-adic valuation of the $\psi$-component equals the number $m_{\ker\psi}(n)$ of split $p\|n$
with $[\mathfrak p]_{\rm o}\in\ker\psi$, up to a bounded term; these primes form a proportion $1/d$
of the split primes (Chebotarev in the Hilbert class field), so $m_{\ker\psi}(n)=(\frac1{2d}+o(1))LL$
in probability. On the annihilated set $v_2(a_C(n))\ge\min_{\psi}m_{\ker\psi}(n)-v_2(h_0)$ over
the survivors, always. For the matching upper bound, all constants are fixed first and the
minimizing kernel is selected last: \cref{lem:strassmann} ($2$-adically, along the
coordinate of a class pair in $G^2$) bounds the probability of cancellation within each fixed
kernel, the uniformity modulo $M$ bounds the probability that two kernels tie within a fixed
margin, and a union bound over the finitely many kernels gives
$v_2(a_C(n))=(\frac1{2d(n)}+o(1))LL$ in probability. With facts (ii)--(iii) this gives
$u_{2^k}(n)/\lambda_N\to(1-1/d(n))\log2$ in probability. The survivors are determined by local
events at split $p^e\|n$ with $e\ge2$ (a nontrivial $\psi$ with $\psi(\mathfrak p)$ of order $o$
dies exactly when $o\mid e+1$), and on the support the exponents at finitely many split primes
are asymptotically independent with $\mathbb P(v_p=e)=(1-\frac1p)p^{-e}$; this gives the law of
$d$ and, for $G_{\rm odd}\cong\Z/\ell$, the stated product. For the first moment,
$u_{2^k}\le(\log_2a-v_2(a_C))/k+O(1)$, the lower bound for $v_2$ holds off a set of $o(S_Q)$
indices, $1/d$ is bounded, and the exceptional set has negligible width by the common facts, so
$W(N)\le(\log2\,\mathbb E_\Delta[1-1/d]+o(1))S_Q(N)\lambda_N$. The limit law of $u/\lambda_N$ is
supported in $[0,(1-|G_{\rm odd}|^{-1})\log2]\subset[0,1)$, so the mixture form applies. For
$2$-groups, $G_{\rm odd}=1$ and \cref{lem:annihilation} gives $r_Q=w(1*\chi_\Delta)/h_0$ on the
genus of $Q$ off $\mathcal E_N$.
\end{proof}

\Cref{thm:D}(iii) follows from \cref{prop:formlower,thm:classgroups}. For class number one,
$r_Q=w(1*\chi_\Delta)$ and \cref{thm:lambertradix} applies directly; for one class per genus
($h_0=1$, for example $\Delta=-20,-24$), $r_Q=w(1*\chi_\Delta)\ind_{\rm genus}$. For non-fundamental
or imprimitive forms only the lower value $1$ and the order are proved here.

\section{Theta values, weight-one forms and periodic weights}\label{sec:theta}

\subsection{Binary theta values}\label{subsec:thetaupper}

Let $Q(m,n)=\alpha m^2+\beta mn+\gamma n^2$ be positive definite and integral with discriminant
$\Delta=\beta^2-4\alpha\gamma<0$, $Q=gQ_0$ with $Q_0$ primitive of discriminant $\Delta_0$, and
$r_Q(j)=\cnt\{(m,n):Q(m,n)=j\}$, so that $\Theta_Q(1/b)=F_b(r_Q)$.

\begin{lemma}[Lattice points]\label{lem:lattice}\evidence{Paper proof}
$\sum_{j\le X}r_Q(j)=2\pi X/\sqrt{|\Delta|}+O_Q(\sqrt X)$ and
$r_Q(j)\le2(2\sqrt{4\alpha j/|\Delta|}+1)$.
\end{lemma}

\begin{proof}
The ellipse $\{Q\le X\}$ has area $2\pi X/\sqrt{|\Delta|}$, and unit squares centred at its lattice
points lie between the ellipses shrunk and enlarged by $\sqrt2/2$. For the pointwise bound,
$Q(m,n)\ge(|\Delta|/4\alpha)n^2$ and each $n$ admits at most two values of $m$.
\end{proof}

\begin{lemma}[Even exponents]\label{lem:bernays}\evidence{Paper proof}
If $(\Delta_0/p)=-1$ then $v_p(Q_0(m,n))$ is even for all $(m,n)\ne(0,0)$. Hence
$B_Q(X)=\cnt\{1\le j\le X:r_Q(j)>0\}\ll_QX/\sqrt{\log X}$.
\end{lemma}

\begin{proof}
For odd $p$ with $(\Delta_0/p)=-1$, the identity $4\alpha Q_0=(2\alpha m+\beta n)^2-\Delta_0n^2$
(or its analogue with $\gamma$) shows that $Q_0$ is anisotropic modulo $p$, so $p\mid Q_0(m,n)$
forces $p\mid m$, $p\mid n$; iterate. For $p=2$, $\Delta_0\equiv5\pmod8$ forces $\alpha,\beta,\gamma$
odd, and $m^2+mn+n^2$ has no nontrivial zero over $\mathbb F_2$. The represented set therefore
lies in a multiplicative set of prime density $\frac12$, and \cref{thm:meanvalues}(a) applies;
$B_Q(X)=B_{Q_0}(X/g)$.
\end{proof}

\begin{theorem}[Theta values]\label{thm:thetaupper}\evidence{Lean-supported paper proof}
For every positive definite integral binary $Q$ and every $b\ge2$, $\Theta_Q(1/b)\in K_b$ and
$1-D_N(\Theta_Q(1/b);b)\ll_{Q,b}\LL/\sqrt{\log N}$.
\end{theorem}

\begin{proof}
\Cref{lem:lattice} gives linear mass and \cref{lem:bernays} gives $S(X)\ll X/\sqrt{\log X}$;
apply \cref{thm:leantransfer} for membership and \cref{thm:largermass} for the rate.
\end{proof}

\begin{proposition}[Extensions]\label{prop:extensions}\evidence{Lean-supported paper proof}
The conclusion of \cref{thm:thetaupper} holds for:
\begin{enumerate}[label=\textup{(\alph*)},nosep]
\item integral shifted lattices $\sum_{(m,n)}b^{-f(m,n)}$ with
$f=Q+l_1m+l_2n+l_0$, $l_i\in\Z$;
\item congruence subsums $\sum_{(m,n)\equiv(r,s)\,(D)}b^{-Q(m,n)}$;
\item weighted sums $\sum\chi(m,n)b^{-Q(m,n)}$ with integer weights $|\chi|\le B$;
\item finite $\Z[1/b]$-linear combinations of the above.
\end{enumerate}
\end{proposition}

\begin{proof}
(a) With $U=2\gamma l_1-\beta l_2$ and $V=2\alpha l_2-\beta l_1$,
$|\Delta|^2f(m,n)=Q(|\Delta|m+U,|\Delta|n+V)+E_0$ with $E_0\in\Z$, so the values of $f$ lie in an affine
image of the represented set, with count $\ll X/\sqrt{\log X}$; the mass is $O(X)$ by area, and the
finitely many nonpositive values contribute integers. (b) and (c) have coefficients bounded by
$Br_Q(j)$ on the same support. (d) is \cref{lem:linear}.
\end{proof}

Harmonic polynomial weights of degree $\nu\ge1$ give cusp forms of weight $1+\nu$ whose mass
grows like $X^{1+\nu/2}$; they fall outside the method (\cref{sec:limits}).

\begin{proposition}[Arithmetic nature]\label{prop:arith}\evidence{Paper proof}
Let $a\ge0$ be an integer sequence with infinitely many nonzero terms, $A(X)=O(X)$ and
$S(X)=o(X)$. Then $F_b(a)$ is irrational. In particular every $\Theta_Q(1/b)$ and every
$L_\chi(b)$ with $\chi$ real nonprincipal is irrational. Every $\Theta_Q(1/b)$, every $L_\chi(b)$
with $\chi$ odd, and $E_b$ are transcendental
\cite{Nesterenko1996,DNNS1996,Bertrand1997,HamiehMurty2016}.
\end{proposition}

\begin{proof}
A rational element of $K_b$ has an eventually periodic orbit ending at a fixed point
\cite{DiscrepancyCompanion}; since the orbit concentrates at $0$, the number would lie in
$\Z[1/b]$, and the positive tails $\tau_n=\sum_{m>n}a(m)b^{n-m}$ would be integers $\ge1$ for all
large $n$. But at least $X-L\,S(X+L)\ge X/2$ of the $n\le X$ have no support in $(n,n+L]$, and their
tails average $\ll b^{-L}$ by \cref{lem:meantail}, so some tail lies in $(0,1)$. For transcendence:
$\Theta_Q$ is a weight-one modular form with integer coefficients and $q=1/b$ is algebraic, so
\cite[Thm.~3.1]{HamiehMurty2016} applies; for odd $\chi$, $L_\chi(b)$ is the value at $q=1/b$ of
the $q$-analogue $L_q(1,\chi)$, transcendental by \cite[Thm.~1.1]{HamiehMurty2016} for primitive
$\chi$ and by Theorem~3.1 there in general (the case of even $\chi$ is open); and $E_b$ is
transcendental because
$1728\,b^{-1}E_b^{24}=Q(1/b)^3-R(1/b)^2$ in Ramanujan's notation, while $Q(1/b)$ and $R(1/b)$ are
algebraically independent by Nesterenko's theorem.
\end{proof}

\subsection{Weight one}\label{subsec:weightone}

\begin{theorem}[Weight-one inputs]\label{thm:weightoneinput}\evidence{External input}
$M_1(\Gamma_1(M))$ is spanned by $g(t\tau)$ for newforms $g$ of weight one and level $M'\mid M$,
$t\mid M/M'$, and by Eisenstein series $E_1^{\psi,\varphi}(t\tau)$ with $\psi\varphi$ odd
\cite[Thms.~4.8.1, 5.8.3]{DiamondShurman}. A weight-one newform corresponds to an odd irreducible
Artin representation $\rho$ with finite image and $a_g(p)=\tr\rho(\mathrm{Frob}_p)$ for $p\nmid M'$
\cite{DeligneSerre1974}; the primes with $\tr\rho(\mathrm{Frob}_p)=0$ have a positive Chebotarev
density, equal to the proportion of involutions in the projective image
\cite[Thm.~18]{Serre1981}.
\end{theorem}

\begin{theorem}[Weight one]\label{thm:weightone}\evidence{Lean-supported paper proof}
Let $f=\sum_{n\ge0}a(n)q^n\in M_1(\Gamma_1(M))$ with all $a(n)\in\Z$. Let $\alpha_f>0$ be the minimum,
over the newform constituents $g$ of $f$, of the density of primes with $\tr\rho_g(\mathrm{Frob}_p)=0$,
and over the Eisenstein constituents $E_1^{\psi,\varphi}$, of $1/m$, where $m$ is the order of
$\varphi\bar\psi$. Then $S_f(X)\ll X/(\log X)^{\alpha_f}$, $A_f(X)\ll X\log X$, and for every $b\ge2$,
$f(1/b)\in K_b$ with $1-D_N(f(1/b);b)\ll_{f,b}\LL/(\log N)^{\alpha_f}$.
\end{theorem}

\begin{proof}
Write $f$ as a finite combination of the basis of \cref{thm:weightoneinput}. For a newform the
Frobenius eigenvalues are roots of unity, so $|a_g(n)|\le\tau(n)$, and for $p$ with trace-zero
Frobenius the Hecke recursion $a(p^{e+1})=a(p)a(p^e)-\chi(p)a(p^{e-1})$ kills all odd powers; so
the support lies in a multiplicative set of prime density $1-\delta_0(g)$. For an Eisenstein series,
$\sigma(n)=\sum_{d\mid n}\psi(n/d)\varphi(d)$ and $\sigma(p)=\psi(p)(1+\omega(p))$ with
$\omega=\varphi\bar\psi$. Since $\psi\varphi$ is odd, $\omega(-1)=-1$, so $\omega$ has even order $m$,
$\omega(p)=-1$ on primes of density $1/m$, and there
$\sigma(p^{2e+1})=\psi(p)^{2e+1}\sum_{i=0}^{2e+1}(-1)^i=0$. Hence every constituent is lacunary, the
support of $a_f$ lies in a finite union of dilates of these sets, and $|a_f(n)|\le C_f\tau(n)$.
\Cref{thm:meanvalues}(a), \cref{thm:leantransfer} (logarithmic mass, $\gamma=1$) and
\cref{thm:largermass} conclude; integrality of all $a(n)$, including $a(0)$, makes the heads integers.
\end{proof}

The exponent $\alpha_f$ is $\frac12$ for $\Theta_Q$, $\frac34$ for the form $\eta(12\tau)^2$ behind
$E_b^2$, and would be $\frac14$ for an integer-coefficient form of tetrahedral or icosahedral type
(for instance the trace over the Hecke field of an exotic newform). For a rational form with
denominator $d$ only $d\,f(1/b)\in K_b$ follows. The mean of $|\tr|$ over the binary polyhedral group
is less than $1$ for every finite projective image, so every weight-one newform has $A_g(X)=o(X)$,
whereas an Eisenstein series with $m\ge4$ has $A(X)\asymp X(\log X)^{\kappa_m-1}$ with
$\kappa_m=\frac1m\sum_j|1+e^{2\pi ij/m}|>1$; the logarithmic-mass transfer absorbs this.

\subsection{Characters of even order and periodic weights}\label{subsec:periodic}

For a weight $w$ put $\mathcal L_w(b)=\sum_{m\ge1}w(m)/(b^m-1)=\sum_n(1*w)(n)b^{-n}$, and for a
Dirichlet character $\chi$ write $c_\chi=1*\chi$. Locally $c_\chi(p^e)=\sum_{i\le e}\chi(p)^i$,
which vanishes exactly when $\chi(p)\ne1$ has order dividing $e+1$.

\begin{theorem}[Even-order characters]\label{thm:evenorder}\evidence{Lean-supported paper proof}
Let $\chi$ have even order $h$. Then $\cnt\{n\le X:c_\chi(n)\ne0\}\sim C_\chi X(\log X)^{-1/h}$, and
$\sum_{n\le X}|c_\chi(n)|\asymp X(\log X)^{\kappa_h-1}$ with $\kappa_h=\frac2h\cot\frac\pi{2h}$; in
particular the mass is linear only for real $\chi$. If $\chi_1,\dots,\chi_J$ have even orders with
maximum $H$, $d_j\ge1$, and $a(n)=\sum_j\lambda_jc_{\chi_j}(n/d_j)$ is integer valued (with
$c(x)=0$ for $x\notin\Z$), then for every $b\ge2$, $F_b(a)\in K_b$ and
$1-D_N(F_b(a);b)\ll\LL/(\log N)^{1/H}$. If $\chi$ has odd order, the support of $c_\chi$ has
positive lower density.
\end{theorem}

\begin{proof}
For even $h$ the primes with $\chi(p)=-1$ form a coset of $\ker\chi$ of density $1/h$, and
$c_\chi(p^e)=0$ for odd $e$ at them; at all other primes $c_\chi(p)\ne0$. The support indicator is
multiplicative with prime density $1-1/h$, and \cref{thm:meanvalues} gives the asymptotic; the zeros
at exponents $e\ge2$ change only the constant. The mass follows from Wirsing's theorem applied to
$|c_\chi|$, since $\chi(p)$ is equidistributed in $\mu_h$. For the combination,
$\supp a\subseteq\bigcup_jd_j\supp c_{\chi_j}$ and $|a(n)|\le(\sum|\lambda_j|)\tau(n)$; apply
\cref{thm:leantransfer} (logarithmic mass) and \cref{thm:largermass}. For odd order,
$c_\chi(n)=\prod_{p\mid n}(1+\chi(p))\ne0$ for squarefree $n$ prime to the modulus.
\end{proof}

\begin{theorem}[Odd periodic weights]\label{thm:oddweights}\evidence{Lean-supported paper proof}
Let $w:\Z\to\Z$ be $M$-periodic and odd, $w(-m)=-w(m)$. If $p$ is a prime with $w(px)=-w(x)$ for all
$x$ (for instance $p\equiv-1\pmod M$), then $(1*w)(n)=0$ whenever $v_p(n)$ is odd. Consequently
$\cnt\{n\le X:(1*w)(n)\ne0\}\ll X(\log X)^{-1/\varphi(M)}$,
$\sum_{n\le X}|(1*w)(n)|\le2\max|w|\cdot X\log(X+2)$, and $\mathcal L_w(b)\in K_b$ for every $b\ge2$.
\end{theorem}

\begin{proof}
Write $n=p^em$ with $p\nmid m$. The divisors are $p^id'$ with $d'\mid m$, and
$w(p^id')=(-1)^iw(d')$, so $(1*w)(n)=\sum_{d'\mid m}w(d')\sum_{i=0}^e(-1)^i=0$ for odd $e$. For
$p\equiv-1\pmod M$ this cancellation is
\leanP{AntisymmetricLambert}{101}{dirichletOneConv_eq_zero_of_odd_padicVal}, the
mass bound is \leanP{AntisymmetricLambert}{271}{absMass_dirichletOneConv_le}, and the reduction of
membership to sparsity for bounded weights is
\leanP{AntisymmetricLambert}{510}{bounded_dirichletOneConv_tendsto_one_of_sparse}. The support
bound follows from \cref{thm:meanvalues}(a) with the primes $\equiv-1\pmod M$, which supplies the
sparsity hypothesis.
\end{proof}

For $w=\ind_r-\ind_{-r}$ with $\gcd(r,M)=1$ only odd characters occur in the Fourier expansion of
$w$, each of even order, and the largest such order is the Carmichael exponent $\lambda(M)$.
\Cref{thm:evenorder} applied to this expansion therefore gives the support bound
$\ll X(\log X)^{-1/\lambda(M)}$, an improvement on the exponent $1/\varphi(M)$ of
\cref{thm:oddweights}; we do not prove here that this exponent is sharp.

Let $\psi_q(x)=\Gamma_q'(x)/\Gamma_q(x)=-\log(1-q)+\log q\sum_{j\ge1}q^{jx}/(1-q^j)$ be the
$q$-digamma function \cite[\S1.10]{GasperRahman}, and
$A_{r,M}(b)=\sum_{m\equiv r\,(M)}1/(b^m-1)$.

\begin{theorem}[$q$-digamma reflections]\label{thm:qdigamma}\evidence{Lean-supported paper proof}
Let $b\ge2$, $M\ge2$, $1\le r<M$ and $q=b^{-M}$. Then
\[
\frac{\psi_q(r/M)-\psi_q(1-r/M)}{\log q}=A_{r,M}(b)-A_{M-r,M}(b)
=\sum_{n\ge1}\bigl(\tau_{r,M}(n)-\tau_{-r,M}(n)\bigr)b^{-n},
\]
where $\tau_{r,M}(n)$ counts divisors $d\equiv r\pmod M$ of $n$. This number lies in $K_b$. If
$2r\ne M$, $g=\gcd(r,M)$ and $M'=M/g$, then
$1-D_N\ll_{b,M,r}\LL/(\log N)^{1/\lambda(M')}$.
\end{theorem}

\begin{proof}
The $-\log(1-q)$ terms cancel, and termwise $q^{k+r/M}/(1-q^{k+r/M})=1/(b^{kM+r}-1)$; the identity is
\leanP{GeneralMassTransfer}{991}{qDigamma_reflection_eq_residueLambert}. The weight
$\ind_r-\ind_{-r}$ is odd, so \cref{thm:oddweights} gives the sparsity hypothesis of
\leanP{AntisymmetricLambert}{525}{qDigamma_reflection_tendsto_one_of_sparse}, which yields
membership. After the dilation by $g$, the remark above and \cref{thm:largermass} give the
rate.
\end{proof}

\begin{proposition}[Residue-class Lambert sums]\label{prop:residueclass}\evidence{Lean-supported paper proof}
Let $H\le G=(\Z/q)^\times$ with $G/H$ elementary abelian of order $2^k$, and
$A_{gH}(b)=\sum_{m\bmod q\in gH}1/(b^m-1)$, $E_q(b)=\sum_{(m,q)=1}1/(b^m-1)$. Then
$A_{gH}(b)-A_{g'H}(b)\in K_b$ and $2^kA_{gH}(b)-E_q(b)\in K_b$ for every $b\ge2$, with
$1-D_N\ll\LL/\sqrt{\log N}$. In base $2$, for $q=8$ and $r\in\{1,3,5,7\}$,
$4A_r-E_{\rm odd}=\chi_{-4}(r)L_{\chi_{-4}}+\chi_8(r)L_{\chi_8}+\chi_{-8}(r)L_{\chi_{-8}}$ and
$|D_N(A_r;2)-D_N(E_{\rm odd};2)|\le\frac2N+\sum_{\chi\ne\chi_0}(1-D_N(L_\chi;2))\ll\LL/\sqrt{\log N}$;
so $A_1,A_3,A_5,A_7$ and $E_{\rm odd}$ share all limit points of $D_N$.
\end{proposition}

\begin{proof}
The characters trivial on $H$ are real, and $\ind_{gH}=2^{-k}\sum_{\chi|_H=1}\chi(g)\chi$, so both
differences are Lambert series of integer weights whose divisor sums are dominated by
$\sum_{\chi\ne\chi_0}1*\chi$; apply \cref{thm:leantransfer,thm:largermass}. The mod-$8$ identities
are \leanP{CharacterLambertTransfer}{299}{mod8_character_identities}, and the finite inequalities
are \leanP{CharacterLambertTransfer}{329}{mod8_odd_discrepancy_diff},
\leanP{CharacterLambertTransfer}{381}{mod8_class_discrepancy_diff} and
\leanP{CharacterLambertTransfer}{407}{mod8_asymptotic_equivalence}; the rate uses
\cref{thm:lambertradix}.
\end{proof}

Nothing here concerns $E_{\rm odd}$, $E_{\rm even}$ or the Erd\H{o}s--Borwein constant
$\sum_m1/(2^m-1)$ individually.

\section{Digit frequencies in the kernel}\label{sec:digits}

For a word $w=w_1\cdots w_L$ over $\{0,\dots,b-1\}$ let $v(w)=\sum_iw_ib^{L-i}$ and
$C_w=[v(w)b^{-L},(v(w)+1)b^{-L})$. The word read at shift $n$ (the digits $n+1,\dots,n+L$) equals
$w$ exactly when $x_n\in C_w$. Put $f_w(N)=\frac1N\cnt\{n<N:x_n\in C_w\}$ and
$\delta_N(z)=1-D_N(z;b)$. The fixed points of $t\mapsto bt$ on $\T$ are $c_m=m/(b-1)$,
$0\le m\le b-2$; we call $c_0=0$ the boundary fixed point and the others interior.

\subsection{The dichotomy}\label{subsec:dichotomy}

\begin{lemma}[Fixed points and cylinders]\label{lem:fixedcyl}\evidence{Paper proof}
Let $L\ge1$. If $b\ge3$ and $1\le m\le b-2$, then every $t\in[0,1)$ with
$\dist{t-c_m}<\eta_{m,L}=\min(m,b-1-m)/((b-1)b^L)$ lies in $C_{m^L}$, and
$\eta_{m,L}\ge1/((b-1)b^L)$. Every $t$ with $\dist t<b^{-L}$ reads $0^L$ if $t<\frac12$ and
$(b-1)^L$ if $t>\frac12$.
\end{lemma}

\begin{proof}
$b^Lc_m=mR_L+m/(b-1)$ with $R_L=(b^L-1)/(b-1)$, so $c_m$ lies in $C_{m^L}=[mR_Lb^{-L},(mR_L+1)b^{-L})$
at distances $m/((b-1)b^L)$ and $(b-1-m)/((b-1)b^L)$ from its endpoints. The second claim is
$C_{0^L}=[0,b^{-L})$ and $C_{(b-1)^L}=[1-b^{-L},1)$.
\end{proof}

\begin{lemma}[Localization]\label{lem:localize}\evidence{Paper proof}
If $N\ge3$ and $\delta=\delta_N(z)<\frac12-\frac1N$, then for some fixed point $c$,
$\cnt\{n<N:\dist{x_n-c}>b\delta/(b-1)\}\le\delta N$.
\end{lemma}

\begin{proof}
Fix $\delta<\delta'\le\frac12-\frac1N$. By \cref{lem:arcs} some arc $J$ of length $\ell$ has
$e+\ell<\delta'$, where $e=1-\mu_N(J)$; let $c'$ be its midpoint. At least
$N-1-2eN>N-1-2\delta'N\ge1$ indices $n\le N-2$ have $x_n,x_{n+1}\in J$, and for such $n$
$\dist{(b-1)c'}\le\dist{bc'-bx_n}+\dist{x_{n+1}-c'}\le(b+1)\ell/2$. If $c$ is the fixed point nearest
to $c'$ then $\dist{(b-1)c'}=(b-1)\dist{c'-c}$, so $\dist{c'-c}\le(b+1)\ell/(2(b-1))$ and every
$x_n\in J$ satisfies $\dist{x_n-c}\le b\ell/(b-1)<b\delta'/(b-1)$. Let $\delta'\downarrow\delta$;
one fixed point recurs infinitely often, and the exceptional sets stabilize.
\end{proof}

\begin{theorem}[Finite dichotomy]\label{thm:dichotomy}\evidence{Paper proof}
Let $L\ge1$, $N\ge4$ and $\delta=\delta_N(z)<b^{-L-1}$. Then exactly one of the following holds:
\textup{(i)} the words of length $L$ other than $0^L$ and $(b-1)^L$ have total frequency at most
$\delta$; \textup{(ii)} $b\ge3$ and $f_{m^L}(N)\ge1-\delta$ for one $1\le m\le b-2$.
\end{theorem}

\begin{proof}
\Cref{lem:localize} gives $c$ with at most $\delta N$ shifts farther than
$r=b\delta/(b-1)<b^{-L}/(b-1)$. If $c=0$, \cref{lem:fixedcyl} shows that every other shift reads
$0^L$ or $(b-1)^L$. If $c=c_m$, then $r<\eta_{m,L}$, so every other shift reads $m^L$. The
alternatives are exclusive since $\delta<\frac12$.
\end{proof}

\begin{theorem}[Classification in $K_b$]\label{thm:classification}\evidence{Paper proof}
Let $z\in K_b$.
\begin{enumerate}[label=\textup{(\alph*)},nosep]
\item A single alternative of \cref{thm:dichotomy}, with a single $m$ in case (ii), holds for every
$L$ and all large $N$.
\item In case (ii), $f_{m^L}(N)\to1$ for every $L$.
\item In case (i), for every $L$, $\sum_{w\ne0^L,(b-1)^L}f_w(N)\to0$,
$f_{0^L}(N)-f_0(N)\to0$, $f_{(b-1)^L}(N)-f_{b-1}(N)\to0$ and $f_0(N)+f_{b-1}(N)\to1$. The limit
points of the length-$L$ word distribution are exactly the measures
$\alpha\delta_{0^L}+(1-\alpha)\delta_{(b-1)^L}$ with $\alpha\in[\liminf f_0(N),\limsup f_0(N)]$.
\item For $b\ge3$, $z$ is not simply normal in base $b$. For every $b\ge2$, $z$ is not normal in
base $b$: the frequency of the word $0(b-1)$ tends to $0$.
\end{enumerate}
\end{theorem}

\begin{proof}
Fix $L$ and $N_0\ge5$ with $\delta_N<\min(b^{-L-1},\frac18)$ for $N\ge N_0$. Since
$|f_w(N+1)-f_w(N)|\le1/(N+1)$, an interior alternative at $N$ ($f_{m^L}\ge\frac78$) forces
$f_{m^L}(N+1)\ge\frac7{10}$, which excludes every other alternative at $N+1$; and the boundary
alternative at $N$ forces $f_{m^L}(N+1)\le\frac18+\frac15<\frac78$. So the alternative is constant.
It does not depend on $L$: an interior alternative at level $L$ gives one at level $1$, and the
boundary alternative at level $L$ makes every interior digit an other word. Parts (b) and (c) follow
from $\delta_N\to0$, since a start reading a word beginning with $0$ but different from $0^L$ reads
an other word. The limit points of $f_0(N)$ form an interval because its steps tend to $0$. For
part (d) with $b\ge3$: in case (i) the $b-2$ interior digits have total frequency at most
$\delta_N\to0\ne(b-2)/b$, while in case (ii) $f_m(N)\to1\ne1/b$; either way $z$ is not simply
normal. The word $0(b-1)$ is an other word in case (i) and differs from $m^2$ in case (ii), so
its frequency tends to $0$.
\end{proof}

In base $2$ only case (i) exists, and for $z\in K_2$ simple normality is equivalent to
$f_{0^L}(N)\to\frac12$ for some, equivalently every, $L$: the orbit must approach $0$ from above and
from below equally often.

\subsection{Oriented tails}\label{subsec:oriented}

For $z=F_b(a)$ let $I_n=\sum_{m\le n}a(m)b^{n-m}\in\Z$ and $\tau_n=b^nz-I_n=\sum_{m>n}a(m)b^{n-m}$,
so $x_n=\fr{\tau_n}$. Put $\rho_+(N)=\frac1N\cnt\{n<N:\tau_n\ge0\}$ and
$B_L(N)=\cnt\{n<N:|\tau_n|\ge b^{-L}\}$.

\begin{lemma}[Oriented tails]\label{lem:oriented}\evidence{Paper proof}
\leavevmode
\begin{enumerate}[label=\textup{(\alph*)},nosep]
\item If $0\le\tau_n<b^{-L}$ the word at shift $n$ is $0^L$; if $-b^{-L}<\tau_n<0$ it is $(b-1)^L$.
\item $|Nf_{0^L}(N)-N\rho_+(N)|\le B_L(N)$, $|Nf_{(b-1)^L}(N)-N(1-\rho_+(N))|\le B_L(N)$, and every
other word occurs at most $B_L(N)$ times.
\item If $0=m_0<m_1<\cdots$ contains the support of $a$ in $[1,\infty)$ and $n\in[m_k,m_{k+1})$, then
$\tau_n=b^{n-m_{k+1}}V_{k+1}$ with $V_{k+1}=a(m_{k+1})+\tau_{m_{k+1}}$; so
$N\rho_+(N)$ is the gap-weighted count of nonnegative $V_{k+1}$, and $\sgn V_{k+1}=\sgn a(m_{k+1})$
whenever $|\tau_{m_{k+1}}|<1$.
\item For all $N,R,L\ge1$,
$B_L(N)\le R\,S(N+R-1)+\frac{b^{L-R}}{b-1}\sum_{m>R}|a(m)|b^{-(m-R-N)_+}$. If
$S(X)\ll X(\log X)^{-\kappa}$ and $A(X)\ll X(\log X)^\beta$, then
$B_L(N)=O_L(N\log\log N/(\log N)^\kappa)$.
\end{enumerate}
\end{lemma}

\begin{proof}
(a) The word is $\lfloor b^L\fr{b^nz}\rfloor$ written with $L$ digits, and $\fr{b^nz}$ equals
$\tau_n$ or $1+\tau_n$. (b) follows from (a). (c) There is no support in $(n,m_{k+1})$, and
$|a(m_{k+1})|\ge1$. (d) is the argument of \cref{thm:finitetransfer} applied to $|\tau_n|$ with
window $R$ and Markov threshold $b^{-L}$; for the rate take $R=L+\lceil(\kappa+\beta)\log_b\log N\rceil$.
\end{proof}

Concentration of $\dist{b^nz}$ at $0$ alone does not give $B_L(N)=o(N)$: for $a(m)=-(b-1)$,
$m\ge1$, one has $\tau_n=-1$ for all $n$.

\begin{theorem}[Nonnegative sparse series]\label{thm:nonneg}\evidence{Paper proof}
Let $a\ge0$ be integers with $\sum_ma(m)b^{-m}<\infty$ and $z=F_b(a)$, and let $b\ge2$, $L\ge1$.
\begin{enumerate}[label=\textup{(\alph*)},nosep]
\item If $a$ satisfies one of the three hypotheses of \cref{thm:leantransfer}
\textup{(}$A(X)=O(X)$ and $S(X)=o(X)$; $A(X)\le C(X+1)^k$ and $S(X)\log X=o(X)$;
$A(X)\le CX(\log(X+2))^\gamma$ and $S(X)\log\log X=o(X)$\textup{)}, then $f_{0^L}(N)\to1$.
\item If $S(X)\ll X(\log X)^{-\kappa}$ and $A(X)\ll X(\log X)^\beta$ with $\kappa>0$, $\beta\ge0$,
then $1-f_{0^L}(N)=O_L(\LL/(\log N)^\kappa)$.
\item If the support lies in the squares and $a(m)\le2m^\gamma$ with $\gamma\ge0$, then
$1-f_{0^L}(N)\le R(\sqrt{N+R}+1)/N$ for every
$R\ge L+\lceil\log_b(2C_{b,\gamma}(N+R)^\gamma)\rceil+1$, where
$C_{b,\gamma}=\sum_{k\ge1}(1+k)^\gamma b^{-k}$.
\end{enumerate}
In each case, if $a$ has infinitely many nonzero terms, $z$ is not simply normal in any base,
and in base $2$ its digit $0$ has frequency $1$.
\end{theorem}

\begin{proof}
Since $a\ge0$, $\tau_n\ge0$ for all $n$, so $\rho_+(N)=1$ and \cref{lem:oriented}(b) gives
$1-f_{0^L}(N)\le B_L(N)/N$. In \cref{lem:oriented}(d) the last sum is at most $\Ahat(N+R)$.
(a) Under linear mass, $\Ahat(N+R)\le C(N+R+1)$ by \cref{lem:meantail}; for fixed $R$,
$\limsup_NB_L(N)/N\le Cb^{L-R}/(b-1)$, and $R$ is arbitrary. Under polynomial mass,
$\Ahat(X)\le C_{b,k}'(X+1)^k$ by Abel summation; take
$R=L+\lceil(k-1)^+\log_bN\rceil+\omega_N$ with $\omega_N\to\infty$ and $\omega_N\le\log N$, so
the Markov term is $\ll b^{-\omega_N}N$ and $RS(N+R-1)\ll(\log N)(S(N)+R)=o(N)$. Under
logarithmic mass, $\Ahat(X)\ll X(\log X)^\gamma$; take
$R=L+\lceil\gamma\log_b\log N\rceil+\omega_N$ with $\omega_N\to\infty$ and
$\omega_N\le\log\log N$, so that both terms are $o(N)$. (b) is \cref{lem:oriented}(d). (c) If no
square lies in $(n,n+R]$, then
$\tau_n\le2\sum_{i>R}(N+i)^\gamma b^{-i}\le2(N+R)^\gamma C_{b,\gamma}b^{-R}<b^{-L}$, using
$N+R+k\le(N+R)(1+k)$ and $\gamma\ge0$, and at most $R(\sqrt{N+R}+1)$ shifts have a square in
$(n,n+R]$. The final sentence holds because $f_0(N)\to1$ forces $f_d(N)\to0$ for $d\ne0$.
\end{proof}

Since $(1*\chi)(p^e)\in\{1,e+1,(1+(-1)^e)/2\}$ for real $\chi$, $1*\chi\ge0$, and
\cref{thm:nonneg} applies to $L_\chi(b)$ for every real nonprincipal $\chi$, to every
$\Theta_Q(1/b)$, to $\Theta_3(1/b)$ and to $D^j\Theta_3(1/b)$ $(j\ge1)$, with $D=q\,d/dq$, in every
base. This proves \cref{thm:F}(ii).

\subsection{Signed quadratic families}\label{subsec:signedfamilies}

By \cref{thm:eulerjacobi}, for $j\ge0$,
\[
D^jE_b=\sum_{k\ge1}(-1)^k\bigl((P_k^-)^jb^{-P_k^-}+(P_k^+)^jb^{-P_k^+}\bigr)+[j=0],\qquad
D^jE_b^3=\sum_{k\ge0}(-1)^k(2k+1)T_k^jb^{-T_k},
\]
and $E_b^2DE_b=\frac13D(E_b^3)=\sum_{k\ge1}(-1)^k\frac{k(k+1)(2k+1)}6b^{-T_k}$. In each case every
nonzero coefficient has the sign $(-1)^k$ of its pair.

\begin{proposition}[Lean statements for derivatives]\label{prop:leanderiv}\evidence{Kernel-verified}
If $A(X)\le C(X+1)^k$ and $S(X)\log X=o(X)$, then $D^jF_b(a)\in K_b$ for every $j\ge0$ and every
$b\ge2$. Every $D^j$ of the pentagonal series, of Jacobi's triangular series and of the series of
$\Theta_3(1/b)$ lies in $K_b$; the coefficients of $E_b^2DE_b$ are exactly one third of those of
$D$ applied to Jacobi's series; and the series $E_b^2DE_b$ lies in $K_b$.
\end{proposition}

\begin{proof}
\leanP{DerivativeTransfer}{96}{derivCoeff_tendsto_one},
\leanP{DerivativeTransfer}{196}{deriv_eulerPentagonal_discrepancy_tendsto_one},
\leanP{DerivativeTransfer}{219}{deriv_eulerCube_discrepancy_tendsto_one},
\leanP{DerivativeTransfer}{291}{deriv_thetaThree_discrepancy_tendsto_one},
\leanP{DerivativeTransfer}{350}{three_mul_eulerSqDerivCoeff} and
\leanP{DerivativeTransfer}{373}{eulerSqDeriv_discrepancy_tendsto_one}.
\end{proof}

\begin{theorem}[Frequencies]\label{thm:signedfreq}\evidence{Lean-supported paper proof}
Let $z$ be $D^jE_b$ or $D^jE_b^3$ $(j\ge0)$ or $E_b^2DE_b$. For every $L\ge1$ and $N\ge2$,
\[
\Bigl|f_{0^L}(N)-\tfrac12\Bigr|+\Bigl|f_{(b-1)^L}(N)-\tfrac12\Bigr|+\sum_{w\ne0^L,(b-1)^L}f_w(N)\le C_{b,j}\frac{L+\log N}{\sqrt N}.
\]
Hence these numbers are simply normal in base $2$ and not $2$-normal, and for $b\ge3$ they are not
simply normal in base $b$. For every $k\ge2$, in base $b^k$ the digits $0$ and $b^k-1$ each have
frequency $\frac12$; in particular they are not simply normal in base $4$.
\end{theorem}

\begin{proof}
Let the support be $m_1<m_2<\cdots$ with coefficients $c_i$ and signs $\sigma_i$, and
$R_i=\sum_{l>i}|c_l|b^{m_i-m_l}$. The gaps beyond $m_i$ are at least $K\asymp\sqrt{m_i}$ while the
coefficients are polynomial in $m_i$, so $R_i\le\frac12$ for $i\ge i_0$. For $n\in[m_{i-1},m_i)$ with
$r=m_i-n\ge\ell_i=L+\lceil\log_b(|c_i|+\frac12)\rceil+1$, \cref{lem:oriented}(c) gives
$\tau_n=b^{-r}(c_i+\theta R_i)$ with $|\theta|\le1$, of sign $\sigma_i$ and modulus $<b^{-L}$; so the
word is $0^L$ or $(b-1)^L$ according to $\sigma_i$. The exceptional shifts number
$\sum_{m_{i-1}<N}\ell_i+m_{i_0}=O_{b,j}(\sqrt N(L+\log N))$. It remains to show that the shifts
whose next support point is positive and those whose next support point is negative differ in
number by $O(\sqrt N)$. In the pentagonal case the gaps preceding $P_K^-$ and $P_K^+$ are $2K-1$ and
$K$, both with sign $(-1)^K$, and $|\sum_{K\le K_0}(-1)^K(3K-1)|\le(3K_0+2)/2$; in the triangular case
the gap before $T_k$ is $k$ with sign $(-1)^k$. For base $b^k$, the base-$b^k$ digit at place
$n+1$ is the base-$b$ word of length $k$ at shift $kn$, and the balance argument applies with an extra
$O(1)$ per gap. The Lean support is \cref{prop:leanderiv}, which gives membership in $K_b$.
\end{proof}

\begin{theorem}[Disjunctivity]\label{thm:disjunctive}\evidence{Paper proof}
For every $b\ge2$, $D^jE_b$ $(j\ge1)$, $D^jE_b^3$ $(j\ge0)$ and $E_b^2DE_b$ are disjunctive in base $b$.
\end{theorem}

\begin{proof}
Let $m=m_i$ with $i\ge i_0$, let $\ell$ be the number of digits of $|c_i|$, and suppose the gap before
$m$ is at least $\ell$ and $|\tau_m|<1$; both hold for large $i$. Then
$\lfloor b^mz\rfloor=b^gI_{m_{i-1}}+c_i+\lfloor\tau_m\rfloor$, so the digits at places
$m-\ell+1,\dots,m$ form the $\ell$-digit expansion of $(c_i+\epsilon)\bmod b^\ell$, with $\epsilon=0$
if $\tau_m>0$ and $\epsilon=-1$ if $\tau_m<0$; the sign of $\tau_m$ is $\sigma_{i+1}$. For $D^jE_b$ take
$m=P_K^-$ with $K$ even: then $\sigma_i=\sigma_{i+1}=+1$ and the window reads $v_K=(P_K^-)^j$. For
$D^jE_b^3$ and $E_b^2DE_b$ take $m=T_k$ with $k$ even: then the window reads $v_k=c_k-1$. In each case
$v_K\to\infty$ increases with $v_{K+2}/v_K\to1$. Given a word $w$, let $W$ have digits $1w$; for large
$s$ some $v_{K+2}$ lies in $[Wb^s,(W+1)b^s)$ and so begins with $1w$.
\end{proof}

\Cref{thm:F}(iii) follows from \cref{thm:signedfreq,thm:disjunctive}; recall from
\cref{cor:missing} that $E_b$ itself is not disjunctive.

\subsection{Euler's square}\label{subsec:squaredigits}

\begin{theorem}[Digit law of $E_b^2$]\label{thm:squaredigits}\evidence{Paper proof}
For every $b\ge2$ and $L\ge1$, $f_{0^L}(N;E_b^2)=\rho_+(N)+o(1)$,
$f_{(b-1)^L}(N;E_b^2)=1-\rho_+(N)+o(1)$, and all other words have frequency $o(1)$, where
$\rho_+(N)$ is the gap-weighted proportion of support points $m$ with
$\sum_{m'\ge m}a_2(m')b^{m-m'}\ge0$. Hence $E_b^2$ is not simply normal for $b\ge3$, and $E_2^2$ is simply
normal in base $2$ if and only if $\rho_+(N)\to\frac12$ for $b=2$.
\end{theorem}

\begin{proof}
By \cref{prop:squareelem}, $S(X)\ll X/\sqrt{\log X}$ and $A(X)=O(X)$, so
\cref{lem:oriented}(d) gives $B_L(N)=O(N\LL/\sqrt{\log N})$; apply \cref{lem:oriented}(b),(c).
\end{proof}

The oriented values are in fact never $0$: a vanishing value would make $b^mE_b^2$ an integer, while
$E_b$ is transcendental by Nesterenko's theorem (\cref{prop:arith}). This remark uses that external
input; \cref{thm:squaredigits} does not.

\begin{theorem}[Wirsing's mean-value theorem {\cite{Wirsing1967}}]\label{thm:wirsing67}\evidence{External input}
Let $g$ be multiplicative with $0\le g(n)\le1$ for all $n$, and let $f$ be real multiplicative with
$|f|\le g$. If $\sum_{p\le x}g(p)\log p\sim\tau x$ with $\tau>0$ and
$\sum_p(g(p)-f(p))/p=\infty$, then $\sum_{n\le x}f(n)=o(\sum_{n\le x}g(n))$.
\end{theorem}

\begin{theorem}[Both constant words]\label{thm:squarebothwords}\evidence{Paper proof}
Granting \cref{thm:etainput}, there is $c>0$ such that for every $b\ge2$ and $L\ge1$,
$c\le\liminf_Nf_{0^L}(N;E_b^2)\le\limsup_Nf_{0^L}(N;E_b^2)\le1-c$.
\end{theorem}

\begin{proof}
Let $U=\supp a_2$ and $s=\sgn a_2$ on $U$. The function $f(m)=\sgn c(m)$ is real multiplicative with
$|f|\le g$, the indicator of $\{12n+1:n\in U\}$, and $\sum_{p\le x}g(p)\log p\sim x/4$. By
\cref{thm:square}(c) the primes with $\epsilon_p=-1$ have Chebotarev density $\frac18$ in the ring
class field, so $\sum_p(g(p)-f(p))/p=\infty$, and \cref{thm:wirsing67} gives
$\sum_{n\le N,\,n\in U}s(n)=o(S_\eta(N))$. Hence $U\cap(N/2,N]$ contains at least
$\kappa'N/(\log N)^{3/4}$ positive points for large $N$. Let $H=\lfloor\eta(\log N)^{3/4}\rfloor$. By
\cref{thm:S}, at most $2C_0HN/(\log N)^{3/2}\le\frac18\kappa'N/(\log N)^{3/4}$ points of $U\cap(N/2,N]$
have another support point within $H$ when $\eta\le\kappa'/(16C_0)$. By Markov's inequality,
$o(N/(\log N)^{3/4})$ points have $|a_2(n)|>(\log N)^{1/3}$, and $O(Nb^{-H})$ isolated points have
$|\tau_m|\ge1$. Every remaining positive point $m$ has a positive oriented value
$a_2(m)+\tau_m$ and produces the word $0^L$ at the
shifts in $[m-H,m-r_0]$, $r_0=L+\log_b(2(\log N)^{1/3})+1$, by \cref{lem:oriented}; these sets are
disjoint, so $Z_{0^L}(N)\ge(H-r_0)\frac14\kappa'N/(\log N)^{3/4}\ge\frac15\eta\kappa'N$. The same
argument with $s=-1$ bounds $Z_{(b-1)^L}(N)$ from below, and $Z_{0^L}+Z_{(b-1)^L}\le N$.
\end{proof}

At $N=2^{18}$ the certified binary digits of $E_2^2$ give $f_0=0.50170$ against $\rho_+=0.50179$; this
is a finite measurement, and the simple normality of $E_2^2$ in base $2$ remains open.

\subsection{Rational scalings}\label{subsec:automaton}

Let $d\ge2$ with $\gcd(d,b)=1$, $a\ne0$ with $\gcd(a,d)=1$, and $t=\ord_d(b)$. Let
$z\in\{E_b,E_b^3\}$ with insertion positions $p_j$ and coefficients $c_j$. Then
$b^n\frac adz=\frac adI_n+\frac ad\tau_n$, and the \emph{atom} at shift $n$ is $\alpha_n=aI_n\bmod d$.
Between insertions $I_{n+1}=bI_n$, and at an insertion $I_p=bI_{p-1}+c_p$.

\begin{lemma}[Nonzero atoms are interior]\label{lem:atoms}\evidence{Paper proof}
For $1\le c\le d-1$ and $L\ge1$, $c/d$ is at distance at least $1/(db^L)$ from $b^{-L}\Z$; so $c/d$ lies
in the interior of the cylinder labelled by the first $L$ digits $\pi_L(c/d)$ of its purely periodic
expansion.
\end{lemma}

\begin{proof}
$|c/d-k/b^L|=|cb^L-kd|/(db^L)$ and $d\nmid cb^L$.
\end{proof}

\begin{theorem}[Signed automaton]\label{thm:automaton}\evidence{Lean-supported paper proof}
Augment each run between insertions with the sign $\sigma$ of the insertion that ends it. The pairs
(atom at the start of the run, $\sigma$) evolve by a purely periodic finite automaton, and every word
$w$ of every length $L$ has a rational limiting frequency
\[
f_w=\sum_{(u,\sigma)}\omega(u,\sigma)\frac1{|O(u)|}\sum_{x\in O(u)}\ind\bigl[\mathrm{word}(ax\bmod d,\sigma\sgn a)=w\bigr],
\]
where $O(u)$ is the orbit of $u$ under multiplication by $b$, $\omega(u,\sigma)$ is the limiting share
of runs of that class, $\mathrm{word}(c,\cdot)=\pi_L(c/d)$ for $c\ne0$, $\mathrm{word}(0,+)=0^L$ and
$\mathrm{word}(0,-)=(b-1)^L$. Moreover $\max_{|w|=L}|f_w(N)-f_w|$ is $O(N^{-1/2})$ for $E_b$ and
$O(\log N/\sqrt N)$ for $E_b^3$.
\end{theorem}

\begin{proof}
For $E_b$ the state after the pair $k$ is $(k\bmod\mathrm{lcm}(2,t),\,I_{P_k^+}\bmod d)$, updated by an
affine bijection of $\Z/d$; pure periodicity is
\leanP{NormalReciprocalCore}{336}{eulerAutomaton_purely_periodic}, via
\leanP{NormalReciprocalCore}{315}{eulerAutomatonStep_bijective}; the cube is analogous with states
modulo $\mathrm{lcm}(t,2,d)$. Let $n$ lie in the run ending at the insertion $p$ with coefficient $c$,
at distance $r=p-n\ge\ell=L+\lceil\log_b(|a|(|c|+\frac12))\rceil+1$, beyond the finitely many
insertions whose remainder exceeds $\frac12$. Then $a\tau_n$ has sign $\sigma\sgn a$ and modulus
$<b^{-L}$. If $\alpha_n\ne0$, then $|x_n-\alpha_n/d|<1/(db^L)$ and \cref{lem:atoms} gives the word
$\pi_L(\alpha_n/d)$; if $\alpha_n=0$, \cref{lem:oriented}(a) gives $0^L$ or $(b-1)^L$ by the sign. The
exceptional shifts number $O(\sqrt N(L+\log(|a|N)))$. Inside a run $\alpha_n$ moves along $u\mapsto ub$
while $\sigma$ is constant, and runs of one type have lengths that are the same linear function of
$k$ for every class, so each class receives its share up to $O(1/\sqrt N)$.
\end{proof}

\begin{theorem}[Odd order]\label{thm:oddorder}\evidence{Paper proof}
If $t=\ord_d(b)$ is odd, then the limit word law of $\frac adE_b$, and of $\frac adE_b^3$ when moreover
$d$ is odd, is invariant under the digit complement $w\mapsto\bar w$, $\bar w_i=b-1-w_i$. In particular,
in base $2$, $\frac adE_2$ and $\frac adE_2^3$ are simply normal whenever $\ord_d(2)$ is odd.
\end{theorem}

\begin{proof}
For $E_b$ write $P(j)=j(3j-1)/2$ for $j\in\Z$, so that the generalized pentagonal numbers up to $P_k^+$
are $P(j)$, $-k\le j\le k$, and $I_{P_k^+}\equiv b^{P_k^+}\sum_{j=-k}^k\mathrm{tm}(j)$ with
$\mathrm{tm}(j)=(-1)^jb^{-P(j)}\bmod d$. Since $t$ is odd, $P(j+t)-P(j)=t(6j+3t-1)/2\equiv0\pmod t$,
so $\mathrm{tm}(j+t)=-\mathrm{tm}(j)$, any $2t$ consecutive terms sum to $0$, and the atom of the run
of index $k+t$ is the negative of that of index $k$, while $b^{P_{k+t}^\pm}=b^{P_k^\pm}$ and the run
signs flip. Negation sends $\pi_L(c/d)$ to its complement and swaps $0^L$ with $(b-1)^L$, and classes
$k$ and $k+t$ have equal weight. For $E_b^3$ let $T=\mathrm{lcm}(t,d)$, odd, and
$\mathrm{tm}(i)=(-1)^i(2i+1)b^{-T_i}\bmod d$. Then $\mathrm{tm}(i+T)=-\mathrm{tm}(i)$, and the
involution $i\mapsto T-1-i$ gives $\mathrm{tm}(T-1-i)=-\mathrm{tm}(i)$, so $\sum_{i<T}\mathrm{tm}(i)=0$
(as $d$ is odd) and the argument concludes as before.
\end{proof}

\begin{proposition}[No higher normality]\label{prop:nohigher}\evidence{Paper proof}
At most $d+1$ words of length $L$ have positive limit frequency in $\frac adE_b$ or $\frac adE_b^3$.
Hence these numbers are not $L$-normal for $b^L>d+1$, and never normal.
\end{proposition}

\begin{proof}
By \cref{thm:automaton} the only words with positive limit frequency are $\pi_L(c/d)$, $c\ne0$, and
$0^L$, $(b-1)^L$.
\end{proof}

For example, $\frac17E_2$ and $\frac17E_2^3$ are simply normal in base $2$ (as $\ord_7(2)=3$), with
$f_{00}=f_{11}=\frac{11}{54}$ for $\frac17E_2$, whereas $\frac13E_2$ has $f_0=\frac7{12}$ and
$\frac12E_3$ has digit frequencies $\frac16,\frac23,\frac16$ in base $3$. In every case examined with
$d<64$, simple normality of $\frac adE_2$ in base $2$ holds exactly when $\ord_d(2)$ is odd; the
converse of \cref{thm:oddorder} is open (\cref{sec:open}).

\section{A uniform pair sieve for Euler's square}\label{sec:sieve}

Let $\mathcal P_0$ be the set of primes $p\ge5$ with $p\not\equiv1\pmod{12}$. By \cref{thm:square}(a),
if $a_2(n)\ne0$ then $v_p(12n+1)$ is even for every $p\in\mathcal P_0$; for such $n$ put
$v(n)=\prod_{p\in\mathcal P_0}p^{v_p(12n+1)/2}$ and $u(n)=(12n+1)/v(n)^2$, a product of primes
$\equiv1\pmod{12}$.

\begin{theorem}[Uniform pair sieve]\label{thm:S}\evidence{Paper proof}
Granting \cref{thm:etainput}, there is an absolute constant $C_0$ such that for all $N\ge2$ and
$1\le H\le N$,
\[
\sum_{h=1}^H\cnt\{1\le n\le N:\ a_2(n)a_2(n+h)\ne0\}\le C_0\frac{HN}{(\log N)^{3/2}}.
\]
\end{theorem}

\begin{lemma}[Selberg's upper bound]\label{lem:selberg}\evidence{Paper proof}
Let $F\in\Z[m]$, let $I$ be a set of $X$ consecutive integers, $\mathcal Q$ a finite set of primes,
$P=\prod_{p\in\mathcal Q}p$, $\rho(d)=\cnt\{m\bmod d:d\mid F(m)\}$ for squarefree $d\mid P$, and assume
$1\le\rho(p)<p$ on $\mathcal Q$. Put $g(d)=\prod_{p\mid d}\rho(p)/(p-\rho(p))$,
$G(\xi)=\sum_{d<\xi,\,d\mid P}g(d)$ and $r_d=\cnt\{m\in I:d\mid F(m)\}-X\rho(d)/d$. Then $|r_d|\le\rho(d)$
and for every $\xi>1$
\[
\cnt\{m\in I:\gcd(F(m),P)=1\}\le\frac X{G(\xi)}+\sum_{d_1,d_2<\xi,\ d_i\mid P}|r_{[d_1,d_2]}|.
\]
\end{lemma}

\begin{proof}
$\rho$ is multiplicative by the Chinese remainder theorem and each of the $\rho(d)$ classes meets $I$
in $\lfloor X/d\rfloor$ or $\lceil X/d\rceil$ points. For real $\lambda_d$ ($d\mid P$, $d<\xi$) with
$\lambda_1=1$, $\ind[\gcd(F(m),P)=1]\le(\sum_{d\mid\gcd(F(m),P)}\lambda_d)^2$, so the count is at most
$XQ(\lambda)+\sum\lambda_{d_1}\lambda_{d_2}r_{[d_1,d_2]}$ with
$Q(\lambda)=\sum\lambda_{d_1}\lambda_{d_2}f([d_1,d_2])$, $f(d)=\rho(d)/d$. With $h=1/f-1$ extended
multiplicatively and $y_t=\sum_{t\mid d}\lambda_df(d)$, $Q(\lambda)=\sum_th(t)y_t^2$ and
$1=\sum_t\mu(t)y_t$, so Cauchy--Schwarz gives $Q\ge1/G(\xi)$ with equality for
$y_t=\mu(t)g(t)/G(\xi)$. For this choice
$\lambda_d=\mu(d)\frac{g(d)}{f(d)}\frac{G_d(\xi/d)}{G(\xi)}$ with $G_d(y)=\sum_{k<y,\,k\mid P,\,(k,d)=1}g(k)$;
since $g(d)/f(d)=\sum_{s\mid d}g(s)$ and the products $sk$ are distinct divisors of $P$ below $\xi$,
$|\lambda_d|\le1$.
\end{proof}

\begin{lemma}[Lower bound for $G$]\label{lem:Glower}\evidence{Paper proof}
In \cref{lem:selberg} assume also $\rho(p)\in\{1,2\}$ and $5\le p\le z$ on $\mathcal Q$, $\log z\ge2$
and $\xi\ge z^{17}$. Then $G(\xi)\ge\frac12\prod_{p\in\mathcal Q}(1-1/p)^{-\rho(p)}$.
\end{lemma}

\begin{proof}
Let $\kappa$ be multiplicative with $\kappa(p^k)=\binom{k+\rho(p)-1}{\rho(p)-1}\le\rho(p)^k$. Then
$g(e)\ge\sum_{\rad n=e}\kappa(n)/n$ for squarefree $e\mid P$, so $G(\xi)\ge\Pi-T$ with
$\Pi=\prod_{p\in\mathcal Q}(1-1/p)^{-\rho(p)}$ and $T=\sum_{n\ge\xi,\,\rad n\mid P}\kappa(n)/n$. With
$\delta=1/\log z$, Rankin's trick gives $T\le\xi^{-\delta}\prod_p(1-p^{\delta-1})^{-\rho(p)}$; since
$\delta\log p\le1$, $p^\delta-1\le(e-1)\delta\log p$ and $1-p^{\delta-1}\ge1-e/5$, and Mertens'
bound $\sum_{p\le z}(\log p)/p\le\log z+2$ gives $T\le\xi^{-\delta}\Pi e^{15.07}\le e^{-1.93}\Pi<\Pi/2$.
\end{proof}

\begin{lemma}[Mertens over $\mathcal P_0$]\label{lem:mertensP0}\evidence{Paper proof}
There is an absolute $B$ with $\prod_{p\le z,\,p\in\mathcal P_0}(1-1/p)^2\le e^{2B}(\log z)^{-3/2}$ for
$z\ge5$.
\end{lemma}

\begin{proof}
For $a\in\{5,7,11\}$, $\sum_{p\le x,\,p\equiv a\,(12)}(\log p)/p=\frac14\log x+O(1)$
\cite[Thm.~7.3]{Apostol}; partial summation and $(1-1/p)^2\le e^{-2/p}$.
\end{proof}

\begin{lemma}[Supported pairs]\label{lem:pairstructure}\evidence{Paper proof}
Let $h\ge1$, $a_2(n)a_2(n+h)\ne0$, $v=v(n)$, $v'=v(n+h)$ and $g=\gcd(v,v')$. Then $g^2\mid h$. For
$v,v'$ prime to $6$ with $g^2\mid h$, the $n$ with $v^2\mid12n+1$ and $v'^2\mid12(n+h)+1$ form one class
$n\equiv n_0\pmod L$, $L=v^2v'^2/g^2$; writing $n=n_0+mL$,
$(12n+1)/v^2=A_1m+B_1$ and $(12(n+h)+1)/v'^2=A_2m+B_2$ with $A_1=12v'^2/g^2$, $A_2=12v^2/g^2$ and
$A_1B_2-A_2B_1=144h/g^2$. For $p\in\mathcal P_0$ with $p\nmid vv'$, the product
$(A_1m+B_1)(A_2m+B_2)$ has $2$ roots modulo $p$ if $p\nmid h$ and $1$ root if $p\mid h$.
\end{lemma}

\begin{proof}
$g^2$ divides $12n+1$ and $12n+12h+1$, hence $12h$, and $\gcd(g,6)=1$. The congruences are solvable
exactly when they agree modulo $g^2$, that is, $g^2\mid h$, and then the solution is unique modulo
$L$. The relation $v^2(A_1m+B_1)+12h=v'^2(A_2m+B_2)$ gives the determinant, and two roots modulo
$p\nmid12vv'$ coincide exactly when $p\mid144h/g^2$, that is, $p\mid h$.
\end{proof}

\begin{proof}[Proof of \cref{thm:S}]
We first show that for $N\ge N_1=\lceil e^{200}\rceil$ and $1\le h\le N$,
\begin{equation}\label{eq:pershift}
P_h(N):=\cnt\{1\le n\le N:a_2(n)a_2(n+h)\ne0\}\le C_1\frac h{\varphi(h)}\frac N{(\log N)^{3/2}},
\end{equation}
with $C_1=4000e^{2B}C_5+1$ and
$C_5=\sum_g(g/\varphi(g))^4g^{-2}\bigl(\sum_w(w/\varphi(w))^2w^{-2}\bigr)^2<\infty$. Put
$V=(\log N)^3$, $z=N^{1/100}$, $\xi=z^{17}$. Pairs with $v(n)>V$ or $v(n+h)>V$ number at most
$2N/V+\sqrt{12N+1}+\sqrt{24N+1}$, since each $v$ prime to $6$ fixes one class modulo $v^2$ and
$\sum_{v>V,\,(v,6)=1}v^{-2}\le1/V$. For the other pairs, fix $(v,v')$ with $v,v'\le V$ and $g^2\mid h$.
By \cref{lem:pairstructure} the relevant $m$ form a set of $X\le2N/L$ consecutive integers at which
$F(m)=(A_1m+B_1)(A_2m+B_2)$ has no prime factor in $\mathcal Q=\{p\in\mathcal P_0:p\le z,\ p\nmid vv'\}$,
and $\rho(p)\in\{1,2\}$ there. \Cref{lem:selberg,lem:Glower,lem:mertensP0}, together with
$|r_e|\le\tau(e)$ and $\sum_{d<\xi}\tau(d)\le\xi(1+\log\xi)$, bound this count by
\[
4000e^{2B}\frac h{\varphi(h)}\frac N{(\log N)^{3/2}}\cdot\frac{(v/\varphi(v))^2(v'/\varphi(v'))^2}L+\xi^2(1+\log\xi)^2 .
\]
Writing $v=gw$, $v'=gw'$ with $\gcd(w,w')=1$, $L=g^2w^2w'^2$, and the sum of the last factor over all
pairs is at most $C_5$, uniformly in $h$. The remainders total at most
$V^2\xi^2(1+\log\xi)^2+2N/V+9\sqrt N\le N(\log N)^{-3/2}$ for $\log N\ge200$. This proves
\eqref{eq:pershift}. Since $h/\varphi(h)=\sum_{d\mid h}\mu(d)^2/\varphi(d)$,
$\sum_{h\le H}h/\varphi(h)\le H\zeta(2)\zeta(3)/\zeta(6)\le1.9436H$. For $N<N_1$ use the trivial bound
$HN\le(\log N_1)^{3/2}HN/(\log N)^{3/2}$.
\end{proof}

\begin{theorem}[Rational scalings of $E_b^2$]\label{thm:squarescaling}\evidence{Paper proof}
Granting \cref{thm:etainput}, for a reduced fraction $a/d$, $\frac adE_b^2\in K_b$ if and only if
$a/d\in\Z[1/b]$.
\end{theorem}

\begin{proof}
As in the proof of \cref{thm:polyscal} it suffices to show $\frac1\ell E_b^2\notin K_b$ for each prime
$\ell\nmid b$; $\ell=2$ is \cref{thm:polyscal}(c). Let $\ell$ be odd and suppose
$\frac1\ell E_b^2\in K_b$. With $I_n$ and $\tau_n$ as in \cref{subsec:oriented}, the transfer bounds
give $|\tau_n|\to0$ in density, so $I_n\equiv c\pmod\ell$ for a density-one set of $n$, for some $c$
with $bc\equiv c$. Let $S_\ell=\{n:\ell\nmid a_2(n)\}$. Between consecutive elements of $S_\ell$,
$I_n\equiv bI_{n-1}\pmod\ell$, so on each such run $I_n\equiv c$ everywhere or nowhere, and two
consecutive runs cannot both sit at $c$. Let $S'$ be the set of $n$ with $12n+1=uv^2$, $u$ squarefree
composed of primes $\equiv1\pmod{12}$ and $v$ composed of primes in $\mathcal P_0$; by
\cref{thm:square}(a), $|a_2(n)|=2^{\omega(u)}$ on $S'$, so $S'\subseteq S_\ell$, and Wirsing's theorem
gives $|S'\cap(N/2,N]|\ge\kappa N(\log N)^{-3/4}$. By \cref{thm:S} applied to $n$ and to $n-h$, with
$H=\lfloor\kappa(\log N)^{3/4}/(4C_0)\rfloor$, at least $\frac12\kappa N(\log N)^{-3/4}$ points of
$S_\ell\cap(N/2,N]$ have both neighbouring gaps at least $H$, and at each of them one of the two
adjacent stretches of length $H$ avoids $c$. Each $n$ lies in at most two such stretches, so
$\cnt\{n<N+H:I_n\not\equiv c\}\ge\frac H2\cdot\frac{\kappa N}{2(\log N)^{3/4}}\ge\frac{\kappa^2}{20C_0}N$,
a contradiction.
\end{proof}

\Cref{thm:C}(iii) follows from \cref{thm:polyscal,thm:squarescaling}. The limit law of
$\frac adE_b^2$ for odd $d>1$ is not determined here.

\section{Limits of sparse and signed transfer}\label{sec:limits}

\subsection{Products and dense series}\label{subsec:products}

\begin{proposition}[Sumset criterion]\label{prop:sumset}\evidence{Lean-supported paper proof}
Let $x=F_b(a)$ and $y=F_b(c)$ with absolutely convergent integer series, and
$\bar A_a(X)=\sum_{0\le m\le X}|a(m)|$. If
$\cnt\bigl((\supp a\cap[0,X])+(\supp c\cap[0,X])\bigr)=o(X)$ and $\bar A_a(X)\bar A_c(X)=O(X)$, then
$xy\in K_b$.
\end{proposition}

\begin{proof}
$xy=F_b(a*c)$, $\supp(a*c)\subseteq\supp a+\supp c$, and the convolution mass bound of
\cref{prop:leaneuler} gives $\sum_{m\le X}|(a*c)(m)|\le\bar A_a(X)\bar A_c(X)$; apply
\cref{thm:leantransfer}.
\end{proof}

\begin{proposition}[Squares leave the kernel]\label{prop:squaring}\evidence{Paper proof}
Let $\mathcal B$ be the set of integers whose base-$4$ digits lie in $\{0,1\}$, $r\ge2$,
$U=\sum_{\beta\in\mathcal B}b^{-r\beta}$, $V=\sum_{\beta\in\mathcal B}b^{-2r\beta}$ and $x=U+V$. Then
$x\in K_b$ and $x^2\notin K_b$.
\end{proposition}

\begin{proof}
Every $n\ge0$ is uniquely $\beta+2\beta'$ with $\beta,\beta'\in\mathcal B$ (Moser--de Bruijn), so
$UV=\sum_nb^{-rn}=b^r/(b^r-1)$. Both $U$ and $V$ have bounded coefficients and support
$\ll X^{1/2}$, so $x\in K_b$. Sums of two elements of $\mathcal B$ have base-$4$ digits in
$\{0,1,2\}$, so $\mathcal B+\mathcal B$ has counting function $\asymp X^{\log3/\log4}=o(X)$, and
\cref{prop:sumset} gives $U^2,V^2\in K_b$. Hence $x^2\equiv2b^r/(b^r-1)\pmod{K_b}$. The reduced
denominator $(b^r-1)/\gcd(2,b^r-1)$ is prime to $b$ and exceeds $b-1$, so this rational is not in
$K_b$ \cite{DiscrepancyCompanion}, and neither is $x^2$.
\end{proof}

\begin{proposition}[Theta powers modulo the kernel]\label{prop:thetapowers}\evidence{Paper proof}
In every base $b\ge2$, with $\Theta=\Theta_3(1/b)$ and $L=L_{\chi_{-4}}(b)$,
$\Theta^4-16L^2\in K_b$ and $\Theta^3-4\Theta L\in K_b$. In base $2$,
$|D_N(\Theta_3(1/2)^4)-D_N(L_{\chi_{-4}}(2)^2)|\le1-D_N(L_{\chi_{-4}}(2))+7/N\ll\LL/\sqrt{\log N}$.
\end{proposition}

\begin{proof}
$\Theta^2=1+4L$, so $\Theta^4=1+8L+16L^2$ and $\Theta^3=\Theta+4\Theta L$, while $L$ and $\Theta$ lie in
$K_b$ (\cref{thm:thetaupper}). Apply \cref{lem:linear}. For the base-$2$ bound, write
$2^4L^2-\Theta^4=-1-z$ with $z=2^3L$ and apply \cref{lem:linear} with $k=4$ and the single term
$z$ (coefficient $-1$; the integer $-1$ does not change $D_N$); this gives $4/N+(1-D_N(z))$, and
\cref{lem:linear} again gives $1-D_N(z)\le1-D_N(L)+3/N$.
\end{proof}

These transfers decide nothing: membership in $K_b$, normality and $D^+<1$ are open for
$\Theta_3(1/b)^3$, $\Theta_3(1/b)^4$, $E_2(1/b)$, $E_4(1/b)$, $E_6(1/b)$, $\Delta(1/b)=b^{-1}E_b^{24}$ and
the partition constant $P_b=1/E_b=\sum p(n)b^{-n}$. For dense coefficient series the exact identity
$\dist{x_n-R_{n,L}b^{-L}}\le\sum_{j>L}|a(n+j)|b^{-j}$, with
$R_{n,L}=\sum_{j\le L}a(n+j)b^{L-j}\bmod b^L$, reduces $D_N$ to finite residues; for $P_b$ the required
window is $(\pi\sqrt{2/3}/\log b)\sqrt N+O(\log N)$ digits, sharp to first order by the
Hardy--Ramanujan asymptotic \cite{HardyRamanujan1918}. Congruences such as Ramanujan's
$p(5k+4)\equiv0\pmod5$ leave no trace in the digits:

\begin{proposition}[Wash-out]\label{prop:washout}\evidence{Paper proof}
Let $Z\subset\Z_{\ge1}$ contain no two consecutive integers. Every real $u$ has a representation
$u=\lfloor u\rfloor+\sum_{m\ge1}a^*(m)b^{-m}$ with integers $0\le a^*(m)\le b^2-1$ and $a^*(m)=0$ for
$m\in Z$.
\end{proposition}

\begin{proof}
Let $d_1d_2\cdots$ be the base-$b$ digits of $\fr u$. Move each digit $d_m$, $m\in Z$, to position
$m+1\notin Z$ via $d_mb^{-m}=(bd_m)b^{-m-1}$; the new coefficient at $m+1$ is
$d_{m+1}+bd_m\le b^2-1$.
\end{proof}

\subsection{Signed digits: completeness}\label{subsec:completeness}

A \emph{local decomposition} of $a$ assigns to each $m$ a finite multiset $\mathcal D_m$ of pairs
$(u,e)$ with $u\in\Z\setminus\{0\}$, $e\ge0$, $\sum_{\mathcal D_m}ub^e=a(m)$ and
$\sum_m\sum_{\mathcal D_m}|u|b^{e-m}<\infty$. Its net sequence is
$A_{\rm net}(t)=\sum_m\sum_{(u,e)\in\mathcal D_m,\ m-e=t}u$.

\begin{theorem}[Net transfer]\label{thm:net}\evidence{Lean-supported paper proof}
If $m-e\ge-E_0$ for all pairs, then $F_b(a)-\sum_{t\ge1}A_{\rm net}(t)b^{-t}\in\Z$. Consequently
$F_b(a)\in K_b$ whenever $A'(t)=A_{\rm net}(t)-c$ satisfies, for some $c\in\Z$, the hypotheses of
any of the three cases of \cref{thm:leantransfer}.
\end{theorem}

\begin{proof}
Absolute convergence permits the rearrangement and the terms with $t\le0$ are integers. Then
$F_b(a)\equiv c/(b-1)+F_b(A')$, with $c/(b-1)$ a fixed point, and $K_b$ is invariant under such
translations \cite{DiscrepancyCompanion}; apply the Lean theorems listed in
\cref{thm:leantransfer} to $A'$.
\end{proof}

\begin{theorem}[Completeness]\label{thm:completeness}\evidence{Paper proof}
$x\in K_b$ if and only if there are $c\in(b-1)^{-1}\Z$ and integers $A(t)$ $(t\ge1)$ with
$|A(t)|\le\lfloor b/2\rfloor$, $\cnt\{t\le N:A(t)\ne0\}=o(N)$ and
$x-c-\sum_{t\ge1}A(t)b^{-t}\in\Z$. One may take $A(t)=I_t-bI_{t-1}$ with
$I_t=\lfloor b^t(x-c)+\frac12\rfloor$. In base $2$, $x\in K_2$ if and only if the binary digits of $x$
change $o(N)$ times among the first $N$.
\end{theorem}

\begin{proof}
``If'' is \cref{thm:net}. Conversely $x=c+w$ with $M_N(w)\to0$ \cite{DiscrepancyCompanion}. With
$\varepsilon_t=I_t-b^tw\in(-\frac12,\frac12]$, $A(t)=\varepsilon_t-b\varepsilon_{t-1}$ has modulus
$<(b+1)/2$, so $|A(t)|\le\lfloor b/2\rfloor$; if $\dist{b^{t-1}w}<1/(2b)$ then $I_t=bI_{t-1}$, so by
Markov's inequality $\cnt\{t\le N:A(t)\ne0\}\le2bNM_N(w)=o(N)$. In base $2$ the canonical sequence is
$A(n)=d_{n+1}-d_n$.
\end{proof}

So every element of $K_b$ has a sparse bounded signed representation computed from its own digits,
and signed transfer with arbitrary cancellation between coefficients is equivalent to membership.
A barrier can only concern decompositions fixed in advance, coefficient by coefficient.

\subsection{The entropy--multiplicity wall}\label{subsec:wall}

For an integer $x\ne0$ let $w_b^\pm(x)$ be the least $r$ such that $x=\sum_{i\le r}u_ib^{e_i}$ with
distinct $e_i\ge0$ and $1\le|u_i|\le b-1$. (For $b=2$ this is the weight of the non-adjacent form
\cite{Reitwiesner1960}.)

\begin{lemma}[Counting signed representations]\label{lem:entropy}\evidence{Paper proof}
\leavevmode
\begin{enumerate}[label=\textup{(\alph*)},nosep]
\item If $x=\sum_{i\le r}u_ib^{e_i}$ with $e_1>\dots>e_r\ge0$ and $1\le|u_i|\le b-1$, then
$|x|\ge b^{e_1-r+1}$.
\item If $x=\sum_ec_eb^e$ is finite, then $w_b^\pm(x)\le\sum_e|c_e|$; so a decomposition of $x$ into
$r$ terms $ub^e$ with $|u|\le B$ has $r\ge w_b^\pm(x)/B$.
\item For $X\ge b$ and $k\ge0$, $\cnt\{0<|x|\le X:w_b^\pm(x)\le k\}\le\sum_{r\le k}\binom Er(2b-2)^r$
with $E=\lfloor\log_bX\rfloor+k$. If $N\ge b^2$, $A\ge1$, $\theta\in[2^{-20},1]$, $X=N^A$ and
$k\le\frac\theta{40A}\log_bN$, the count is at most $N^\theta$.
\end{enumerate}
\end{lemma}

\begin{proof}
(a) $|x|\ge b^{e_1}-(b-1)(b^{e_1-1}+\dots+b^{e_1-r+1})=b^{e_1-r+1}$. (b) Normalize carries upward:
replacing $c_e=qb+s$ ($|s|<b$, $s$ of the sign of $c_e$) by $s$ and adding $q$ to $c_{e+1}$ lowers
$\sum|c_e|$ by $(b-1)|q|$. (c) By (a), exponents are at most $\log_bX+k-1$; choose $r$ positions
and signed digits. With $\ell=\log_bN$ and $c=k/\ell$, $\sum_{r\le k}\binom Er\le2^{(A+c)\ell H(c/(A+c))}$
with $H(p)\le p\log_2(e/p)$, and the resulting exponent is at most $\theta\ell\log_2b$ at
$c=\theta/(40A)$.
\end{proof}

\begin{theorem}[Entropy--multiplicity wall]\label{thm:wall}\evidence{Paper proof}
Let $b\ge2$, $A\ge1$, $\theta\in[2^{-20},1]$, $c=\theta/(40A)$, $N\ge b^2$, and $\mathcal T\subseteq[1,N]$
with $0<|a(n)|\le N^A$ on $\mathcal T$. Let $\mu=\max_y\cnt\{n\in\mathcal T:a(n)=y\}$. Then
$\sum_{n\in\mathcal T}w_b^\pm(a(n))\ge c\log_bN\,(|\mathcal T|-\mu N^\theta)$, and every local
decomposition with $|u|\le B$ uses at least $(c/B)\log_bN\,(|\mathcal T|-\mu N^\theta)$ pairs on
$\mathcal T$.
\end{theorem}

\begin{proof}
Each value $y$ with $w^\pm(y)\le c\log_bN$ is taken at most $\mu$ times, and there are at most $N^\theta$
such values by \cref{lem:entropy}(c). The last claim is \cref{lem:entropy}(b).
\end{proof}

\begin{theorem}[CM forms of weight $2$, $3$, $4$ {\cite{Ribet1977,Serre1981,Serre1985}}]\label{thm:cminput}\evidence{External input}
Let $g$ be one of $\eta(6\tau)^4$, $\eta(4\tau)^6$, $\eta(3\tau)^8$, of weight $k\in\{2,3,4\}$. Then
$g$ is a newform with complex multiplication by $F=\Q(\sqrt{-3})$, $\Q(i)$, $\Q(\sqrt{-3})$
respectively: for split $p\nmid$ level, $p=\alpha\bar\alpha$ in $\mathcal O_F$ and
$a(p)=2\Re(\chi(\alpha)\alpha^{k-1})$ with $\chi$ taking values in a finite set of roots of unity;
$a(p)=0$ for inert $p$; $|a(n)|\le\tau(n)n^{(k-1)/2}$; and the support of $a$ has counting function
$\asymp X/\sqrt{\log X}$.
\end{theorem}

\begin{corollary}[Signed width of modular coefficients]\label{cor:wall}\evidence{Paper proof}
In every radix $b\ge2$, with $W_\pm(N)=\sum_{n\le N,\,a(n)\ne0}w_b^\pm(a(n))$:
\textup{(a)} for $D(\Theta_3(1/b)^2)=\sum_nnr_2(n)b^{-n}$, $W_\pm(N)\asymp N\sqrt{\log N}$;
\textup{(b)} for $D(E_b^2)$, granting \cref{thm:etainput}, $W_\pm(N)\asymp N(\log N)^{1/4}$;
\textup{(c)} for $E_b^4$, $E_b^6$, $E_b^8$, $W_\pm(N)\asymp N\sqrt{\log N}$. The same lower bounds,
divided by $B$, hold for the number of pairs of every local decomposition with $|u|\le B$.
\end{corollary}

\begin{proof}
Upper bounds: $W_\pm\le\sum_{a(n)\ne0}(1+\log_b|a(n)|)$ with the support bounds and
$|a(n)|\le N^{O(1)}$. (a) Take $\mathcal T=\supp r_2\cap[1,N]$, $A=2$, $\theta=\frac14$; if
$nr_2(n)=y$ then $r_2(n)$ is $4$ times a divisor count, so $\mu\le\max_{n\le N}\tau(n)=N^{o(1)}$, and
\cref{thm:wall} gives $W_\pm\gg\log_bN\cdot N/\sqrt{\log N}$. (b) Take $\mathcal T=\supp a_2\cap[1,N]$,
of size $\gg N(\log N)^{-3/4}$; $n|a_2(n)|\le N^2$, and $\mu\le\max_{m\le N^2}\tau(m)=N^{o(1)}$. (c) In the
index of the newform $g$ of \cref{thm:cminput}, let $\mathcal T_g$ be the set of $n=pm\le Y$ with $p$
a split prime in $(Y^{1/2},Y^{3/4}]$ with $a(p)\ne0$ and $a(m)\ne0$ (the split primes with $a(p)=0$
are $O(\sqrt Y)$ in number by the count below with $t=0$); then $a(n)=a(p)a(m)\ne0$, $p$ is determined by
$n$, and $|\mathcal T_g|\gg Y/\sqrt{\log Y}$ by the support asymptotic. For $t\ne0$, the integer
points of the box $|x|,|y|\le2\sqrt Y$ on the curve $2\Re(\zeta(x+y\omega)^{k-1})=t$ number $O(\sqrt Y)$,
so $\cnt\{p\le Y:a(p)=t\}\le C_g\sqrt Y$, and
$\cnt\{n\in\mathcal T_g:a(n)=y\}\le C_gY^{1/2}\cdot Y^{3/8}$. With $A=k$ and $\theta=\frac1{16}$,
$\mu Y^\theta=o(|\mathcal T_g|)$. The coefficients of $E_b^r$ are those of $g$ along an arithmetic
progression of indices, so the bound pulls back.
\end{proof}

For comparison, $F=\sum_m(2^{w(m)}-1)2^{-m}$, with $w(m)$ the binary length of $m$, lies in $K_2$
although its raw width is $\sim N\log_2N$: the decomposition $2^{w(m)}-2^0$ has net sequence supported
on $\{0\}\cup\{2^s-s-1\}$ because the $-1$ at index $m$ cancels the $+1$ contributed by another
index. \Cref{cor:wall} says that for the four modular examples every coefficient-by-coefficient
decomposition places $\gg N\sqrt{\log N}$ signed digits, so membership could only come from massive
cancellation between different coefficients; by \cref{thm:completeness} that is exactly the
membership question. Weight one is the threshold at which $|a(n)|\le\tau(n)=n^{o(1)}$ and the typical
footprint of a coefficient is bounded. At $N=2^{14}$ the binary orbits of $E_2^4$ and $E_2^6$ have
$D_N$ near $0.01$, the size expected for independent samples; this is a finite measurement.

\Cref{thm:G} follows from \cref{thm:completeness}, \cref{cor:wall} and
\cref{prop:squaring,prop:thetapowers}.

\section{Open problems}\label{sec:open}

\begin{enumerate}[nosep]
\item Is $E_b^4\in K_b$? By \cref{thm:completeness} this is equivalent to a lower bound for the canonical
net density of its digits, and \cref{cor:wall} rules out coefficient-by-coefficient arguments.
\item Prove $D^+(P_b)<1$, or $P_b\notin K_b$, in some base.
\item Determine the constant of \cref{thm:primedensity} above the threshold $\delta_b^*$. It is fixed by
the frequency of long base-$b$ runs of $0$ or $b-1$ in $M\cdot2^m$.
\item Exact constants for theta values of non-fundamental and imprimitive forms.
\item Is $E_2^2$ simply normal in base $2$? Equivalently, does the gap-weighted proportion of positive
oriented values tend to $\frac12$?
\item Is the converse of \cref{thm:oddorder} true in base $2$: is $\frac adE_2$ simply normal only when
$\ord_d(2)$ is odd? Is some rational scaling of $E_b$ simply normal in a base $b\ge3$?
\item Determine the limit law of $\frac adE_b^2$ for odd $d>1$.
\item Second-order terms of the deficits of $E_b$ and $E_b^3$, which oscillate with $\fr{\log_bN}$.
\item Is $\Theta_Q(1/b)$ disjunctive for every $Q$ and $b$? Is $E_b^2$ disjunctive?
\end{enumerate}

\section{Verification and data}\label{sec:verification}

The Lean declarations used are listed below with the statements that rest on them. All links point
to commit \texttt{347ec3fe631a3e7ab2ae25d53508a3a4c9ea9fa8}.

\begin{center}
\small
\begin{tabular}{@{}p{0.52\linewidth}p{0.42\linewidth}@{}}
\toprule
Declaration & Used in\\
\midrule
\leanP{SharpTranslation}{386}{sharp_translation_isGreatest} & \cref{thm:sharp}, \eqref{eq:sharp}\\
\leanP{SharpTranslation}{123}{arc_le}, \leanP{SharpTranslation}{167}{interval_eq_arc_excess} & \cref{lem:arcs}\\
\leanP{Phase18Synthetic}{524}{discrepancy_deficit_le} & \cref{lem:deficit}\\
\leanP{Phase18Synthetic}{768}{sparse_farCount_le} & \cref{thm:eulerdeficit} (upper bound)\\
\leanP{CoefficientTransfer}{686}{one_sub_discrepancy_ge_of_bulk_and_far} & \cref{lem:circle}\\
\leanP{CharacterLambertTransfer}{175}{radix_pow_linear_transfer}, \leanP{Phase18Synthetic}{681}{discrepancy_add_int} & \cref{lem:linear}\\
\leanP{CoefficientTransfer}{392}{finite_sparse_mass_transfer} and three companions & \cref{thm:leantransfer}\\
\leanP{GeneralMassTransfer}{253}{finite_poly_mass_transfer} and three companions & \cref{thm:leantransfer}\\
\leanP{CoefficientTransfer}{588}{radixVal_mul_radixUnit}, \decl{not_dvd_radixUnit}, \leanP{CoefficientTransfer}{617}{displaced_term_eq} & \cref{lem:radixval}\\
\texttt{RadixWidthLaw}: seven declarations & \cref{thm:displacedwidth,lem:onesource,thm:widthlaw,lem:firstmoment,lem:valuations}\\
\texttt{EulerPentagonal}: ten declarations & \cref{prop:leaneuler}\\
\texttt{DerivativeTransfer}: six declarations & \cref{prop:leanderiv}\\
\texttt{AntisymmetricLambert}: four declarations & \cref{thm:oddweights,thm:qdigamma}\\
\leanP{GeneralMassTransfer}{991}{qDigamma_reflection_eq_residueLambert} & \cref{thm:qdigamma}\\
\texttt{CharacterLambertTransfer}: four mod-$8$ declarations & \cref{prop:residueclass}\\
\leanP{NormalReciprocalCore}{336}{eulerAutomaton_purely_periodic}, \leanP{NormalReciprocalCore}{315}{eulerAutomatonStep_bijective} & \cref{thm:automaton}\\
\bottomrule
\end{tabular}
\end{center}

These declarations are finite or qualitative statements about explicitly defined series and real
numbers. Lean does not check Euler's or Jacobi's identity, any statement about $E_b^2$, any analytic
number theory input, the exact constants, the digit formulas or the sieve of \cref{sec:sieve}. Exact
scripts in the repository reproduce every finite number quoted in the text; they illustrate the
theorems and are used in no proof. The theorem-by-theorem evidence map, the reproduction commands
and the literature ledger accompany the source of this manuscript (directory
\texttt{papers/11-generating-functions-digits}, files \texttt{VERIFICATION.md} and
\texttt{LEDGER.md}). Human specialist review of the written proofs is outstanding.

\section{Use of AI tools}\label{sec:ai}

This manuscript was produced in an AI-assisted research project. Large language models proposed
conjectures and proof strategies, wrote the proofs and the text, implemented the exact computations
and the Lean~4 formalizations, searched the literature and reviewed drafts across several model
providers. The repository owner directed the project and is responsible for its release. The Lean
kernel checked the declarations listed in \cref{sec:verification}. The written proofs, literature
comparisons and priority statements have been checked only by AI models from different providers.
They have not been checked by an independent human expert. Cross-model review is not peer review.

\bibliographystyle{plainnat}
\bibliography{references}

\end{document}
